% III : formuler les idées essentielles d’un texte à résumer — Français 1re Tech — Cours signé OnBuch+
% Programme officiel MINESEC. Compiler avec : tectonic N19-iii-formuler-idees-essentielles-texte-resumer.tex
\newcommand{\DOCMATIERE}{Français}
\newcommand{\DOCNIVEAU}{1re Tech}
\newcommand{\DOCMODULE}{Français – classe de Première technique}
\newcommand{\DOCLECON}{N19}
\newcommand{\DOCTITRE}{III : formuler les idées essentielles d’un texte à résumer}
% OnBuch+ — préambule « cours signés » ENRICHI (compilé avec tectonic / XeLaTeX)
% Système visuel « Pop » d'OnBuch+ : fond crème (jamais blanc), bordures encre
% épaisses, ombres « dures » décalées, titres Archivo Black en capitales.
% Version MATHÉMATIQUES : graphes (pgfplots/tikz), boîtes pédagogiques
% (définition, propriété, méthode, attention, le savais-tu, exercice résolu,
% exercice, TP, bilan), page de garde et signature OnBuch+ ancrée en bas.
%
% Macros attendues dans main.tex AVANT \input{preamble} :
%   \DOCMATIERE  \DOCNIVEAU  \DOCMODULE  \DOCLECON (numéro)  \DOCTITRE
\documentclass[11pt]{article}

\usepackage{fontspec}
\usepackage[french]{babel}
\usepackage[a4paper,margin=2cm,top=3.4cm,bottom=2.8cm,headheight=1.4cm,footskip=1.3cm]{geometry}
\usepackage{amsmath,amssymb,amsfonts}
\usepackage{mathtools} % \xmapsto, \coloneqq... souvent utilisés par les modèles
\usepackage{cancel} % simplifications barrées (bilans de forces)
% \abs{z} (module, valeur absolue) et \norm{u} (norme), utilisés par les modèles
\providecommand{\abs}{}\let\abs\relax\DeclarePairedDelimiter{\abs}{\lvert}{\rvert}
\providecommand{\norm}{}\let\norm\relax\DeclarePairedDelimiter{\norm}{\lVert}{\rVert}
\usepackage[table]{xcolor} % [table] : fournit aussi \rowcolors (alternance de lignes)
\usepackage{pagecolor}
\usepackage{tikz}
\usetikzlibrary{arrows.meta,positioning,calc,shapes.geometric,shapes.misc,shapes.arrows,decorations.pathmorphing,decorations.pathreplacing,decorations.markings,patterns,shadows,mindmap,trees,fit,backgrounds,matrix,intersections,angles,quotes}
\usepackage{pgfplots}
\usepackage[version=4]{mhchem} % équations nucléaires \ce{^{235}_{92}U}
\usepackage[european,siunitx]{circuitikz} % schémas électriques
\pgfplotsset{compat=1.18}
\usepgfplotslibrary{fillbetween,groupplots} % aires entre courbes (calcul intégral)
\usepackage{enumitem}
\usepackage{fancyhdr}
% [most] charge toutes les bibliothèques tcolorbox (listings, minted, raster,
% documentation...) et gonfle inutilement la mémoire TeX sur un document long
% (~300 boîtes) : on ne charge que ce qui est réellement utilisé.
\usepackage[skins,breakable]{tcolorbox}
\usepackage{graphicx}
\usepackage{colortbl}
\usepackage{booktabs}
\usepackage{tabularx}
\usepackage{multirow}
\usepackage{diagbox} % case d'en-tête diagonale des tableaux à double entrée
% Colonnes extensibles centrée / gauche / droite (C, L, R), souvent utilisées par les modèles
\newcolumntype{C}{>{\centering\arraybackslash}X}
\newcolumntype{L}{>{\raggedright\arraybackslash}X}
\newcolumntype{R}{>{\raggedleft\arraybackslash}X}
\usepackage{array}
\usepackage{multicol}
\usepackage{siunitx}
% parse-numbers=false : accepte aussi \SI{2,5 \times 10^{-3}}{...}
\sisetup{locale=FR,output-decimal-marker={,},group-minimum-digits=4,parse-numbers=false}
\DeclareSIUnit\ohm{\ensuremath{\symup{\Omega}}} % U+2126 absent de la police
\DeclareSIUnit\division{div} % réglages d'oscilloscope (ms/div, V/div)
\DeclareSIUnit\tr{tr} % tours (vitesse de rotation, ex. \SI{1500}{\tr\per\minute})
\DeclareSIUnit\molar{M} % concentration molaire (ex. \SI{3}{\molar}, \SI{1}{\micro\molar})
\DeclareSIUnit\an{an} % année (le modèle confond parfois avec \year, un compteur TeX) — ex. \SI{5}{\kilo\meter\per\an}
\DeclareSIUnit\FCFA{FCFA} % franc CFA (ex. \SI{500}{\FCFA\per\kilo\gram})
\DeclareSIUnit\degreeBrix{\textdegree Brix} % teneur en sucre d'un sirop/jus (réfractométrie)
\DeclareSIUnit\month{mois} % le modèle utilise parfois \month, un compteur TeX, comme unité de durée
\DeclareSIUnit\microliter{\micro\liter} % le modèle écrit parfois \microliter en un seul token
\DeclareSIUnit\million{millions} % dénombrement (ex. \SI{15}{\million\per\mL} spermatozoïdes)
\usepackage{qrcode}
% \degree : commande fréquemment utilisée par les modèles pour un angle ou une
% température en dehors de \SI{}{} (non fournie par les paquets chargés).
\providecommand{\degree}{^{\circ}}
% \qed (fin de démonstration), utilisé par les modèles sans amsthm
\providecommand{\qed}{\hfill$\square$}
% \barre : double barre des valeurs interdites dans un tableau de variations (tkz-tab)
\providecommand{\barre}{\ensuremath{\|}}
% environnement proof (amsthm non chargé) utilisé par les modèles
\ifdefined\proof\else\newenvironment{proof}[1][Démonstration]{\par\noindent\textit{#1.}\ }{\qed\par}\fi
\usepackage{lastpage}
\usepackage{hyperref}
\hypersetup{hidelinks}

% ============================================================
% Palette « Pop » OnBuch+ — mêmes hex que la classe P (plus_ui.dart)
% ============================================================
\definecolor{poporange}{HTML}{F59321}
\definecolor{popdark}{HTML}{E07A0C}
\definecolor{popcream}{HTML}{FFF6E8}
\definecolor{popink}{HTML}{1C1714}
\definecolor{poppurple}{HTML}{7A5AE0}
\definecolor{popgreen}{HTML}{2E9E6A}
\definecolor{poppink}{HTML}{E0457B}
\definecolor{popblue}{HTML}{2F7DE1}
\definecolor{popgold}{HTML}{C98A12}
\definecolor{popmuted}{HTML}{8A7F76}
\definecolor{popline}{HTML}{EADFD2}
\definecolor{popgray}{HTML}{8A7F76}
\colorlet{popgrey}{popgray}
% Alias tolérants (le modèle nomme parfois les couleurs différemment)
\colorlet{popwhite}{white}
\colorlet{popblack}{popink}
\colorlet{popred}{poppink}
\colorlet{popyellow}{popgold}
% Teintes claires (fonds de tableaux/encadrés) — le modèle nomme parfois
% les couleurs avec un suffixe L (ex. popblueL) sans les avoir définies.
\colorlet{poporangeL}{poporange!20}
\colorlet{popdarkL}{popdark!20}
\colorlet{poppurpleL}{poppurple!20}
\colorlet{popgreenL}{popgreen!20}
\colorlet{poppinkL}{poppink!20}
\colorlet{popblueL}{popblue!20}
\colorlet{popgoldL}{popgold!20}
\colorlet{popredL}{popred!20}
\colorlet{popgrayL}{popgray!20}
% Teintes claires pour fonds de tableaux / zones
\colorlet{popblueL}{popblue!10!white}
\colorlet{poporangeL}{poporange!14!white}
\colorlet{popgreenL}{popgreen!12!white}
\colorlet{poppurpleL}{poppurple!12!white}
\colorlet{poppinkL}{poppink!12!white}
\colorlet{popgoldL}{popgold!14!white}

\pagecolor{popcream}

% ============================================================
% Typographie
% ============================================================
\defaultfontfeatures{Ligatures=TeX}
\setmainfont{PlusJakartaSans-Medium.ttf}[Path=../../../fonts/,BoldFont=PlusJakartaSans-Bold.ttf]
% Maths en sans-serif (Fira Math) avec chiffres et lettres droites de Plus
% Jakarta, pour que les formules se fondent dans le texte.
\usepackage[math-style=ISO,bold-style=ISO]{unicode-math}
\setmathfont{FiraMath-Regular.otf}[Path=../../../fonts/]
\setmathfont{PlusJakartaSans-Medium.ttf}[Path=../../../fonts/,range={up/num,up/{Latin,latin},\mathup}]
\usepackage{icomma} % virgule décimale sans espace en mode math
% Caractères mathématiques Unicode absents des polices de texte (Plus Jakarta,
% Archivo Black) : rendus par leur équivalent mathématique, en texte comme en
% formule — sinon ils s'affichent en carré vide (ex. « Arithmétique dans ℤ »).
\usepackage{newunicodechar}
\newunicodechar{ℝ}{\ensuremath{\mathbb{R}}}
\newunicodechar{ℤ}{\ensuremath{\mathbb{Z}}}
\newunicodechar{ℂ}{\ensuremath{\mathbb{C}}}
\newunicodechar{ℕ}{\ensuremath{\mathbb{N}}}
\newunicodechar{ℚ}{\ensuremath{\mathbb{Q}}}
\newunicodechar{𝔻}{\ensuremath{\mathbb{D}}}
\newunicodechar{α}{\ensuremath{\alpha}}
\newunicodechar{β}{\ensuremath{\beta}}
\newunicodechar{γ}{\ensuremath{\gamma}}
\newunicodechar{δ}{\ensuremath{\delta}}
\newunicodechar{ε}{\ensuremath{\varepsilon}}
\newunicodechar{θ}{\ensuremath{\theta}}
\newunicodechar{λ}{\ensuremath{\lambda}}
\newunicodechar{μ}{\ensuremath{\mu}}
\newunicodechar{σ}{\ensuremath{\sigma}}
\newunicodechar{τ}{\ensuremath{\tau}}
\newunicodechar{φ}{\ensuremath{\varphi}}
\newunicodechar{ψ}{\ensuremath{\psi}}
\newunicodechar{ω}{\ensuremath{\omega}}
\newunicodechar{Σ}{\ensuremath{\Sigma}}
% majuscules produites par \MakeUppercase dans les titres (α-aminés -> Α)
\newunicodechar{Α}{\ensuremath{\alpha}}
\newunicodechar{Β}{\ensuremath{\beta}}
\newunicodechar{Γ}{\ensuremath{\Gamma}}
\newunicodechar{Δ}{\ensuremath{\Delta}}
\newunicodechar{Ε}{\ensuremath{\varepsilon}}
\newunicodechar{Θ}{\ensuremath{\Theta}}
\newunicodechar{Λ}{\ensuremath{\Lambda}}
\newunicodechar{Μ}{\ensuremath{\mu}}
\newunicodechar{Τ}{\ensuremath{\tau}}
\newunicodechar{Φ}{\ensuremath{\Phi}}
\newunicodechar{Ψ}{\ensuremath{\Psi}}
\newunicodechar{Ω}{\ensuremath{\Omega}}
\newunicodechar{↦}{\ensuremath{\mapsto}}
\newunicodechar{∈}{\ensuremath{\in}}
\newunicodechar{∉}{\ensuremath{\notin}}
\newunicodechar{∧}{\ensuremath{\wedge}}
\newunicodechar{⇔}{\ensuremath{\Leftrightarrow}}
\newunicodechar{⟺}{\ensuremath{\Longleftrightarrow}}
\newunicodechar{⇒}{\ensuremath{\Rightarrow}}
\newunicodechar{⟹}{\ensuremath{\Longrightarrow}}
\newunicodechar{∘}{\ensuremath{\circ}}
\newunicodechar{∪}{\ensuremath{\cup}}
\newunicodechar{∩}{\ensuremath{\cap}}
\newunicodechar{≡}{\ensuremath{\equiv}}
\newunicodechar{⊂}{\ensuremath{\subset}}
\newunicodechar{⊥}{\ensuremath{\perp}}
\newunicodechar{∃}{\ensuremath{\exists}}
\newunicodechar{∀}{\ensuremath{\forall}}
\newunicodechar{∛}{\ensuremath{\sqrt[3]{\,}}}
\newunicodechar{∜}{\ensuremath{\sqrt[4]{\,}}}
\newunicodechar{⃗}{\ensuremath{{}^{\to}}}
\newunicodechar{ⁿ}{\ifmmode^{n}\else\textsuperscript{\ensuremath{n}}\fi}
\newunicodechar{ˣ}{\ifmmode^{x}\else\textsuperscript{\ensuremath{x}}\fi}
\newunicodechar{ᵇ}{\ifmmode^{b}\else\textsuperscript{\ensuremath{b}}\fi}
\newunicodechar{ʳ}{\ifmmode^{r}\else\textsuperscript{\ensuremath{r}}\fi}
\newunicodechar{ᵘ}{\ifmmode^{u}\else\textsuperscript{\ensuremath{u}}\fi}
\newunicodechar{ʸ}{\ifmmode^{y}\else\textsuperscript{\ensuremath{y}}\fi}
\newunicodechar{ᵃ}{\ifmmode^{a}\else\textsuperscript{\ensuremath{a}}\fi}
\newunicodechar{ᵉ}{\ifmmode^{e}\else\textsuperscript{\ensuremath{e}}\fi}
\newunicodechar{ᵏ}{\ifmmode^{k}\else\textsuperscript{\ensuremath{k}}\fi}
\newunicodechar{ᵖ}{\ifmmode^{p}\else\textsuperscript{\ensuremath{p}}\fi}
\newunicodechar{ᴺ}{\ifmmode^{N}\else\textsuperscript{\ensuremath{N}}\fi}
\newunicodechar{ᶜ}{\ifmmode^{c}\else\textsuperscript{\ensuremath{c}}\fi}
\newunicodechar{ʰ}{\ifmmode^{h}\else\textsuperscript{\ensuremath{h}}\fi}
\newunicodechar{ᵠ}{\ifmmode^{q}\else\textsuperscript{\ensuremath{q}}\fi}
\newunicodechar{ₙ}{\ifmmode_{n}\else\textsubscript{\ensuremath{n}}\fi}
\newunicodechar{ₐ}{\ifmmode_{a}\else\textsubscript{\ensuremath{a}}\fi}
\newunicodechar{ᵢ}{\ifmmode_{i}\else\textsubscript{\ensuremath{i}}\fi}
\newunicodechar{ᵣ}{\ifmmode_{r}\else\textsubscript{\ensuremath{r}}\fi}
\newunicodechar{ₖ}{\ifmmode_{k}\else\textsubscript{\ensuremath{k}}\fi}
\newunicodechar{ᵦ}{\ifmmode_{\beta}\else\textsubscript{\ensuremath{\beta}}\fi}
\newunicodechar{ₓ}{\ifmmode_{x}\else\textsubscript{\ensuremath{x}}\fi}
\newfontfamily\popdisplay{ArchivoBlack-Regular.ttf}[Path=../../../fonts/]
\newfontfamily\poplogo{SpaceGrotesk-Bold.ttf}[Path=../../../fonts/]
\newfontfamily\popsemi{PlusJakartaSans-SemiBold.ttf}[Path=../../../fonts/]
\newfontfamily\popxbold{PlusJakartaSans-ExtraBold.ttf}[Path=../../../fonts/]
\newfontfamily\popxboldls{PlusJakartaSans-ExtraBold.ttf}[Path=../../../fonts/,LetterSpace=3.0]

\setlist[enumerate]{leftmargin=1.7em,itemsep=5pt,topsep=4pt}
\setlist[itemize]{leftmargin=1.7em,itemsep=4pt,topsep=4pt,label={\color{poporange}\textbullet}}
\setlength{\parindent}{0pt}
\setlength{\parskip}{5pt}
\renewcommand{\arraystretch}{1.25}

% pgfplots : style Pop par défaut
\pgfplotsset{
  every axis/.append style={axis line style={popink,line width=1pt},tick style={popink},
    grid style={popline,line width=0.5pt},label style={font=\small},tick label style={font=\footnotesize}},
  popcurve/.style={poporange,line width=1.8pt,no markers,smooth},
}
% Même style disponible dans un \draw TikZ simple (hors pgfplots)
\tikzset{popcurve/.style={poporange,line width=1.8pt}}
\tikzset{popgrid/.style={popline,very thin}}
\colorlet{popcurve}{poporange} % le nom du style est parfois utilisé comme couleur

% ============================================================
% Pill / squiggle
% ============================================================
\newcommand{\poppill}[3]{%
  \tikz[baseline=-0.55ex]\node[draw=popink,line width=1.2pt,rounded corners=2.6pt,%
    fill=#2,inner xsep=6.5pt,inner ysep=3pt]{%
    \popxboldls\fontsize{7}{7}\selectfont\color{#3}\MakeUppercase{#1}};%
}
\newcommand{\popsquiggle}[1]{%
  \tikz[baseline=-0.4ex]{
    \draw[#1,line width=2.6pt,line cap=round]
      plot [smooth] coordinates {(0,0.09) (0.34,0.01) (0.68,0.21) (1.02,0.01) (1.36,0.10)};
  }%
}
% Mot-clé surligné (marqueur orange sous le texte)
\newcommand{\cle}[1]{{\popxbold\color{popdark}#1}}

% ============================================================
% En-tête / pied
% ============================================================
\pagestyle{fancy}
\fancyhf{}
\renewcommand{\headrulewidth}{0pt}
\renewcommand{\footrulewidth}{0pt}
\lhead{\raisebox{-0.15ex}{\poplogo\fontsize{13}{13}\selectfont\color{popink}OnBuch\color{poporange}+}}
\rhead{\poppill{\DOCMATIERE\ \textbullet\ \DOCNIVEAU\ \textbullet\ Leçon \DOCLECON}{white}{popink}}
\lfoot{\popxbold\fontsize{7.6}{7.6}\selectfont\color{popmuted}\MakeUppercase{Cours signé OnBuch+}\ \textbullet\ {\popsemi\DOCTITRE}}
\rfoot{\tikz[baseline=-0.6ex]\node[rounded corners=8.5pt,fill=popink,minimum height=17pt,minimum width=17pt,inner xsep=5pt,inner sep=0]{\popxbold\fontsize{8}{8}\selectfont\color{popcream}\thepage\,/\,\pageref*{LastPage}};}

% ============================================================
% Titres
% ============================================================
\newcommand{\chaptitre}[2]{%
  \begin{tcolorbox}[enhanced,colback=poporange,colframe=popink,boxrule=2.6pt,arc=11pt,
      drop shadow=popink,left=15pt,right=15pt,top=12pt,bottom=11pt,before skip=3pt,after skip=18pt]
    \poppill{#2}{popink}{popcream}\par\vspace{9pt}
    {\raggedright\hyphenpenalty=10000\exhyphenpenalty=10000\popdisplay\fontsize{22}{24}\selectfont\color{popcream}\MakeUppercase{#1}\par}
  \end{tcolorbox}
}
% Section numérotée : pastille orange avec le numéro + titre Archivo
\newcounter{popsec}
\newcommand{\coursec}[1]{%
  \stepcounter{popsec}\setcounter{popsub}{0}%
  \par\vspace{16pt}\noindent
  \begin{minipage}{\linewidth}
  \tikz[baseline=-0.7ex]\node[draw=popink,line width=1.6pt,fill=poporange,rounded corners=4pt,minimum size=22pt,inner sep=0]{\popdisplay\fontsize{12}{12}\selectfont\color{popcream}\thepopsec};%
  \hspace{8pt}{\popdisplay\fontsize{15}{17}\selectfont\color{popink}\MakeUppercase{#1}}\par
  \vspace{4pt}\noindent{\color{poporange}\rule{36pt}{3.6pt}}
  \end{minipage}\par\nopagebreak\vspace{8pt}
}
\newcounter{popsub}
\newcommand{\courssub}[1]{%
  \stepcounter{popsub}%
  \par\vspace{9pt}\noindent{\popxbold\fontsize{11.5}{13}\selectfont\color{popink}{\color{poporange}\thepopsec.\thepopsub}\hspace{6pt}#1}\par\nopagebreak\vspace{4pt}
}

% ============================================================
% Boîtes pédagogiques — même recette : fond blanc, bordure épaisse colorée,
% ombre dure encre, badge plein. Titre optionnel [#1].
% ============================================================
\tcbset{popbase/.style={breakable,enhanced,colback=white,boxrule=2.3pt,arc=9pt,
  shadow={3pt}{-3pt}{0pt}{popink},left=14pt,right=14pt,top=11pt,bottom=11pt,before skip=11pt,after skip=15pt,
  fontupper=\popsemi\fontsize{10.6}{14.5}\selectfont\color{popink},
  fontlower=\popsemi\fontsize{10.6}{14.5}\selectfont\color{popink}}}
% Coche dessinée (le glyphe ✓ n'existe pas dans les polices du document)
\newcommand{\popcheck}[1]{\tikz[baseline=-0.5ex]\draw[#1,line width=1.6pt,line cap=round,line join=round](-0.13,0.0)--(-0.04,-0.09)--(0.13,0.1);}
\providecommand{\coche}{\popcheck{popgreen}}
\providecommand{\resultat}[1]{\par\smallskip\noindent\fcolorbox{popgreen}{popgreenL}{\parbox{\dimexpr\linewidth-2\fboxsep-2\fboxrule\relax}{\textbf{Résultat.} #1}}\par\smallskip}
\newcommand{\popboxhead}[3]{\poppill{#1}{#2}{white}\ifx\relax#3\relax\else\hspace{6pt}{\popxbold\fontsize{10.6}{12}\selectfont\color{popink}#3}\fi\par\vspace{7pt}}

\newtcolorbox{aretenir}[1][]{popbase,colframe=popgreen,before upper={\popboxhead{À retenir}{popgreen}{#1}}}
\newtcolorbox{exemplebox}[1][]{popbase,colframe=poporange,before upper={\popboxhead{Exemple}{poporange}{#1}}}
\newtcolorbox{definition}[1][]{popbase,colframe=popblue,before upper={\popboxhead{Définition}{popblue}{#1}}}
\newtcolorbox{propriete}[1][]{popbase,colframe=poppurple,before upper={\popboxhead{Propriété}{poppurple}{#1}}}
\newtcolorbox{methode}[1][]{popbase,colframe=popgold,before upper={\popboxhead{Méthode}{popgold}{#1}}}
\newtcolorbox{attention}[1][]{popbase,colframe=poppink,before upper={\popboxhead{Attention}{poppink}{#1}}}
\newtcolorbox{experience}[1][]{popbase,colframe=popblue,colback=popblueL,before upper={\popboxhead{Expérience}{popblue}{#1}}}
% « Le savais-tu ? » : carte violette pleine (contexte, culture, Cameroun)
\newtcolorbox{savaistu}[1][]{popbase,colback=poppurple,colframe=popink,
  fontupper=\popsemi\fontsize{10.4}{14.2}\selectfont\color{white},
  before upper={\poppill{Le savais-tu\,?}{white}{poppurple}\ifx\relax#1\relax\else\hspace{6pt}{\popxbold\color{white}#1}\fi\par\vspace{7pt}}}
% Exercice résolu : énoncé en haut, \tcblower puis la solution sur fond crème
\newcounter{popexo}\newcounter{popexr}
\newtcolorbox{exoresolu}[1][]{popbase,colframe=poporange,colbacklower=popcream,
  segmentation style={popink,line width=1.2pt,dashed},
  before upper={\stepcounter{popexr}\popboxhead{Exercice résolu \thepopexr}{poporange}{#1}},
  before lower={\poppill{Solution}{popink}{popcream}\par\vspace{6pt}}}
% Exercice (non résolu en place) — la correction est dans la partie Corrigés
\newtcolorbox{exercice}[1][]{popbase,colframe=popink,
  before upper={\stepcounter{popexo}\popboxhead{Exercice \thepopexo}{popink}{#1}}}
% Corrigé
\newtcolorbox{corrige}[1][]{popbase,colframe=popgreen,colback=popgreenL,
  before upper={\popboxhead{Corrigé}{popgreen}{#1}}}
% Objectifs / prérequis / bilan
\newtcolorbox{objectifs}{popbase,colframe=popink,colback=poporangeL,
  before upper={\popboxhead{Objectifs de la leçon}{popink}{}}}
\newtcolorbox{prerequis}{popbase,colframe=popink,colback=popblueL,
  before upper={\popboxhead{Prérequis}{popblue}{}}}
\newtcolorbox{bilan}{popbase,colframe=popink,colback=white,boxrule=2.8pt,
  before upper={\popboxhead{Fiche bilan}{poporange}{L'essentiel à connaître pour l'examen}}}
% Figure encadrée avec légende
\newtcolorbox{popfigure}[1][]{enhanced,breakable=false,colback=white,colframe=popline,boxrule=1.4pt,arc=8pt,
  left=10pt,right=10pt,top=10pt,bottom=8pt,before skip=10pt,after skip=12pt,halign=center,
  lower separated=false,halign lower=center,
  fontlower=\popsemi\fontsize{9}{11.5}\selectfont\color{popmuted},
  #1}
\newcommand{\legende}[1]{\par\vspace{6pt}{\popsemi\fontsize{9}{11.5}\selectfont\color{popmuted}#1}}

% ============================================================
% Signature OnBuch+
% ============================================================
\newcommand{\poplogoinline}[1]{{\poplogo\fontsize{#1}{#1}\selectfont\color{popink}OnBuch\color{poporange}+}}
\newcommand{\signatureonbuch}{%
  \begin{tcolorbox}[enhanced,colback=popink,colframe=popink,boxrule=2.2pt,arc=10pt,
      drop shadow=poporange,left=14pt,right=14pt,top=10pt,bottom=10pt,before skip=0pt,after skip=0pt]
    \tikz[baseline=-0.7ex]\node[circle,fill=popgreen,draw=popcream,line width=1.3pt,minimum size=22pt,inner sep=0]{\popcheck{white}};%
    \hspace{9pt}\begin{minipage}[c]{0.62\linewidth}
      {\popxbold\fontsize{12}{13}\selectfont\color{popcream}Signé\ \textbullet\ {\poplogo OnBuch\color{poporange}+}}\par\vspace{2pt}
      {\raggedright\popsemi\fontsize{8}{10.5}\selectfont\color{popline}\MakeUppercase{Équipe pédagogique OnBuch+}\\\MakeUppercase{Conforme au programme officiel MINESEC}\par}
    \end{minipage}\hfill
    {\poplogo\fontsize{22}{22}\selectfont\color{popcream}OnBuch\color{poporange}+}
  \end{tcolorbox}%
}

% ============================================================
% Page de garde
% #1 titre · #2 matière · #3 niveau · #4 module · #5 n° leçon · #6 durée ·
% #7 description / objectifs (texte ou itemize)
% ============================================================
\newcommand{\pagegarde}[7]{%
  \thispagestyle{empty}
  \newgeometry{a4paper,margin=1.7cm,top=1.6cm,bottom=1.4cm}%
  \begin{tikzpicture}[remember picture,overlay]
    \fill[poporange!18] ([xshift=-3.2cm,yshift=-3.6cm]current page.north east) circle (4.6cm);
    \fill[poppurple!14] ([xshift=2.4cm,yshift=7.6cm]current page.south west) circle (3.8cm);
  \end{tikzpicture}%
  {\poplogo\fontsize{30}{30}\selectfont\color{popink}OnBuch\color{poporange}+}\hfill
  \raisebox{0.6ex}{\poppill{Cours signé \textbullet\ Premium}{popink}{popcream}}\par
  \vspace{1pt}\popsquiggle{poppurple}\par
  \vspace{8pt}
  \begin{tcolorbox}[enhanced,colback=poporange,colframe=popink,boxrule=2.8pt,arc=13pt,
      drop shadow=popink,left=18pt,right=18pt,top=12pt,bottom=13pt,after skip=10pt]
    \poppill{#2 \textbullet\ #3}{popink}{popcream}\hspace{5pt}\poppill{#4}{white}{popink}\par\vspace{8pt}
    {\popdisplay\fontsize{14}{14}\selectfont\color{popink}LEÇON #5}\par\vspace{4pt}
    {\raggedright\hyphenpenalty=10000\exhyphenpenalty=10000\popdisplay\fontsize{25}{27}\selectfont\color{popcream}\MakeUppercase{#1}\par}
    \vspace{8pt}
    {\popxbold\fontsize{9.5}{9.5}\selectfont\color{popink}\MakeUppercase{Durée conseillée : #6}}
  \end{tcolorbox}
  \begin{tcolorbox}[enhanced,colback=white,colframe=popink,boxrule=2.2pt,arc=10pt,
      drop shadow=popink,left=16pt,right=16pt,top=10pt,bottom=10pt,after skip=8pt]
    \poppill{Dans cette leçon}{poppurple}{white}\par\vspace{5pt}
    \popsemi\fontsize{9.3}{12.6}\selectfont\color{popink}#7
  \end{tcolorbox}
  \noindent\begin{minipage}[t]{0.487\linewidth}
    \begin{tcolorbox}[enhanced,colback=popgreen,colframe=popink,boxrule=2.2pt,arc=10pt,
        drop shadow=popink,left=11pt,right=11pt,top=10pt,bottom=10pt,height=3.1cm,valign=center]
      \tcbox[colback=white,colframe=popink,boxrule=1.3pt,arc=5pt,boxsep=0pt,nobeforeafter,
        left=3pt,right=3pt,top=3pt,bottom=3pt,valign=center]{\qrcode[height=1.9cm]{https://play.google.com/store/apps/details?id=com.luvvix.onbuch&hl=fr}}%
      \hspace{8pt}\begin{minipage}[c]{0.5\linewidth}
        {\popdisplay\fontsize{11}{12.5}\selectfont\color{white}\MakeUppercase{Télécharge OnBuch}}\par\vspace{4pt}
        {\raggedright\popsemi\fontsize{8.4}{11}\selectfont\color{white}Gratuit \textbullet\ Google Play\\Tuteur IA, annales, concours, résultats\par}
      \end{minipage}
    \end{tcolorbox}
  \end{minipage}\hfill
  \begin{minipage}[t]{0.487\linewidth}
    \begin{tcolorbox}[enhanced,colback=poppurple,colframe=popink,boxrule=2.2pt,arc=10pt,
        drop shadow=popink,left=11pt,right=11pt,top=10pt,bottom=10pt,height=3.1cm,valign=center]
      \tikz[baseline=-0.7ex]\node[circle,fill=white,draw=popink,line width=1.3pt,minimum size=1.6cm,inner sep=0]{\popdisplay\fontsize{24}{24}\selectfont\color{poppurple}+};%
      \hspace{8pt}\begin{minipage}[c]{0.6\linewidth}
        {\popdisplay\fontsize{11}{12.5}\selectfont\color{white}\MakeUppercase{Passe à OnBuch+}}\par\vspace{4pt}
        {\raggedright\popsemi\fontsize{8.4}{11}\selectfont\color{white}Tous les cours signés, Tuteur IA illimité, fiches PDF — onglet OnBuch+ de l'app\par}
      \end{minipage}
    \end{tcolorbox}
  \end{minipage}
  \vspace{8pt}
  \popcontenu
  \vfill
  \signatureonbuch
  \restoregeometry
  \newpage
}

% Rangée « Ce cours contient » (page de garde)
\newcommand{\poptile}[3]{% #1 couleur #2 gros texte #3 légende
  \begin{tcolorbox}[enhanced,colback=#1,colframe=popink,boxrule=2pt,arc=9pt,shadow={2.5pt}{-2.5pt}{0pt}{popink},
      left=6pt,right=6pt,top=9pt,bottom=9pt,halign=center,valign=center,height=1.75cm,width=\linewidth,nobeforeafter]
    {\popdisplay\fontsize{12.5}{13}\selectfont\color{white}\MakeUppercase{#2}}\par\vspace{3pt}
    {\centering\popsemi\fontsize{7.8}{9.5}\selectfont\color{white}#3\par}
  \end{tcolorbox}}
\newcommand{\popcontenu}{%
  \par\noindent{\popxboldls\fontsize{7.5}{7.5}\selectfont\color{popmuted}CE COURS SIGNÉ CONTIENT}\par\vspace{7pt}
  \noindent\begin{minipage}[t]{0.235\linewidth}\poptile{popblue}{Cours}{complet, illustré, pas à pas}\end{minipage}\hfill
  \begin{minipage}[t]{0.235\linewidth}\poptile{poporange}{Méthodes}{et exercices résolus}\end{minipage}\hfill
  \begin{minipage}[t]{0.235\linewidth}\poptile{poppink}{Exercices}{type examen + corrigés}\end{minipage}\hfill
  \begin{minipage}[t]{0.235\linewidth}\poptile{popgreen}{Fiche bilan}{l'essentiel à retenir}\end{minipage}\par}

% Sommaire
\newtcolorbox{sommaire}{enhanced,colback=white,colframe=popink,boxrule=2.2pt,arc=10pt,
  drop shadow=popink,left=18pt,right=18pt,top=13pt,bottom=13pt,before skip=6pt,after skip=20pt,
  before upper={\poppill{Sommaire}{poppurple}{white}\par\vspace{9pt}\popsemi\fontsize{10.6}{14.5}\selectfont\color{popink}}}

% Signature de fin de cours — poussée tout en bas de la dernière page
\newcommand{\findecours}{%
  \par\vfill\nopagebreak
  \begin{center}\popsquiggle{poporange}\end{center}\vspace{4pt}
  \signatureonbuch
}
\AtBeginDocument{\def\llbracket{\mathopen{[\mkern-3.5mu[}}\def\rrbracket{\mathclose{]\mkern-3.5mu]}}} % intervalles d'entiers (glyphe ⟦ absent de la police)
% Glyphes absents de Fira Math : redéfinis à partir de glyphes disponibles
\AtBeginDocument{%
  \def\setminus{\mathbin{\backslash}}\let\smallsetminus\setminus
  \def\vdots{\mathord{\vbox{\baselineskip3.2pt\lineskiplimit0pt\kern4pt\hbox{.}\hbox{.}\hbox{.}}}}%
  \def\ddots{\mathinner{\mkern1mu\raise6.5pt\hbox{.}\mkern2mu\raise3.6pt\hbox{.}\mkern2mu\raise0.7pt\hbox{.}\mkern1mu}}%
  \def\triangle{\mathord{\scalebox{0.9}{$\symup{\Delta}$}}}%
  \def\nabla{\mathord{\raisebox{\depth}{\scalebox{1}[-1]{$\symup{\Delta}$}}}}%
  \def\checkmark{\popcheck{popdark}}%
  \def\star{\mathbin{\ast}}\def\bigstar{\mathord{\ast}}%
  \def\diamond{\mathbin{\scalebox{0.6}{$\Diamond$}}}%
  \def\bigcup{\mathop{\mathchoice{\raisebox{-0.3em}{\scalebox{1.7}{$\cup$}}}{\scalebox{1.25}{$\cup$}}{\cup}{\cup}}\displaylimits}%
  \def\bigcap{\mathop{\mathchoice{\raisebox{-0.3em}{\scalebox{1.7}{$\cap$}}}{\scalebox{1.25}{$\cap$}}{\cap}{\cap}}\displaylimits}%
  \def\lesseqqgtr{\mathrel{\lessgtr}}\def\bigoplus{\mathop{\oplus}\displaylimits}%
}
\providecommand{\ssi}{si et seulement si} % abréviation fréquente
\AtBeginDocument{\providecommand{\wideparen}{\overparen}} % arc de cercle

%% FIGFIX-BEGIN v3 — correctifs globaux des figures (docs/onbuch-generation/tools/figfix)
\usepackage{adjustbox}\usepackage{etoolbox}
% toute figure TikZ/pgfplots plus large que la ligne est réduite à la largeur disponible
\BeforeBeginEnvironment{tikzpicture}{\begin{adjustbox}{max width=\linewidth}}
\AfterEndEnvironment{tikzpicture}{\end{adjustbox}}
% légendes pgfplots lisibles par-dessus les courbes
\pgfplotsset{every axis/.append style={legend style={fill=white,fill opacity=0.92,text opacity=1,draw=black!15,inner sep=3pt}}}
% étiquettes d'arêtes (syntaxe edge["..."]) avec fond blanc
\tikzset{every edge quotes/.append style={fill=white,inner sep=1.2pt}}
% bulles de cartes mentales : texte à la ligne dans la bulle, bulle un peu plus grande
\tikzset{every concept/.append style={text width=2.0cm,align=center,minimum size=2.3cm,inner sep=2pt}}
%% FIGFIX-END
\begin{document}

\pagegarde{III : formuler les idées essentielles d’un texte à résumer}{Français}{1re Tech}{Français – classe de Première technique}{N19}{9 séances (fiche de progression harmonisée)}{Cette leçon apprend à l'élève de Première technique à formuler les idées essentielles d'un texte à résumer : identifier la structure argumentative, repérer les tonalités, utiliser la synonymie et l'antonymie pour reformuler fidèlement. Elle s'appuie sur la lecture méthodique du Lion et la Perle de Wole Soyinka et prépare directement à l'épreuve de contraction de texte du Baccalauréat technique.
\vspace{4pt}
\begin{multicols}{2}\small
\begin{itemize}
\item À la fin de cette leçon, je sais identifier les caractéristiques d'un texte argumentatif (thèse, arguments, exemples, connecteurs).
\item À la fin de cette leçon, je sais dégager la structure argumentative d'un texte (introduction, développement, conclusion).
\item À la fin de cette leçon, je sais distinguer une idée essentielle d'une idée secondaire, d'un exemple ou d'une répétition.
\item À la fin de cette leçon, je sais reconnaître les tonalités comique, dramatique, satirique et tragique dans un texte.
\item À la fin de cette leçon, je sais utiliser la synonymie et l'antonymie pour reformuler une idée sans la déformer.
\item À la fin de cette leçon, je sais reformuler les idées essentielles d'un texte en respectant un nombre de mots imposé.
\end{itemize}\end{multicols}}

\begin{sommaire}
\begin{multicols}{2}
\begin{enumerate}
\item Le texte argumentatif : caractéristiques
\item La structure argumentative du texte
\item Les tonalités : comique, dramatique, satirique, tragique
\item Lecture méthodique du Lion et la Perle (n°4 et n°5)
\item Synonymie et antonymie : outils de la reformulation
\item La reformulation des idées essentielles : méthode complète du résumé
\item Compte rendu, remédiation et entraînement à l'examen
\item Activité d'intégration
\item Méthodes et exercices résolus
\item Exercices
\item Corrigés des exercices
\item Fiche bilan
\end{enumerate}
\end{multicols}
\end{sommaire}

\begin{objectifs}
\begin{itemize}
\item À la fin de cette leçon, je sais identifier les caractéristiques d'un texte argumentatif (thèse, arguments, exemples, connecteurs).
\item À la fin de cette leçon, je sais dégager la structure argumentative d'un texte (introduction, développement, conclusion).
\item À la fin de cette leçon, je sais distinguer une idée essentielle d'une idée secondaire, d'un exemple ou d'une répétition.
\item À la fin de cette leçon, je sais reconnaître les tonalités comique, dramatique, satirique et tragique dans un texte.
\item À la fin de cette leçon, je sais utiliser la synonymie et l'antonymie pour reformuler une idée sans la déformer.
\item À la fin de cette leçon, je sais reformuler les idées essentielles d'un texte en respectant un nombre de mots imposé.
\item À la fin de cette leçon, je sais rédiger un résumé cohérent, objectif et rédigé à la troisième personne.
\item À la fin de cette leçon, je sais justifier ma démarche de contraction (choix des idées conservées, suppressions opérées).
\item À la fin de cette leçon, je sais appliquer la méthode du résumé à des situations concrètes de la vie courante (compte rendu de réunion, synthèse d'article de presse camerounaise).
\end{itemize}
\end{objectifs}

\begin{prerequis}
\begin{itemize}
\item La notion de paragraphe et de plan de texte (classe de Seconde)
\item Les connecteurs logiques (d'abord, ensuite, cependant, donc...) vus en début d'année
\item Les types de textes (narratif, descriptif, argumentatif) étudiés en classe de Seconde
\item La lecture méthodique n°1 à n°3 du Lion et la Perle (séances précédentes)
\item Le rappel grammatical : la phrase simple et la phrase complexe
\end{itemize}
\end{prerequis}

\newpage

\coursec{Le texte argumentatif : caractéristiques}

\textbf{Situation d'accroche.} À la radio RFI ou sur CRTV, le journaliste résume en 30 secondes un long discours officiel prononcé à Yaoundé ; sur WhatsApp, un élève de Nkongsamba résume en trois lignes un article de deux pages pour ses camarades de classe. Comment font-ils pour garder l'essentiel sans trahir l'auteur ? Lors des examens techniques (Probatoire et Baccalauréat), l'épreuve de contraction de texte exige exactement cette compétence : repérer la \cle{structure argumentative} d'un texte, en dégager les \cle{idées essentielles} et les reformuler en un nombre de mots imposé. Cette leçon vous donne la méthode complète, étape par étape.

\textbf{Problématique.} Qu'est-ce qu'un texte argumentatif, quels sont ses éléments constitutifs, et comment reconnaître ses marques pour pouvoir ensuite le résumer fidèlement ?

\courssub{Qu'est-ce qu'un texte argumentatif ?}

Commençons par l'intuition. Quand votre père vous dit : « Va à l'école, \emph{car} l'école ouvre toutes les portes ; regarde ton oncle devenu ingénieur », il ne raconte pas une histoire : il défend une idée (« il faut aller à l'école »), il donne une raison (« l'école ouvre les portes ») et il cite un cas concret (l'oncle ingénieur). Il \emph{argumente}. Un texte argumentatif fonctionne exactement de la même manière, mais à l'écrit et de façon organisée.

\begin{definition}[Le texte argumentatif]
Le \cle{texte argumentatif} est un texte dans lequel un \cle{énonciateur} (celui qui parle ou écrit) défend une \cle{thèse} à l'aide d'\cle{arguments} et d'\cle{exemples} pour \cle{convaincre} ou \cle{persuader} un \cle{destinataire} (celui qui reçoit le message).
\end{definition}

On rencontre le texte argumentatif partout : dans les éditoriaux des journaux (\emph{Cameroon Tribune}, \emph{Le Messager}, \emph{Mutations}), dans les discours politiques, dans les plaidoyers d'avocats, dans les publicités, dans les essais des écrivains africains comme Mongo Beti ou Ferdinand Oyono, et même dans les débats de quartier. C'est le type de texte le plus fréquent à l'épreuve de contraction de texte du Probatoire et du Baccalauréat technique.

\begin{savaistu}[Le texte argumentatif à l'examen]
Au Baccalauréat technique, le texte à résumer est presque toujours un texte argumentatif tiré de la presse ou d'une œuvre africaine. Savoir repérer sa structure (thèse, arguments, exemples) est donc la première compétence à maîtriser : sans elle, impossible de dégager les idées essentielles.
\end{savaistu}

\courssub{Les éléments constitutifs : thèse, arguments, exemples}

Tout texte argumentatif repose sur trois éléments qui s'emboîtent comme les étages d'une maison.

\textbf{1. La thèse.} C'est \emph{l'idée défendue} par l'auteur, le point de vue qu'il veut faire partager. Exemple : « La scolarisation des jeunes filles est indispensable au développement du Cameroun. » La thèse peut être \cle{explicite} (clairement formulée) ou \cle{implicite} (sous-entendue, le lecteur doit la déduire).

\textbf{2. Les arguments.} Ce sont les \emph{raisons} qui justifient la thèse. Un argument répond à la question : « Pourquoi l'auteur a-t-il raison ? » Exemple : « Une fille instruite fait reculer la pauvreté de toute sa famille. »

\textbf{3. Les exemples.} Ce sont des \emph{cas particuliers, des faits concrets, des chiffres, des anecdotes} qui illustrent et rendent crédible un argument. Exemple : « Dans la région de l'Extrême-Nord, les villages qui ont ouvert des collèges pour filles ont vu le taux de mariages précoces diminuer. »

\begin{attention}[Argument ou exemple ?]
Erreur fréquente : confondre un \cle{argument} (raison \emph{générale}) avec un \cle{exemple} (cas \emph{particulier} qui illustre l'argument). Test simple : si l'énoncé peut être remplacé par « par exemple... », c'est un exemple ; s'il répond à « pourquoi ? », c'est un argument. Dans un résumé, on garde la thèse et les arguments ; on supprime presque toujours les exemples.
\end{attention}

\begin{popfigure}
\begin{center}
\begin{tikzpicture}[
  these/.style={rectangle, draw=popblue, fill=popblueL, rounded corners, thick, text width=6.2cm, align=center, font=\small\bfseries},
  arg/.style={rectangle, draw=poporange, fill=poporangeL, rounded corners, thick, text width=3.6cm, align=center, font=\small},
  ex/.style={rectangle, draw=popgreen, fill=popgreenL, rounded corners, text width=3.2cm, align=center, font=\footnotesize},
  lien/.style={thick, popdark}
]
\node[these, align=center] (t) at (0,0){THÈSE\\ L'idée défendue par l'auteur};
\node[arg, align=center] (a1) at (-5,-2.2){Argument 1\\ (raison générale)};
\node[arg, align=center] (a2) at (0,-2.2){Argument 2\\ (raison générale)};
\node[arg, align=center] (a3) at (5,-2.2){Argument 3\\ (raison générale)};
\node[ex] (e1) at (-5,-4.2) {Exemple : fait, chiffre, anecdote};
\node[ex] (e2) at (0,-4.2) {Exemple : cas particulier};
\node[ex] (e3) at (5,-4.2) {Exemple : témoignage};
\draw[lien] (t) -- (a1);
\draw[lien] (t) -- (a2);
\draw[lien] (t) -- (a3);
\draw[lien] (a1) -- (e1);
\draw[lien] (a2) -- (e2);
\draw[lien] (a3) -- (e3);
\end{tikzpicture}
\end{center}
\legende{Schéma en arbre de la structure argumentative : la thèse au sommet, les arguments en branches, les exemples en feuilles. Dans un résumé, on conserve le tronc et les branches ; on coupe les feuilles.}
\end{popfigure}

\courssub{Convaincre et persuader : deux façons d'agir sur le lecteur}

L'auteur d'un texte argumentatif poursuit l'un de ces deux buts, souvent les deux à la fois :

\begin{itemize}
\item \cle{Convaincre} : s'adresser à la \emph{raison} du lecteur, par des arguments logiques, des preuves, des chiffres, des démonstrations. C'est la démarche du scientifique ou de l'éditorialiste sérieux.
\item \cle{Persuader} : s'adresser aux \emph{émotions} du lecteur (peur, pitié, fierté, colère), par des images frappantes, des mots chargés affectivement, des témoignages touchants. C'est la démarche du publicitaire ou de l'orateur.
\end{itemize}

\begin{exemplebox}[Convaincre ou persuader ?]
\begin{itemize}
\item « Le paludisme tue chaque année des milliers d'enfants au Cameroun ; or la moustiquaire imprégnée coûte moins de 2000 FCFA. » $\rightarrow$ l'auteur \emph{convainc} (chiffres, logique).
\item « Imaginez cette mère de Maroua veillant, seule, le corps brûlant de fièvre de son petit garçon... » $\rightarrow$ l'auteur \emph{persuade} (émotion, image pathétique).
\end{itemize}
\end{exemplebox}

\courssub{Les connecteurs logiques de l'argumentation}

Les \cle{connecteurs logiques} sont les mots-charnières qui articulent les idées entre elles. Ils sont les panneaux indicateurs du texte argumentatif : en les repérant, on suit le raisonnement de l'auteur. On les classe selon leur fonction.

\begin{center}
\begin{tabularx}{\linewidth}{|l|X|X|}
\hline
\rowcolor{popblueL} \textbf{Fonction} & \textbf{Connecteurs} & \textbf{Emploi} \\
\hline
\cle{Addition} & de plus, en outre, également, aussi, d'abord... ensuite & ajouter un argument à un autre \\
\hline
\cle{Opposition} & mais, cependant, pourtant, néanmoins, en revanche, alors que & introduire une objection ou une nuance \\
\hline
\cle{Cause} & car, puisque, parce que, en effet, vu que & justifier, donner la raison \\
\hline
\cle{Conséquence} & donc, par conséquent, ainsi, c'est pourquoi, dès lors & tirer une conclusion, annoncer un résultat \\
\hline
\end{tabularx}
\end{center}

\begin{exemplebox}[Connecteurs en action]
Dans la phrase : « L'école est la clé du développement \textbf{car} elle forme les citoyens ; \textbf{ainsi}, le Cameroun a construit des milliers de salles de classe », le mot \emph{car} introduit l'\textbf{argument} (la raison) et le mot \emph{ainsi} introduit l'\textbf{exemple} (la conséquence concrète). Repérer ces deux connecteurs suffit à comprendre l'architecture de la phrase.
\end{exemplebox}

\courssub{Les procédés énonciatifs : qui parle, et avec quels mots ?}

Au-delà des connecteurs, l'auteur laisse dans son texte des traces de sa présence et de son point de vue. Ce sont les \cle{procédés énonciatifs}.

\textbf{1. Les pronoms personnels.} Le \emph{je} engage l'auteur personnellement (« je dénonce », « j'affirme »). Le \emph{nous} associe le lecteur à la cause et crée une complicité (« nous devons agir »). Le \emph{on} généralise. Le \emph{ils} désigne souvent un adversaire ou un groupe visé.

\textbf{2. Les modalisateurs.} Ce sont les mots qui marquent le degré de certitude ou d'obligation : \emph{il faut}, \emph{il est évident que}, \emph{certes}, \emph{sans doute}, \emph{peut-être}, \emph{assurément}. Ils révèlent l'attitude de l'énonciateur face à sa propre thèse : fermeté (« il faut absolument ») ou prudence (« il semble que »).

\textbf{3. Les champs lexicaux valorisants ou dévalorisants.} Le choix des mots n'est jamais neutre. Parler de « jeunesse dynamique et promise à un bel avenir » (lexique valorisant : \emph{dynamique, promesse, avenir}) n'a pas le même effet que parler de « jeunesse désœuvrée et sacrifiée » (lexique dévalorisant : \emph{désœuvrée, sacrifiée}). Le lexique trahit le jugement de l'auteur.

\courssub{Repérer la thèse : la compétence-clé}

Tout le travail du résumé dépend de cette première opération : trouver la thèse. Trois indices aident à la localiser :

\begin{enumerate}
\item \textbf{La place} : la thèse se trouve \emph{souvent en introduction} (l'auteur annonce d'emblée son point de vue) ou \emph{en conclusion} (l'auteur aboutit à son idée après l'avoir démontrée).
\item \textbf{Les marques linguistiques} : elle est fréquemment introduite par des formules comme « il faut », « nous devons », « il est évident que », « il est temps de », « nul ne peut nier que ».
\item \textbf{Le cas implicite} : parfois la thèse n'est écrite nulle part ; il faut alors la dégager en réunissant les arguments (« que veut me faire admettre l'auteur, au fond ? »).
\end{enumerate}

\begin{methode}[Repérer la thèse en trois questions]
\begin{enumerate}
\item Se demander : « \emph{Que veut me faire admettre l'auteur ?} »
\item Chercher les marques : « il faut », « nous devons », « il est évident que », souvent placées en introduction ou en conclusion.
\item Vérifier : tous les arguments du texte doivent pouvoir se rattacher à la thèse trouvée ; si l'un d'eux reste orphelin, la thèse est mal formulée.
\end{enumerate}
\end{methode}

\courssub{Document d'application : un éditorial sur la scolarisation des jeunes filles}

Lisons cet extrait d'éditorial inspiré de la presse camerounaise, sur le modèle des articles de \emph{Cameroon Tribune} :

\begin{exemplebox}[Extrait d'éditorial]
« Il faut le dire avec force : tant que les jeunes filles de l'Extrême-Nord resteront loin des bancs de l'école, notre région restera loin du développement. En effet, une fille instruite fait reculer la pauvreté de toute sa famille : à Mokolo, les villages qui ont ouvert des collèges pour filles ont vu les mariages précoces diminuer. De plus, l'école protège la santé des mères et des enfants. Certes, les traditions pèsent lourd ; cependant, aucune coutume ne justifie de priver la moitié de nos enfants de leur avenir. Nous devons donc, parents, chefs traditionnels et élus, faire de la scolarisation des filles notre combat commun. »
\end{exemplebox}

Analysons ce texte selon la grille apprise :

\begin{center}
\begin{tabularx}{\linewidth}{|l|X|}
\hline
\rowcolor{popblueL} \textbf{Élément} & \textbf{Repérage dans le texte} \\
\hline
Thèse & « il faut... scolariser les jeunes filles » (annoncée dès l'introduction, reprise en conclusion : « notre combat commun ») \\
\hline
Argument 1 & « une fille instruite fait reculer la pauvreté » (introduit par \emph{en effet}) \\
\hline
Exemple & « à Mokolo, les villages... mariages précoces en baisse » (cas particulier) \\
\hline
Argument 2 & « l'école protège la santé des mères et des enfants » (introduit par \emph{de plus}) \\
\hline
Objection + réfutation & « Certes, les traditions pèsent lourd ; \emph{cependant}... » (opposition) \\
\hline
Énonciation & \emph{nous} (complicité), \emph{il faut} (modalisateur d'obligation), lexique valorisant : \emph{avenir, développement, combat} \\
\hline
\end{tabularx}
\end{center}

\begin{exoresolu}[Identifier thèse, argument et exemple]
\textbf{Énoncé.} Dans la phrase suivante, indiquez la thèse, l'argument et l'exemple : « Nous devons protéger la forêt camerounaise, car elle régule le climat de toute la sous-région ; ainsi, la disparition des forêts du Sud a déjà modifié le régime des pluies. »
\tcblower
\textbf{Solution détaillée.}
\begin{enumerate}
\item \textbf{Thèse} : « Nous devons protéger la forêt camerounaise. » On la reconnaît au modalisateur \emph{nous devons}, marque classique de la thèse.
\item \textbf{Argument} : « elle régule le climat de toute la sous-région. » Le connecteur de cause \emph{car} introduit la raison générale qui justifie la thèse.
\item \textbf{Exemple} : « la disparition des forêts du Sud a déjà modifié le régime des pluies. » Le connecteur \emph{ainsi} introduit un cas concret et localisé (le Sud du Cameroun) qui illustre l'argument.
\end{enumerate}
Vérification : l'argument répond à « pourquoi protéger la forêt ? » et l'exemple répond à « où le voit-on concrètement ? ». La distinction est donc correcte.
\end{exoresolu}

\begin{aretenir}[L'essentiel de la section]
\begin{itemize}
\item Un texte argumentatif défend une \cle{thèse} grâce à des \cle{arguments} (raisons générales) illustrés par des \cle{exemples} (cas particuliers).
\item Il cherche à \cle{convaincre} (par la raison) ou à \cle{persuader} (par l'émotion).
\item Les \cle{connecteurs logiques} (addition, opposition, cause, conséquence) balisent le raisonnement.
\item Les \cle{procédés énonciatifs} (pronoms, modalisateurs, champs lexicaux) révèlent la présence et le point de vue de l'auteur.
\item Pour repérer la thèse : « que veut me faire admettre l'auteur ? » ; elle se situe souvent en introduction ou en conclusion, annoncée par « il faut », « nous devons », « il est évident que ».
\end{itemize}
\end{aretenir}

\begin{exercice}[À vous de jouer]
Relevez dans un article de journal de votre choix (ou dans l'éditorial sur Mokolo ci-dessus) : 1) la thèse ; 2) deux arguments ; 3) un exemple ; 4) trois connecteurs logiques en précisant leur fonction ; 5) deux modalisateurs. Rédigez ensuite une phrase qui résume la thèse en vos propres mots.
\end{exercice}

\coursec{La structure argumentative du texte}

Dans la section précédente, nous avons appris à reconnaître un texte argumentatif : un texte qui défend une thèse, qui cherche à convaincre ou à persuader, et qui s'appuie sur des arguments et des exemples. Mais reconnaître un texte argumentatif ne suffit pas pour en faire la \cle{contraction de texte} (épreuve du Probatoire). Pour formuler les idées essentielles d'un texte, il faut comprendre \emph{comment} ce texte est construit : où commence la thèse, où se trouvent les arguments, où se situe la conclusion. C'est précisément l'objet de cette section : découvrir la \cle{structure argumentative}, c'est-à-dire le plan interne du texte, et apprendre à la dégager méthodiquement. Cette compétence est la clé de la contraction de texte au Probatoire et au Baccalauréat technique.

\courssub{Les trois parties de la structure argumentative}

\begin{definition}[La structure argumentative]
La \cle{structure argumentative} est l'organisation logique des idées d'un texte argumentatif. Elle comprend une \cle{thèse} (l'idée défendue par l'auteur), des \cle{arguments} (les raisons qui soutiennent la thèse), des \cle{exemples} (les illustrations concrètes des arguments) et une \cle{conclusion} (le bilan final). Cette organisation se déploie en trois grandes parties : l'introduction, le développement et la conclusion.
\end{definition}

Imaginons un avocat devant un tribunal de Yaoundé. Il ne commence jamais par \emph{assener} ses preuves sans prévenir : il annonce d'abord ce qu'il va démontrer, puis il présente ses arguments un à un, et il termine par un rappel frappant de sa position. Le texte argumentatif fonctionne exactement de la même manière.

\textbf{1. L'introduction : le sujet et la thèse.} L'introduction présente le \cle{sujet} (le problème dont on va parler) et énonce la \cle{thèse} (la position de l'auteur). Elle se situe généralement dans le ou les premiers paragraphes. Des formules comme « il s'agit de », « le problème est le suivant », « nous soutenons que », « je pense que » signalent souvent la thèse.

\textbf{2. Le développement : les arguments.} Le développement est la partie la plus longue du texte. Il est constitué d'\cle{arguments}, c'est-à-dire de raisons logiques qui prouvent la thèse. Chaque argument est généralement illustré par un ou plusieurs \cle{exemples} (faits, statistiques, anecdotes, citations, références historiques). En général, un argument occupe un paragraphe ou un groupe de paragraphes.

\textbf{3. La conclusion : le bilan et l'ouverture.} La conclusion tire le bilan du raisonnement (« ainsi », « en définitive », « il apparaît donc que », « c'est pourquoi ») et peut proposer une \cle{ouverture} : une solution, un appel à l'action, une perspective nouvelle, une citation de synthèse.

\begin{popfigure}
\begin{center}
\begin{tikzpicture}[
 bloc/.style={draw=popblue, fill=popblueL, rounded corners=2pt, text width=3.1cm, minimum height=1.05cm, align=center, font=\small},
 ex/.style={draw=poporange, fill=poporangeL, rounded corners=2pt, text width=2.5cm, minimum height=0.7cm, align=center, font=\footnotesize},
 fleche/.style={-{Stealth[length=2.5mm]}, thick, popdark}]
\node[bloc, fill=poppurpleL, draw=poppurple, align=center] (intro) at (0,0){\textbf{Introduction}\\ sujet + thèse};
\node[bloc] (a1) at (4.3,0) {\textbf{Argument 1}};
\node[ex] (e1) at (7.6,0) {exemple(s)};
\node[bloc] (a2) at (4.3,-1.9) {\textbf{Argument 2}};
\node[ex] (e2) at (7.6,-1.9) {exemple(s)};
\node[bloc] (a3) at (4.3,-3.8) {\textbf{Argument 3}};
\node[ex] (e3) at (7.6,-3.8) {exemple(s)};
\node[bloc, fill=popgreenL, draw=popgreen, align=center] (ccl) at (0,-3.8){\textbf{Conclusion}\\ bilan + ouverture};
\draw[fleche] (intro) -- (a1);
\draw[fleche] (a1) -- (e1);
\draw[fleche] (a1.south) -- (a2.north);
\draw[fleche] (a2) -- (e2);
\draw[fleche] (a2.south) -- (a3.north);
\draw[fleche] (a3) -- (e3);
\draw[fleche] (a3.west) -- (ccl.east);
\node[font=\footnotesize\itshape, popmuted] at (3.8,0.75) {Développement};
\draw[popline, dashed] (3.3,0.6) -- (3.3,-4.4);
\end{tikzpicture}
\end{center}
\legende{Schéma en blocs de la structure argumentative : l'introduction annonce la thèse, le développement enchaîne les arguments illustrés d'exemples, la conclusion tire le bilan.}
\end{popfigure}

\begin{aretenir}[L'essentiel]
Tout texte argumentatif bien construit suit le schéma : \textbf{introduction} (sujet + thèse) $\rightarrow$ \textbf{développement} (arguments + exemples) $\rightarrow$ \textbf{conclusion} (bilan + ouverture). Repérer ces trois parties, c'est déjà faire la moitié du travail de contraction.
\end{aretenir}

\courssub{Les plans argumentatifs fréquents}

Les trois parties se combinent selon des plans types qu'il faut savoir reconnaître, car le correcteur du Probatoire et du Baccalauréat technique les attend aussi dans vos propres productions (dissertation, commentaire).

\textbf{1. Le plan dialectique : thèse / antithèse / synthèse.} On examine d'abord une position (thèse), puis la position opposée (antithèse), enfin on dépasse l'opposition par une position équilibrée ou nuancée (synthèse). Ce plan convient aux sujets controversés : « Faut-il interdire les réseaux sociaux aux élèves ? », « L'argent fait-il le bonheur ? ».

\textbf{2. Le plan analytique : constat / causes / conséquences / solutions.} On part d'un fait observé (constat), on en cherche les causes, on en mesure les conséquences, puis on propose des solutions. C'est le plan typique des articles de presse sur un problème de société : l'insécurité routière, le chômage des jeunes, la déforestation, la cybercriminalité.

\textbf{3. Le plan thématique.} On examine successivement les différents aspects d'une même question (aspect économique, aspect social, aspect culturel, aspect politique), sans opposition frontale ni problème unique à résoudre. Ce plan convient aux sujets descriptifs, explicatifs ou aux essais de société.

\begin{popfigure}
\begin{center}
\begin{tikzpicture}[
 etape/.style={draw=popblue, fill=popblueL, rounded corners=2pt, text width=2.5cm, minimum height=0.62cm, align=center, font=\footnotesize},
 titre/.style={font=\small\bfseries, popdark},
 fleche/.style={-{Stealth[length=2mm]}, thick, popdark}]
% Plan dialectique
\node[titre] at (0,0.4) {Plan dialectique};
\node[etape] (d1) at (0,-0.4) {Thèse};
\node[etape] (d2) at (0,-1.35) {Antithèse};
\node[etape, fill=popgreenL, draw=popgreen] (d3) at (0,-2.3) {Synthèse};
\draw[fleche] (d1) -- (d2); \draw[fleche] (d2) -- (d3);
% Plan analytique
\node[titre] at (4.6,0.4) {Plan analytique};
\node[etape] (a1) at (4.6,-0.2) {Constat};
\node[etape] (a2) at (4.6,-0.95) {Causes};
\node[etape] (a3) at (4.6,-1.7) {Conséquences};
\node[etape, fill=popgreenL, draw=popgreen] (a4) at (4.6,-2.45) {Solutions};
\draw[fleche] (a1) -- (a2); \draw[fleche] (a2) -- (a3); \draw[fleche] (a3) -- (a4);
% Plan thématique
\node[titre] at (9.2,0.4) {Plan thématique};
\node[etape] (t0) at (9.2,-0.2) {Sujet};
\node[etape] (t1) at (7.6,-1.35) {Aspect 1};
\node[etape] (t2) at (9.2,-1.35) {Aspect 2};
\node[etape] (t3) at (10.8,-1.35) {Aspect 3};
\node[etape, fill=popgreenL, draw=popgreen] (t4) at (9.2,-2.45) {Conclusion};
\draw[fleche] (t0) -- (t1); \draw[fleche] (t0) -- (t2); \draw[fleche] (t0) -- (t3);
\draw[fleche] (t1) -- (t4); \draw[fleche] (t2) -- (t4); \draw[fleche] (t3) -- (t4);
\end{tikzpicture}
\end{center}
\legende{Les trois plans argumentatifs fréquents : le plan dialectique oppose deux positions avant de les dépasser ; le plan analytique remonte du constat aux solutions ; le plan thématique examine les aspects d'un sujet.}
\end{popfigure}

\begin{savaistu}[Le plan dialectique aux examens camerounais]
Le plan dialectique (thèse / antithèse / synthèse) est très utilisé dans les sujets de dissertation du Probatoire et du Baccalauréat camerounais, notamment lorsque le sujet est formulé comme une question ouverte : « L'école est-elle la seule voie de la réussite ? » Savoir reconnaître ce plan dans un texte à contracter, c'est aussi apprendre à le construire dans vos propres dissertations.
\end{savaistu}

\courssub{Repérer les articulations du texte}

Comment le lecteur repère-t-il concrètement le passage d'une partie à l'autre ? Deux indices se combinent.

\textbf{1. Les connecteurs logiques (ou opérateurs argumentatifs).} Ce sont des mots-charnières qui annoncent la fonction de ce qui suit :
\begin{itemize}
\item pour introduire : \emph{d'abord, en premier lieu, il convient de rappeler, premièrement} ;
\item pour ajouter un argument : \emph{ensuite, de plus, en outre, par ailleurs, de surcroît, deuxième} ;
\item pour opposer / nuancer : \emph{mais, cependant, néanmoins, en revanche, au contraire, si... toutefois} (souvent le passage de la thèse à l'antithèse) ;
\item pour illustrer : \emph{par exemple, ainsi, c'est le cas de, notamment, comme en témoigne} ;
\item pour conclure : \emph{donc, enfin, en définitive, en somme, c'est pourquoi, ultimement}.
\end{itemize}

\textbf{2. Les paragraphes (alinéas).} Dans un texte bien construit, chaque paragraphe développe une seule idée directrice. Le changement de paragraphe signale donc presque toujours le passage à une nouvelle étape du raisonnement. C'est pourquoi la première chose à faire face à un texte à contracter est de \cle{numéroter les paragraphes}.

\begin{methode}[Dégager la structure d'un texte en quatre étapes]
\begin{enumerate}
\item \textbf{Numéroter} les paragraphes du texte (1, 2, 3...).
\item \textbf{Souligner} la phrase-clé (ou phrase-topic) de chaque paragraphe, celle qui en résume l'idée directrice.
\item \textbf{Regrouper} les paragraphes qui développent la même idée ou la même étape (introduction, argument 1, argument 2, conclusion...).
\item \textbf{Nommer} chaque groupe en une phrase : on obtient le plan du texte, c'est-à-dire sa structure argumentative.
\end{enumerate}
\end{methode}

\courssub{Hiérarchiser les idées : principales et secondaires}

Toutes les phrases d'un texte n'ont pas la même importance. La contraction de texte repose entièrement sur cette distinction.

\begin{definition}[Idées principales et idées secondaires]
Les \cle{idées principales} sont celles qui font avancer le raisonnement : la thèse, les arguments, la conclusion. Elles constituent l'\cle{ossature} du texte.
Les \cle{idées secondaires} sont celles qui illustrent, précisent, ornent ou répètent : les \cle{exemples}, les \cle{répétitions} (reformulations internes), les \cle{digressions} (hors-sujet momentané), les \cle{citations} d'autorité et les \cle{statistiques détaillées} (chiffres bruts).
\end{definition}

Un bon test : si l'on supprime la phrase, le raisonnement tient-il encore debout ? Si oui, c'était une idée secondaire. Si le raisonnement s'effondre ou devient incompréhensible, c'était une idée principale.

\begin{attention}[Erreur fréquente à la contraction]
L'erreur la plus fréquente consiste à garder les \textbf{exemples}, les \textbf{citations} et les \textbf{chiffres précis} dans la contraction. Ils doivent être \textbf{supprimés} ou \textbf{généralisés} (ex. : « de nombreux élèves » au lieu de « 15 élèves au lycée de Douala ») : ce sont des idées secondaires qui ne font qu'illustrer l'argument. De même, on élimine les répétitions (la même idée exprimée deux fois) et les digressions. Ne gardez que la thèse, les arguments et la conclusion.
\end{attention}

\textbf{Technique du soulignage.} À la première lecture, soulignez au crayon la phrase-clé de chaque paragraphe. Elle se trouve souvent au début du paragraphe (phrase d'annonce / phrase-topic) ou à la fin (phrase de bilan / transition). Les phrases qui commencent par « par exemple », « ainsi », « c'est le cas », « selon X » ne sont presque jamais des phrases-clés.

\textbf{Technique du résumé de paragraphe (marges).} Après le soulignage, formulez dans la marge, en une phrase et avec \textbf{vos propres mots}, l'idée de chaque paragraphe. Ces formulations marginales deviendront la matière première de votre contraction. C'est ici que s'applique la \cle{reformulation} (synonymie, antonymie, changement de structure syntaxique), que nous étudierons dans la section suivante.

\begin{exemplebox}[Le quart du texte : repère chiffré de la contraction]
Un texte de 400 mots se contracte généralement en \textbf{100 mots environ} (soit \textbf{25\%}, tolérance habituelle \textbf{20\% à 30\%}). Cela représente en pratique 4 à 5 idées essentielles : la thèse, deux ou trois arguments et la conclusion. Si votre contraction dépasse largement le tiers, c'est que vous avez conservé des idées secondaires ; si elle est beaucoup plus courte, c'est que vous avez perdu une idée principale.
\end{exemplebox}

\courssub{Application guidée : découpage d'un article de presse camerounais}

Appliquons la méthode à un article de presse reconstitué, de type éditorial, sur un sujet d'actualité nationale.

\begin{exoresolu}[Découper un article en grandes parties et identifier le plan]
\textbf{Énoncé.} Voici le plan paragraphe par paragraphe d'un article intitulé « Le téléphone portable à l'école : fléau ou outil ? » (article de 8 paragraphes, environ 400 mots). Pour chaque paragraphe, on a noté sa phrase-clé résumée :
\begin{itemize}
\item §1 : De plus en plus d'élèves camerounais possèdent un téléphone portable, y compris en classe.
\item §2 : La question se pose donc : le portable est-il un danger ou un outil pour l'école ?
\item §3 : D'abord, le portable distrait les élèves et fait chuter la concentration.
\item §4 : Par exemple, dans un lycée de Douala, un professeur a confisqué 15 téléphones en une seule journée.
\item §5 : Ensuite, le portable favorise la triche pendant les évaluations.
\item §6 : Cependant, le portable peut aussi servir à faire des recherches documentaires.
\item §7 : Ainsi, certaines applications gratuites donnent accès à des dictionnaires et des cours.
\item §8 : En définitive, il faut encadrer l'usage du portable plutôt que l'interdire totalement.
\end{itemize}
Découpez cet article en grandes parties, identifiez le plan argumentatif utilisé et proposez une contraction d'environ 100 mots.
\tcblower
\textbf{Solution détaillée.}

\textbf{Étape 1 — Regroupement des paragraphes (Macro-structure).}
\begin{itemize}
\item §1 et §2 : Présentation du sujet (phénomène de société) + Problématique + Thèse implicite (le débat est ouvert). $\rightarrow$ \textbf{Introduction}.
\item §3 à §7 : Enchaînement d'arguments contraires. $\rightarrow$ \textbf{Développement}.
\item §8 : Bilan + Proposition équilibrée. $\rightarrow$ \textbf{Conclusion}.
\end{itemize}

\textbf{Étape 2 — Analyse fine du développement (Micro-structure).}
\begin{itemize}
\item §3 et §5 : Deux arguments \textbf{contre} le portable (distraction, triche). Le §4 (exemple Douala) illustre §3. $\rightarrow$ \textbf{Thèse} : « Le portable est un fléau ».
\item §6 et §7 : Le connecteur « Cependant » (§6) signale le retournement. Deux arguments \textbf{pour} le portable (recherche, apps éducatives). Le §7 (exemple apps) illustre §6. $\rightarrow$ \textbf{Antithèse} : « Le portable est un outil ».
\item §8 : « En définitive » + proposition d'encadrement (ni interdiction totale, ni liberté totale). $\rightarrow$ \textbf{Synthèse}.
\end{itemize}

\textbf{Étape 3 — Identification du plan.}
L'article suit un \textbf{plan dialectique} : Thèse (§3-§5) / Antithèse (§6-§7) / Synthèse (§8).

\textbf{Étape 4 — Contraction proposée (environ 100 mots).}
« De plus en plus d'élèves camerounais utilisent le téléphone portable en classe, ce qui pose la question de sa place à l'école. Certains y voient un fléau : il distrait les élèves et favorise la triche lors des évaluations. D'autres soulignent au contraire qu'il peut devenir un outil pédagogique, offrant un accès aux dictionnaires et aux cours en ligne. En définitive, plutôt que de l'interdire totalement, ce qui serait inefficace, il convient d'en encadrer strictement l'usage par une réglementation claire. »

\textbf{Observations :} L'exemple du lycée de Douala (§4) et la liste des applications (§7) ont été \textbf{supprimés} (idées secondaires). Les arguments sont \textbf{reformulés} (ex. : « fait chuter la concentration » $\rightarrow$ « distrait » ; « confisqué 15 téléphones » $\rightarrow$ supprimé/généralisé).
\end{exoresolu}

\begin{exercice}[Entraînement : Analyse structurelle et contraction]
Choisissez un article d'opinion (éditorial, tribune, chronique) dans un journal camerounais récent (\emph{Mutations, Cameroon Tribune, Le Messager, etc.}) d'une longueur de 350 à 450 mots.
\begin{enumerate}
\item Numérotez les paragraphes et soulignez la phrase-clé de chacun.
\item Identifiez l'introduction, le développement (en distinguant les arguments principaux) et la conclusion.
\item Relevez trois connecteurs logiques et précisez leur fonction (addition, opposition, illustration, conclusion).
\item Identifiez le plan argumentatif (dialectique, analytique, thématique) et justifiez votre choix.
\item Listez deux idées secondaires (exemples, chiffres, citations) à éliminer.
\item Rédigez la contraction de l'article (100 mots $\pm$ 10\%) en appliquant la reformulation (vos propres mots).
\end{enumerate}
\end{exercice}

\begin{aretenir}[Bilan de la section : Structure et Hiérarchisation]
\begin{itemize}
\item \textbf{Structure type :} Introduction (Sujet + Thèse) $\rightarrow$ Développement (Arguments + Exemples) $\rightarrow$ Conclusion (Bilan + Ouverture).
\item \textbf{Plans types :} Dialectique (Thèse/Antithèse/Synthèse), Analytique (Constat/Causes/Conséquences/Solutions), Thématique (Aspects multiples).
\item \textbf{Repérage :} Numéroter paragraphes $\rightarrow$ Souligner phrases-clés $\rightarrow$ Regrouper $\rightarrow$ Nommer.
\item \textbf{Hiérarchisation :} Garder \textbf{idées principales} (thèse, arguments, conclusion) ; Supprimer \textbf{idées secondaires} (exemples, détails, répétitions, digressions).
\item \textbf{Calibrage :} Cible = 25\% du texte original (tolérance 20-30\%).
\item \textbf{Prochaine étape :} La \textbf{reformulation} (synonymie, antonymie, nominalisation, passage actif/passif) pour écrire la contraction avec ses propres mots.
\end{itemize}
\end{aretenir}

\coursec{Les tonalités : comique, dramatique, satirique, tragique}

Dans la section précédente, nous avons appris à repérer la \cle{structure argumentative} d'un texte : sa thèse, ses arguments, ses exemples. Mais un texte argumentatif ne convainc pas seulement par la logique : il touche aussi le lecteur par l'\cle{émotion}. Un auteur peut défendre la même idée en faisant rire, en indignant ou en angoissant. Cette « couleur émotionnelle » du texte, c'est ce que l'on appelle la \cle{tonalité} (ou \cle{registre}). Savoir la repérer est indispensable pour comprendre l'intention de l'auteur, et donc pour formuler correctement les idées essentielles d'un texte à résumer.

\courssub{Qu'est-ce que la tonalité ?}

\begin{definition}[La tonalité (ou registre)]
La \cle{tonalité} d'un texte est l'\textbf{impression dominante} qu'il cherche à produire chez le lecteur : le \textbf{rire} (tonalité comique), l'\textbf{indignation} ou la moquerie (tonalité satirique), l'\textbf{émotion} et la tension (tonalité dramatique), l'\textbf{angoisse} face au destin (tonalité tragique). Elle résulte du choix des mots, des situations, du rythme des phrases et des figures de style.
\end{definition}

L'intuition est simple : quand on écoute un conteur, on comprend vite s'il veut nous amuser, nous dénoncer un abus, nous émouvoir ou nous glacer d'effroi. Le texte écrit fonctionne de la même manière. La tonalité n'est donc pas une idée du texte : c'est la \textbf{manière} dont les idées sont présentées.

\begin{attention}[Tonalité et idées essentielles]
La tonalité ne fait pas partie des idées essentielles : \textbf{elle ne se résume pas}. En revanche, elle aide à comprendre l'\textbf{intention} de l'auteur (divertir, dénoncer, émouvoir). C'est pourquoi on la signale \textbf{brièvement dans l'introduction du résumé}, par exemple : « Dans ce texte à la tonalité satirique, l'auteur dénonce... ».
\end{attention}

\courssub{La tonalité comique}

\begin{definition}[Tonalité comique]
La tonalité \cle{comique} est celle qui provoque le \textbf{rire} ou le \textbf{sourire} du lecteur. Son but premier est de \textbf{divertir}, de créer la détente.
\end{definition}

On distingue traditionnellement quatre procédés comiques :

\begin{itemize}
\item le \textbf{comique de mots} : jeux de mots, calembours, mots inventés, quiproquos (un personnage comprend un mot pour un autre), répétitions amusantes ;
\item le \textbf{comique de gestes} : gestes maladroits, chutes, mimiques exagérées, disputes physiques sans gravité ;
\item le \textbf{comique de situation} : quiproquo, malentendu, déguisement, arrivée inattendue d'un personnage au mauvais moment ;
\item le \textbf{comique de caractère} : un défaut exagéré devient ridicule (l'avare, le vaniteux, le paresseux, le menteur invétéré).
\end{itemize}

\begin{exemplebox}[Comique de caractère]
Chez Molière, dans \emph{L'Avare}, Harpagon cache son argent dans son jardin et soupçonne tout le monde de vouloir le voler. Son avarice, poussée à l'extrême, devient ridicule et fait rire : c'est du comique de caractère.
\end{exemplebox}

\courssub{La tonalité satirique}

\begin{definition}[Tonalité satirique]
La tonalité \cle{satirique} utilise le rire comme une \textbf{arme} : elle \textbf{dénonce} les défauts, les vices et les abus d'une personne, d'un groupe ou d'une société, par la \textbf{moquerie}, l'\textbf{ironie} (dire le contraire de ce qu'on pense) et l'\textbf{exagération} (hyperbole). Elle est toujours liée à une \textbf{critique sociale}.
\end{definition}

Le raisonnement de l'auteur satirique est le suivant : si je dis simplement « ce comportement est mauvais », le lecteur peut m'ignorer ; mais si je le rends \textbf{ridicule}, le lecteur rit, et en riant il \textbf{juge} et condamne ce comportement. Le rire devient alors un instrument de critique.

\begin{attention}[Piège fréquent : confondre comique et satirique]
Le \textbf{comique} fait rire pour \textbf{divertir}. Le \textbf{satirique} fait rire pour \textbf{dénoncer}. Dans la satire, il y a toujours une \textbf{cible} (un vice, un abus, un pouvoir) et une \textbf{intention critique}. Demande-toi toujours : « L'auteur veut-il seulement m'amuser, ou veut-il que je condamne quelque chose ? »
\end{attention}

\begin{savaistu}[Wole Soyinka, satiriste engagé]
Wole Soyinka, dramaturge nigérian, \textbf{prix Nobel de littérature en 1986}, utilise dans \emph{Le Lion et la Perle} la satire pour critiquer à la fois le \textbf{traditionalisme aveugle} (incarné par Baroka, le chef qui refuse tout progrès) et la \textbf{modernité superficielle} (incarnée par Lakunle, l'instituteur qui imite l'Occident sans le comprendre). Le rire du spectateur devient ainsi un jugement sur les deux excès.
\end{savaistu}

\begin{exemplebox}[Satire dans \emph{Le Lion et la Perle}]
Lakunle refuse de payer le « prix de la fiancée » pour épouser Sidi, au nom du « progrès », mais il récite des idées occidentales sans les comprendre vraiment. Soyinka le rend ridicule : c'est une satire de la modernité mal digérée. Inversement, Baroka, le vieux chef rusé, est moqué pour son attachement intéressé aux traditions.
\end{exemplebox}

\courssub{La tonalité dramatique}

\begin{definition}[Tonalité dramatique]
La tonalité \cle{dramatique} naît d'un \textbf{conflit intense} entre des personnages ou des forces opposées. Elle crée chez le lecteur une \textbf{tension}, une forte \textbf{émotion}, et le sentiment de l'\textbf{impuissance} des personnages face à leur situation. Les passions (amour, haine, jalousie, ambition) y sont poussées à l'extrême.
\end{definition}

Dans un texte dramatique, le lecteur se demande avec inquiétude : « Comment le personnage va-t-il s'en sortir ? » Le conflit peut encore, en principe, se résoudre : c'est ce qui distingue le dramatique du tragique. Les indices typiques sont les exclamations, les questions angoissées, les phrases courtes et saccadées, les champs lexicaux de la souffrance et de la lutte.

\begin{exemplebox}[Conflit dramatique]
Dans \emph{Le Lion et la Perle}, la lutte entre Lakunle et Baroka pour conquérir Sidi crée une tension dramatique : deux visions du monde s'affrontent, et Sidi est prise entre les deux. Le spectateur ressent le poids de ce choix.
\end{exemplebox}

\courssub{La tonalité tragique}

\begin{definition}[Tonalité tragique]
La tonalité \cle{tragique} repose sur la \textbf{fatalité} : un \textbf{destin inéluctable} écrase le héros, quoi qu'il fasse. La \textbf{mort} ou la catastrophe est \textbf{annoncée d'avance} (prophéties, présages), et le personnage lutte en vain contre une \textbf{force supérieure} (les dieux, le destin, une malédiction). Elle inspire la \textbf{pitié} et la \textbf{terreur}.
\end{definition}

La différence essentielle avec le dramatique : dans le tragique, l'issue est \textbf{certaine} et \textbf{fatale}. Le héros ne peut pas échapper à son destin ; plus il lutte, plus il s'en rapproche. Chez Sophocle, Œdipe fuit pour éviter la prédiction de l'oracle — et c'est en fuyant qu'il l'accomplit.

\courssub{Tableau récapitulatif et échelle d'intensité}

\begin{popfigure}
\begin{center}
\begin{tikzpicture}[scale=1]
\draw[fill=popblueL, draw=popblue] (0,2.4) rectangle (3,3.4);
\node[font=\bfseries\small] at (1.5,2.9) {Tonalité};
\draw[fill=popblueL, draw=popblue] (3,2.4) rectangle (6.6,3.4);
\node[font=\bfseries\small] at (4.8,2.9) {Procédés};
\draw[fill=popblueL, draw=popblue] (6.6,2.4) rectangle (10.2,3.4);
\node[font=\bfseries\small] at (8.4,2.9) {Effet sur le lecteur};
\draw[fill=popblueL, draw=popblue] (10.2,2.4) rectangle (13.8,3.4);
\node[font=\bfseries\small] at (12,2.9) {Champ lexical};
% ligne comique
\draw[draw=popline] (0,1.2) rectangle (3,2.4);
\node[font=\small\bfseries, text=popgreen] at (1.5,1.8) {Comique};
\draw[draw=popline] (3,1.2) rectangle (6.6,2.4);
\node[font=\scriptsize, text width=3.3cm, align=center] at (4.8,1.8) {jeux de mots, quiproquo, gestes, caractère};
\draw[draw=popline] (6.6,1.2) rectangle (10.2,2.4);
\node[font=\scriptsize, text width=3.3cm, align=center] at (8.4,1.8) {rire, sourire, détente};
\draw[draw=popline] (10.2,1.2) rectangle (13.8,2.4);
\node[font=\scriptsize, text width=3.3cm, align=center] at (12,1.8) {ridicule, absurde, drôle};
% ligne satirique
\draw[draw=popline] (0,0) rectangle (3,1.2);
\node[font=\small\bfseries, text=poporange] at (1.5,0.6) {Satirique};
\draw[draw=popline] (3,0) rectangle (6.6,1.2);
\node[font=\scriptsize, text width=3.3cm, align=center] at (4.8,0.6) {ironie, moquerie, exagération};
\draw[draw=popline] (6.6,0) rectangle (10.2,1.2);
\node[font=\scriptsize, text width=3.3cm, align=center] at (8.4,0.6) {rire + indignation, jugement critique};
\draw[draw=popline] (10.2,0) rectangle (13.8,1.2);
\node[font=\scriptsize, text width=3.3cm, align=center] at (12,0.6) {vice, hypocrisie, abus};
% ligne dramatique
\draw[draw=popline] (0,-1.2) rectangle (3,0);
\node[font=\small\bfseries, text=poppurple] at (1.5,-0.6) {Dramatique};
\draw[draw=popline] (3,-1.2) rectangle (6.6,0);
\node[font=\scriptsize, text width=3.3cm, align=center] at (4.8,-0.6) {conflit, tension, phrases saccadées};
\draw[draw=popline] (6.6,-1.2) rectangle (10.2,0);
\node[font=\scriptsize, text width=3.3cm, align=center] at (8.4,-0.6) {émotion, inquiétude, suspense};
\draw[draw=popline] (10.2,-1.2) rectangle (13.8,0);
\node[font=\scriptsize, text width=3.3cm, align=center] at (12,-0.6) {souffrance, lutte, passion};
% ligne tragique
\draw[draw=popline] (0,-2.4) rectangle (3,-1.2);
\node[font=\small\bfseries, text=popink] at (1.5,-1.8) {Tragique};
\draw[draw=popline] (3,-2.4) rectangle (6.6,-1.2);
\node[font=\scriptsize, text width=3.3cm, align=center] at (4.8,-1.8) {fatalité, présages, destin annoncé};
\draw[draw=popline] (6.6,-2.4) rectangle (10.2,-1.2);
\node[font=\scriptsize, text width=3.3cm, align=center] at (8.4,-1.8) {pitié, terreur, angoisse};
\draw[draw=popline] (10.2,-2.4) rectangle (13.8,-1.2);
\node[font=\scriptsize, text width=3.3cm, align=center] at (12,-1.8) {destin, mort, malédiction};
\end{tikzpicture}
\end{center}
\legende{Tableau des quatre tonalités : procédés, effets sur le lecteur et champs lexicaux typiques.}
\end{popfigure}

\begin{popfigure}
\begin{center}
\begin{tikzpicture}[scale=1]
\draw[->, very thick, popmuted] (0,0) -- (13,0) node[right, font=\small] {intensité émotionnelle};
\foreach \x in {0,1,...,12} \draw[popmuted] (\x,0.08) -- (\x,-0.08);
% comique
\fill[popgreen] (1.5,0) circle (0.14);
\node[above, font=\small\bfseries, text=popgreen] at (1.5,0.15) {Comique};
\node[below, font=\scriptsize, text width=2.6cm, align=center] at (1.5,-0.15) {rire léger, divertissement};
% satirique
\fill[poporange] (5,0) circle (0.14);
\node[above, font=\small\bfseries, text=poporange] at (5,0.15) {Satirique};
\node[below, font=\scriptsize, text width=2.6cm, align=center] at (5,-0.15) {rire critique, indignation};
% dramatique
\fill[poppurple] (8.5,0) circle (0.14);
\node[above, font=\small\bfseries, text=poppurple] at (8.5,0.15) {Dramatique};
\node[below, font=\scriptsize, text width=2.6cm, align=center] at (8.5,-0.15) {tension, émotion forte};
% tragique
\fill[popink] (12,0) circle (0.14);
\node[above, font=\small\bfseries, text=popink] at (12,0.15) {Tragique};
\node[below, font=\scriptsize, text width=2.6cm, align=center] at (12,-0.15) {angoisse, fatalité, mort};
\end{tikzpicture}
\end{center}
\legende{Frise graduée de l'intensité émotionnelle : du rire léger du comique à l'angoisse du tragique.}
\end{popfigure}

\courssub{Comment repérer la tonalité d'un texte ?}

\begin{methode}[Identifier la tonalité en trois étapes]
\begin{enumerate}
\item \textbf{Relever le champ lexical dominant} : les mots du rire et du ridicule (comique), du vice et de l'hypocrisie (satirique), de la souffrance et du conflit (dramatique), du destin et de la mort (tragique).
\item \textbf{Observer les figures d'insistance et le rythme} : hyperboles, répétitions, exclamations, ironie (dire le contraire de ce qu'on pense), phrases courtes et saccadées (tension) ou longues et solennelles (tragique).
\item \textbf{Examiner la situation des personnages} : quiproquo ou ridicule (comique), personnage moqué pour ses vices (satirique), conflit intense à l'issue incertaine (dramatique), héros écrasé par une force supérieure (tragique).
\end{enumerate}
\end{methode}

\begin{exoresolu}[Identifier la tonalité d'un passage]
\textbf{Énoncé.} Lis ce passage inspiré de \emph{Le Lion et la Perle} : « Baroka, le vieux Lion, fit courir le bruit de son impuissance. Tout le village le plaignait ; Sadiku elle-même vint annoncer à Sidi que le chef n'était plus un homme. Mais c'était un piège : le Lion, loin d'être faible, préparait sa plus belle ruse. » Identifie la tonalité dominante et justifie ta réponse par deux indices.

\tcblower
\textbf{Solution détaillée.}
\begin{enumerate}
\item \textbf{Champ lexical} : les mots « ruse », « piège », « fit courir le bruit » appartiennent au lexique de la \textbf{tromperie} et du \textbf{faux-semblant}.
\item \textbf{Situation} : Baroka se fait passer pour impuissant alors qu'il est en pleine possession de ses moyens : c'est un \textbf{comique de situation} (tout le monde est trompé, le lecteur seul comprend).
\item \textbf{Intention} : Soyinka ne fait pas seulement rire : il \textbf{dénonce} la ruse et l'hypocrisie du pouvoir traditionnel. Il y a donc aussi une visée \textbf{satirique}.
\end{enumerate}
\textbf{Conclusion} : la tonalité est \textbf{comique} (situation trompeuse, rire du lecteur) avec une dimension \textbf{satirique} (critique des faux-semblants et de la ruse du chef). Dans l'introduction d'un résumé, on écrirait : « Dans cet extrait à la tonalité comique et satirique, Soyinka montre comment Baroka use de la ruse pour parvenir à ses fins. »
\end{exoresolu}

\begin{aretenir}[L'essentiel sur les tonalités]
\begin{itemize}
\item La \textbf{tonalité} est l'émotion dominante qu'un texte cherche à produire chez le lecteur.
\item \textbf{Comique} : faire rire pour divertir (mots, gestes, situations, caractères).
\item \textbf{Satirique} : faire rire pour \textbf{dénoncer} (ironie, moquerie, exagération).
\item \textbf{Dramatique} : créer la tension par un conflit intense à l'issue incertaine.
\item \textbf{Tragique} : montrer un héros écrasé par un destin inéluctable.
\item Pour le résumé : la tonalité \textbf{ne se résume pas}, mais on la \textbf{signale dans l'introduction} car elle éclaire l'intention de l'auteur.
\end{itemize}
\end{aretenir}

\begin{exercice}[À toi de jouer]
Pour chacune des situations suivantes, indique la tonalité dominante et justifie ta réponse en une phrase :
\begin{enumerate}
\item Un personnage glisse sur une peau de banane devant tout le village.
\item Un écrivain décrit un ministre « si honnête qu'il n'a jamais refusé une enveloppe ».
\item Deux frères se disputent violemment l'héritage de leur père mourant.
\item Un oracle annonce au héros qu'il tuera son père ; malgré toutes ses précautions, la prédiction s'accomplit.
\end{enumerate}
\end{exercice}

\begin{corrige}[Exercice : les tonalités]
\begin{enumerate}
\item \textbf{Comique} (comique de gestes : la chute ridicule fait rire, sans visée critique).
\item \textbf{Satirique} : l'ironie (« si honnête » signifie en réalité le contraire) dénonce la corruption ; le rire sert ici à critiquer un vice.
\item \textbf{Dramatique} : conflit intense entre deux personnages, tension forte, issue incertaine.
\item \textbf{Tragique} : le destin annoncé s'accomplit malgré les efforts du héros ; la fatalité écrase le personnage.
\end{enumerate}
\end{corrige}

Nous savons désormais reconnaître la couleur émotionnelle d'un texte. Cette compétence sera immédiatement utile dans la section suivante, consacrée à la lecture méthodique du \emph{Lion et la Perle}, où le rire, la satire et la tension dramatique se mêlent constamment.

\coursec{Lecture méthodique du Lion et la Perle (n°4 et n°5)}

Après avoir identifié les \cle{tonalités} (comique, dramatique, satirique, tragique) qui colorent le discours théâtral, nous passons maintenant à leur application concrète sur l'œuvre au programme. La \cle{lecture méthodique} est l'outil qui permet de passer de l'impression de lecture à l'analyse rigoureuse, étape indispensable avant toute \cle{contraction de texte} ou \cle{commentaire}.

\begin{definition}[La lecture méthodique]
La \textbf{lecture méthodique} est une lecture guidée qui suit un plan d'analyse fixe pour explorer un extrait littéraire en profondeur. Elle structure l'explication de texte en cinq temps obligatoires :
\begin{enumerate}
    \item \textbf{La situation de l'extrait :} localisation dans l'œuvre (acte, scène), personnages présents, contexte immédiat (ce qui précède et ce qui suit).
    \item \textbf{Le thème :} sujet principal abordé dans l'extrait (le conflit, le débat, l'émotion).
    \item \textbf{Les procédés d'écriture :} figures de style, registres, lexique, syntaxe, didascalies, organisation du dialogue.
    \item \textbf{La tonalité :} couleur affective dominante (comique, satirique, dramatique, tragique, etc.) et ses effets sur le spectateur/lecteur.
    \item \textbf{Le sens de l'extrait :} portée générale, idées essentielles, message de l'auteur, lien avec l'ensemble de l'œuvre.
\end{enumerate}
Cette démarche garantit une analyse exhaustive et ordonnée, attendue à l'examen du Probatoire.
\end{definition}

\begin{aretenir}[Rappel : Le texte argumentatif au théâtre]
Un texte argumentatif théâtral ne se limite pas à une suite de répliques ; il met en scène un \textbf{conflit d'idées} porté par des personnages-antagonistes.
\begin{itemize}
    \item \textbf{Thèse :} Position défendue par un personnage (ex: Baroka \emph{pour} la tradition dynamisée).
    \item \textbf{Antithèse :} Position adverse (ex: Lakunle \emph{contre} la coutume, \emph{pour} la modernité occidentale).
    \item \textbf{Arguments :} Preuves logiques, affectives ou d'autorité (proverbes, faits, comparaisons).
    \item \textbf{Connecteurs logiques :} \emph{donc, or, mais, en effet, par conséquent, cependant} (marquent l'enchaînement).
    \item \textbf{Synthèse :} Issue du conflit (victoire, compromis, échec) qui révèle le message de l'auteur.
\end{itemize}
La \cle{contraction de texte} consiste à extraire cette charpente argumentative en gommant l'anecdotique (didascalies, politesses, répétitions).
\end{aretenir}

\courssub{Rappel de l'œuvre : Wole Soyinka, \emph{Le Lion et la Perle}}

\begin{savaistu}[Wole Soyinka et la symbolique du titre]
\emph{Le Lion et la Perle} (\emph{The Lion and the Jewel}, 1959) est une comédie du dramaturge nigérian \textbf{Wole Soyinka}, prix Nobel de littérature 1986. L'action se déroule à \textbf{Ilujinle}, village yoruba imaginaire. Le titre résume le duel central :
\begin{itemize}
    \item \textbf{Le Lion :} \cle{Baroka}, le Bale (chef du village), figure de la \cle{tradition}, de la force, de la ruse et de l'autorité patriarcale.
    \item \textbf{La Perle :} \cle{Sidi}, la belle du village, symbole de la beauté, de la jeunesse, de l'enjeu du conflit et du pouvoir de choix (apparent).
    \item \textbf{Le tiers exclu du titre mais central :} \cle{Lakunle}, l'instituteur « moderne », représentant la \cle{modernité} occidentale, le rationalisme maladroit et le rejet des coutumes.
\end{itemize}
La pièce met en scène le choc des cultures non pas comme une opposition binaire simple, mais comme une négociation complexe où la ruse traditionnelle finit par triompher de la modernité théorique.
\end{savaistu}

\begin{popfigure}
\begin{center}
\begin{tikzpicture}[
    nodeStyle/.style={draw=popdark, thick, fill=popcream, rounded corners=8pt, minimum width=3.5cm, minimum height=2cm, align=center, font=\small},
    arrowStyle/.style={->, >=Stealth, thick, popblue},
    valueStyle/.style={font=\footnotesize\itshape, text=poppurple, align=center}
]
% Nœuds personnages
\node[nodeStyle, align=center] (Baroka) at (0,0){\textbf{BAROKA}\\(Le Lion)\\\textit{Le Bale d'Ilujinle}};
\node[nodeStyle, align=center] (Sidi) at (5,4.33){\textbf{SIDI}\\(La Perle)\\\textit{La Belle du village}};
\node[nodeStyle, align=center] (Lakunle) at (10,0){\textbf{LAKUNLE}\\\textit{L'instituteur}};
% Flèches relations
\draw[arrowStyle, bend left=15] (Baroka) -- node[valueStyle, above left, align=center]{Conquête\\Ruse\\Séduction} (Sidi);
\draw[arrowStyle, bend right=15] (Sidi) -- node[valueStyle, below left, align=center]{Choix final\\Triomphe de la tradition} (Baroka);
\draw[arrowStyle, bend left=15] (Lakunle) -- node[valueStyle, above right, align=center]{Discours moderne\\Refus de la dot\\Séduction ratée} (Sidi);
\draw[arrowStyle, bend right=15] (Sidi) -- node[valueStyle, below right, align=center]{Rejet\\Moquerie} (Lakunle);
\draw[arrowStyle, dashed, poporange] (Baroka) -- node[valueStyle, below, align=center]{Rivalité\\Tradition vs Modernité} (Lakunle);
% Valeurs associées
\node[valueStyle, align=center] at (0,-2.2){\textbf{VALEURS :}\\Tradition, Autorité,\\Ruse, Polygamie,\\Continuité};
\node[valueStyle, align=center] at (5,6.5){\textbf{ENJEU :}\\Beauté, Jeunesse,\\Fécondité,\\Honneur du village};
\node[valueStyle, align=center] at (10,-2.2){\textbf{VALEURS :}\\Modernité, Monogamie,\\Rationalisme,\\Rejet coutumes,\\Maladresse};
\end{tikzpicture}
\end{center}
\legende{Schéma triangulaire des relations et valeurs dans \emph{Le Lion et la Perle}. Le trait plein = relation directe ; le trait pointillé = opposition idéologique.}
\end{popfigure}

\courssub{Lecture méthodique n°4 : Extrait à tonalité comique et satirique}

\textbf{Situation de l'extrait.}\par
Il s'agit généralement de la \textbf{Scène 1 (Matin)} ou du début de la \textbf{Scène 2 (Midi)}. Lakunle tente de convaincre Sidi de l'épouser \textbf{sans payer la dot}, arguant que cette coutume est « barbare », « humiliante » et « archaïque ». Baroka, quant à lui, n'est pas physiquement présent dans la première partie de cette scène, mais son ombre plane ; la rumeur de sa ruse (la feinte impuissance) commence à circuler. L'extrait type pour cette lecture porte sur les \emph{prétentions modernistes de Lakunle} face au bon sens paysan de Sidi.

\textbf{Thème.}\par
Le \textbf{conflit tradition / modernité} vu sous l'angle du mariage et de la dot. Lakunle veut imposer une vision occidentale (mariage d'amour, monogamie, égalité) à une société régie par la coutume yoruba.

\textbf{Procédés d'écriture et analyse du comique/satirique.}\par
Le comique naît du \cle{décalage} entre le discours de Lakunle et la réalité vécue par Sidi.

\begin{exemplebox}[Procédés comiques chez Lakunle]
\begin{itemize}
    \item \textbf{L'exagération (hyperbole) et le pédantisme lexical :} Lakunle utilise un vocabulaire savant, abstrait, déconnecté du concret (« \emph{une coutume barbare, sauvage, rétrograde} », « \emph{acheter une femme} »). Il parle en théoricien, pas en amant.
    \item \textbf{Le quiproquo / L'incompréhension mutuelle :} Sidi répond sur le terrain concret (honneur, famille, valeur marchande de sa beauté), Lakunle sur le terrain idéologique. Ils ne s'entendent pas.
    \item \textbf{L'ironie dramatique :} Le spectateur sait (ou devine) que Lakunle est ridicule. Ses gestes (mimer la modernité, refuser de porter le seau d'eau) le rendent grotesque.
    \item \textbf{La satire sociale :} Soyinka tourne en dérision le « moderne » déraciné qui singe l'Occidental sans en avoir la culture ni les moyens, et qui méprise sa propre culture.
\end{itemize}
\end{exemplebox}

\textbf{Tonalité.}\par
\textbf{Comique} (bouffonnerie de Lakunle, situations cocasses) et \textbf{Satirique} (critique acide de l'assimilationnisme bête, de la modernité vide de sens). Le rire est un outil de lucidité : il dénonce l'aliénation culturelle.

\textbf{Sens de l'extrait.}\par
Lakunle perd la partie rhétorique. Sidi refuse de devenir « la risée du village » (une femme sans dot n'a pas de valeur sociale). L'extrait pose l'échec de la modernité \emph{imposée par le haut} et sans ancrage populaire.

\courssub{Lecture méthodique n°5 : Extrait à tonalité dramatique ou tragique}

\textbf{Situation de l'extrait.}\par
Généralement la \textbf{Scène 3 (Soir)} : la célèbre scène de la \textbf{ruse de Baroka} dans sa chambre. Baroka a fait courir le bruit de son impuissance. Sidi, triomphante, vient le narguer. Baroka la piège par la conversation, l'alcool, la lutte (force physique) et la séduction intellectuelle. Lakunle attend dehors, impuissant, réduit au rôle de spectateur. La scène se termine par la victoire de Baroka (la nuit d'amour) et le choix de Sidi au matin (Scène 4/Épilogue) : elle épouse Baroka.

\textbf{Thème.}\par
\textbf{La ruse comme pouvoir}, la \cle{condition de la femme africaine} (objet d'échange puis sujet de choix), le \textbf{triomphe de la tradition vivante} sur la modernité morte.

\textbf{Procédés d'écriture et analyse du dramatique/tragique.}\par

\begin{exemplebox}[Du comique au tragique : le basculement]
\begin{itemize}
    \item \textbf{La tension dramatique (suspense) :} Sidi entre dans la gueule du lion. Le spectateur sait la ruse (ironie dramatique), Sidi non. Peur pour l'héroïne.
    \item \textbf{Le registre pathétique / tragique :} La perte de virginité de Sidi (hors mariage coutumier) la « souille » aux yeux de la tradition. Elle n'a plus le choix : « \emph{Je suis une femme, je ne peux pas revenir en arrière} ». La contrainte sociale pèse de tout son poids.
    \item \textbf{La force du corps et de la parole :} Baroka domine par le corps (lutte, âge mûr viril) et l'esprit (proverbes, vision d'avenir : le timbre-poste, le progrès \emph{maîtrisé}). Il incarne une tradition \emph{dynamique}, non figée.
    \item \textbf{L'effacement de Lakunle :} Il devient une silhouette pathétique, prêt à épouser Sidi « sans dot » \emph{après} sa défloration par Baroka, par lâcheté ou opportunisme. La modernité est démasquée comme un leurre.
\end{itemize}
\end{exemplebox}

\textbf{Tonalité.}\par
\textbf{Dramatique} (confrontation à haut risque, enjeu vital pour l'héroïne) et \textbf{Tragique} (destin scellé, irreversibilité de l'acte, sacrifice de la liberté individuelle sur l'autel de la coutume). La fin n'est pas « heureuse » au sens occidental, elle est \emph{acceptée} comme l'ordre des choses.

\textbf{Sens de l'extrait.}\par
Baroka gagne parce qu'il \textbf{intègre le changement} (le timbre, la machine) sans rompre le lien social (la dot, le respect des aînés). Sidi choisit la sécurité du connu et l'honneur rétabli. Lakunle reste à la marge.

\courssub{Dégagement des idées essentielles : Synthèse comparative}

Pour réussir la \cle{contraction de texte} ou le \cle{résumé}, il faut isoler les \cle{idées forces} qui traversent les deux extraits. Le tableau ci-dessous structure la pensée argumentative des deux rivaux.

\begin{popfigure}
\begin{center}
\begin{tabularx}{\linewidth}{|>{\bfseries}l|X|X|}
\hline
\rowcolor{popblueL}
\textbf{Thème / Argument} & \textbf{BAROKA (Tradition vivante)} & \textbf{LAKUNLE (Modernité livresque)} \\
\hline
\textbf{Vision du mariage} & Alliance entre familles, continuité lignagère, honneur social. La dot scelle le pacte. & Contrat d'amour entre deux individus, monogamie, égalité des sexes. La dot = achat/vente. \\
\hline
\textbf{Rapport à la femme} & Épouse = partenaire dans le foyer, mère des fils, respectée si dot payée. Polygamie assumée. & « Camarade », « égales », mais traité en enfant incapable de comprendre la « vérité » moderne. \\
\hline
\textbf{Argument sur la dot} & « \emph{La dot prouve la valeur de l'homme et l'honneur de la fille.} » Outil de cohésion sociale. & « \emph{La dot est une coutume barbare, humiliante, qui traite la femme comme une marchandise.} » \\
\hline
\textbf{Rapport au progrès} & \textbf{Intégration sélective :} « \emph{Je veux le timbre-poste, la machine, mais à mon service.} » Progrès maîtrisé. & \textbf{Imitation servile :} Rejet total du passé, singe l'Occidental (costume, langage, gestes) sans en avoir la technique. \\
\hline
\textbf{Moyen d'action} & \textbf{Ruse, parole, force physique, autorité légitime.} Joue sur les codes culturels (proverbes, rites). & \textbf{Discours abstrait, leçons de morale, gestuelle ridicule.} Ne maîtrise pas les codes locaux. \\
\hline
\textbf{Issue (Extrait 5)} & \textbf{Victoire :} Conquiert Sidi, assure sa descendance, modernise le village à sa façon. & \textbf{Échec :} Perd Sidi, humilié, réduit à accepter les miettes (épouser la « souillée » sans dot). \\
\hline
\end{tabularx}
\end{center}
\legende{Tableau comparatif des arguments de Baroka et Lakunle. Ce tableau est une grille de lecture pour le résumé : chaque ligne peut devenir une phrase de la contraction.}
\end{popfigure}

\begin{aretenir}[Idées essentielles à retenir pour l'examen]
Trois idées majeures structurent l'œuvre et doivent apparaître dans tout résumé ou commentaire :
\begin{enumerate}
    \item \textbf{Le conflit Tradition / Modernité n'est pas manichéen :} La tradition (Baroka) n'est pas figée, elle sait absorber le progrès. La modernité (Lakunle) est présentée comme aliénante quand elle nie les racines.
    \item \textbf{La ruse (l'intelligence pratique) triomphe de la théorie :} Baroka use de \emph{métis} (intelligence rusée, adaptation) ; Lakunle use de \emph{logos} abstrait, déconnecté du réel.
    \item \textbf{La condition féminine :} Sidi est l'enjeu, mais elle a une \cle{agency} (capacité d'agir) finale : elle \emph{choisit} Baroka pour sauver son honneur et sa place dans la communauté. Elle n'est pas seulement une victime.
\end{enumerate}
\end{aretenir}

\courssub{La reformulation des idées essentielles : Méthode complète du résumé}

Passer de l'analyse à la \cle{reformulation} (étape clé de la contraction) demande de gommer l'anecdotique pour ne garder que l'argumentatif.

\begin{methode}[Résumer un extrait de théâtre en 5 étapes (Entraînement Probatoire)]
\begin{enumerate}
    \item \textbf{Identifier le conflit central :} Qui veut quoi ? Qui s'oppose à qui ? (Ex: Lakunle veut épouser Sidi \emph{sans dot} ; Sidi refuse \emph{par honneur} ; Baroka veut Sidi \emph{par ruse}).
    \item \textbf{Lister les arguments de chaque camp :} Utiliser le tableau comparatif (dot, progrès, femme, pouvoir). Noter les verbes d'action : \emph{condamner, refuser, imposer, ruser, conquérir, choisir}.
    \item \textbf{Supprimer le narratif (le « récit ») :} Ne pas écrire : « Lakunle arrive, il porte un seau, il parle, Sidi rit, Baroka entre... ». Écrire : « Lakunle \textbf{condamne} la dot ; Sidi \textbf{défend} l'honneur ; Baroka \textbf{manœuvre} et \textbf{vainc}. »
    \item \textbf{Reformuler par la synonymie et l'antonymie :} Remplacer les mots du texte par des termes génériques ou des concepts (voir boîte Synonymie/Antonymie). \emph{« Coutume barbare » $\rightarrow$ « tradition honnie » ; « Perle » $\rightarrow$ « enjeu symbolique ».}
    \item \textbf{Structurer le résumé :} Introduction (sujet + thèse) $\rightarrow$ Développement (arguments groupés par idée, non par personnage) $\rightarrow$ Conclusion (issue + sens). Respecter la limite de lignes (souvent 10-15 lignes au Probatoire, 5 lignes à l'entraînement).
\end{enumerate}
\end{methode}

\begin{attention}[Piège majeur : Raconter au lieu de Résumer]
\textbf{Erreur fréquente :} Produire un « résumé-récit » (chronologie des actions : « Il dit ça, elle répond ça, il fait ça »).
\textbf{Exigence du Bac/Probatoire :} Un \textbf{résumé d'idées} (structure argumentative : « Le texte oppose deux visions du mariage : l'une traditionaliste qui valorise la dot comme gage d'honneur, l'autre moderniste qui la rejette comme aliénation. La ruse du traditionaliste l'emporte... »).
\textbf{Astuce :} Vérifiez vos verbes. Verbes d'action scénique (\emph{entrer, dire, rire, frapper}) = Récit (Mauvais). Verbes intellectuels/argumentatifs (\emph{affirmer, nier, démontrer, contester, illustrer, triompher}) = Résumé (Bon).
\end{attention}

\begin{exoresolu}[Entraînement : Résumer l'argument de Baroka en 5 lignes maximum]
\textbf{Consigne :} À partir de la scène 3 (soir), résumer en 5 lignes l'argumentation de Baroka pour conquérir Sidi.

\textbf{Texte support (résumé de la scène) :} Baroka reçoit Sidi. Il feint la vieillesse et l'impuissance. Il discute de progrès (timbre-poste), boit avec elle, la lutte et la vainc physiquement. Il lui révèle qu'il a payé la dot pour une autre femme le matin même (provocation). Il la séduit par sa vision d'un village modernisé mais respectueux des coutumes. Sidi cède.

\textbf{Application de la méthode :}
\begin{enumerate}
    \item \textbf{Conflit :} Baroka vs Sidi (séduction / pouvoir).
    \item \textbf{Arguments :} 1. Force virile intacte (lutte). 2. Intelligence moderne (timbre/progrès). 3. Respect de la tradition (dot payée pour une autre = puissance). 4. Vision d'avenir partagé.
    \item \textbf{Reformulation (Synonymie) :} Feinte $\rightarrow$ Stratégie ; Lutte $\rightarrow$ Preuve de vitalité ; Timbre $\rightarrow$ Symbole de modernité maîtrisée ; Dot $\rightarrow$ Preuve de statut.
\end{enumerate}

\tcblower

\textbf{Proposition de correction (5 lignes) :}
\begin{quote}
\textbf{Baroka use d'une ruse savante pour conquérir Sidi.} \textbf{Il feint l'impuissance pour la désarmer, puis prouve sa virilité par la lutte.} \textbf{Il oppose une modernité maîtrisée (le timbre-poste) à la modernité stérile de Lakunle.} \textbf{En affichant sa puissance coutumière (paiement de la dot), il impose son autorité.} \textbf{Sidi, vaincue physiquement et séduite intellectuellement, accepte l'alliance traditionnelle.}
\end{quote}

\textbf{Analyse de la correction :} Pas de « Il dit », « Elle répond ». Des noms de concepts (\emph{ruse, virilité, modernité maîtrisée, autorité, alliance}) et des verbes d'enchaînement logique (\emph{use, feint, prouve, oppose, impose, accepte}).
\end{exoresolu}

\begin{exoresolu}[Entraînement : Résumer l'argument de Lakunle en 5 lignes maximum]
\textbf{Consigne :} Résumer en 5 lignes la position de Lakunle face à Sidi (Scène 1).

\tcblower

\textbf{Proposition de correction :}
\begin{quote}
\textbf{Lakunle refuse le mariage coutumier et la dot, qu'il juge barbares et humiliants pour la femme.} \textbf{Il propose une union moderne, fondée sur l'amour, l'égalité et la monogamie, à l'occidentale.} \textbf{Il tente d'imposer sa vision par le discours théorique et le rejet des rites (refus de porter l'eau).} \textbf{Sidi oppose le bon sens paysan et l'honneur social : sans dot, elle n'est « rien » aux yeux du village.} \textbf{Lakunle, incompris et ridicule, échoue à convertir Sidi à sa modernité abstraite.}
\end{quote}
\end{exoresolu}

\courssub{Outils lexicaux : Synonymie et Antonymie au service de la reformulation}

La qualité du résumé dépend de la richesse lexicale. Le programme impose de maîtriser la \cle{synonymie} (mots de sens voisin) et l'\cle{antonymie} (mots de sens opposé) pour éviter la citation servile (copier-coller) et montrer l'appropriation du sens.

\begin{propriete}[Synonymie et Antonymie dans \emph{Le Lion et la Perle}]
\begin{itemize}
    \item \textbf{Synonymie (reformulation) :}
    \begin{itemize}
        \item \emph{Dot} $\leftrightarrow$ \textbf{prestation nuptiale, gage d'honneur, prix de la mariée, coutume matrimoniale}.
        \item \emph{Barbare / Rétrograde} (Lakunle) $\leftrightarrow$ \textbf{archaïque, dépassée, coutumière, traditionnelle, ancestrale}.
        \item \emph{Ruse / Fourberie} (Baroka) $\leftrightarrow$ \textbf{stratégie, manœuvre, habileté, intelligence pratique, métis}.
        \item \emph{Perle} $\leftrightarrow$ \textbf{joyau, trésor, beauté, enjeu, symbole}.
    \end{itemize}
    \item \textbf{Antonymie (structuration du conflit) :}
    \begin{itemize}
        \item \textbf{Tradition} $\neq$ \textbf{Modernité} (mais attention : Baroka = Tradition \textbf{modernisatrice}).
        \item \textbf{Parole/Acte} (Baroka) $\neq$ \textbf{Théorie/Geste vide} (Lakunle).
        \item \textbf{Honneur collectif} (Sidi/Baroka) $\neq$ \textbf{Liberté individuelle} (Lakunle).
        \item \textbf{Virilité/Force} $\neq$ \emph{Impuissance/Féminité} (feinte de Baroka / réalité de Lakunle).
    \end{itemize}
\end{itemize}
\textbf{Exercice d'application :} Dans votre résumé, remplacez « Lakunle dit que la dot est barbare » par « Lakunle \textbf{condamne} la \textbf{prestation nuptiale} qu'il juge \textbf{archaïque} ».
\end{propriete}

\courssub{Lien avec la contraction de texte : L'extrait d'examen}

\begin{methode}[De la lecture méthodique à la contraction (Épreuve type Probatoire)]
L'examen propose souvent un extrait de \emph{Le Lion et la Perle} (ou texte apparenté) comme support de \cle{contraction de texte} (résumé en 1/3 ou 1/4, ou nombre de lignes imposé).
\begin{enumerate}
    \item \textbf{Lecture 1 (Rapide) :} Identifier l'extrait (Quelle scène ? Quels personnages ? Quel conflit ?).
    \item \textbf{Lecture 2 (Active - Méthodique) :} Appliquer le plan : Situation $\rightarrow$ Thème $\rightarrow$ Procédés $\rightarrow$ Tonalité $\rightarrow$ Sens. \textbf{Surligner les mots-clés arguments} (pas les didascalies, pas les répliques anecdotiques).
    \item \textbf{Rédaction du brouillon :} Appliquer la \textbf{Méthode en 5 étapes} (ci-dessus). Grouper les arguments par \emph{idées forces} (ex: Idée 1: La dot ; Idée 2: Le progrès ; Idée 3: La femme). Utiliser la \textbf{synonymie}.
    \item \textbf{Comptage et Ajustement :} Compter les mots/lignes. Couper les exemples, les détails, les adjectifs qualificatifs. Garder la \textbf{charpente logique} (Connecteurs logiques : \emph{donc, or, mais, en effet, par conséquent}).
    \item \textbf{Copie propre :} Vérifier l'orthographe, la syntaxe, l'absence de « je », le style impersonnel.
\end{enumerate}
\end{methode}

\begin{exemplebox}[Exemple de sujet d'entraînement type Probatoire]
\textbf{Sujet :} « Vous résumerez en \textbf{12 lignes maximum} (environ 100 mots) l'extrait suivant de \emph{Le Lion et la Perle} (Scène 3, soirée), en faisant ressortir l'opposition entre la ruse de Baroka et la naïveté de Sidi. »

\textbf{Pistes de correction (Idées forces à inclure) :}
\begin{enumerate}
    \item Situation : Baroka reçoit Sidi qui vient le narguer (croit à son impuissance).
    \item Ruse 1 : Feinte de vieillesse / discussion sur le progrès (timbre) = modernité assumée.
    \item Ruse 2 : Lutte physique = preuve de virilité / domination corporelle.
    \item Ruse 3 : Révélation dot payée = puissance sociale / provocation.
    \item Basculement : Sidi, désarmée, séduite, perd sa virginité.
    \item Conséquence : Triomphe de la tradition vivante / Échec de la modernité livresque (Lakunle absent/impuissant).
\end{enumerate}
\end{exemplebox}

\courssub{Compte rendu, remédiation et entraînement final}

\begin{aretenir}[Check-list de révision pour la Section 4 (Probatoire)]
\begin{itemize}
    \item [ ] Je sais situer les extraits n°4 (Matin/Midi : Lakunle/Sidi) et n°5 (Soir : Baroka/Sidi) dans la pièce.
    \item [ ] Je connais les 5 étapes de la \textbf{lecture méthodique} (Situation, Thème, Procédés, Tonalité, Sens).
    \item [ ] Je sais identifier et nommer les \textbf{procédés comiques} (quiproquo, exagération, ironie, satire) chez Lakunle.
    \item [ ] Je sais identifier les \textbf{procédés dramatiques/tragiques} (suspense, ironie dramatique, pathétique, irreversibilité) chez Baroka/Sidi.
    \item [ ] Je suis capable de remplir le \textbf{tableau comparatif} Baroka / Lakunle (Dot, Progrès, Femme, Pouvoir, Issue).
    \item [ ] Je sais \textbf{reformuler} (synonymie/antonymie) sans raconter l'histoire.
    \item [ ] Je maîtrise la \textbf{méthode du résumé de théâtre} (Conflit $\rightarrow$ Arguments $\rightarrow$ Idées $\rightarrow$ Structure $\rightarrow$ Comptage).
    \item [ ] Je connais le \textbf{piège « Récit vs Résumé »} et je sais utiliser les verbes argumentatifs.
\end{itemize}
\end{aretenir}

\begin{exercice}[Entraînement autonome - À rendre pour la séance 30 (Remédiation)]
\begin{enumerate}
    \item \textbf{Analyse de tonalité :} Relevez dans l'extrait 4 (Scène 1) trois répliques de Lakunle relevant du \textbf{comique de mots} et trois relevant de la \textbf{satire}. Justifiez chaque choix.
    \item \textbf{Étude de la ruse :} Dans l'extrait 5 (Scène 3), identifiez les \textbf{trois temps de la séduction-ruse} de Baroka (Conversation $\rightarrow$ Lutte $\rightarrow$ Révélation). Quelle est la fonction de l'alcool (vin de palme) dans la scène ?
    \item \textbf{Contraction guidée :} Résumez en \textbf{8 lignes exactement} le conflit central de la pièce tel qu'il apparaît dans les deux extraits étudiés. Utilisez au moins 4 paires synonymiques/antonymiques vues en cours.
    \item \textbf{Expression orale :} Préparez une intervention de 2 minutes : « \emph{En quoi Baroka est-il un « moderne » plus efficace que Lakunle ?} » (Utilisez le tableau comparatif).
\end{enumerate}
\end{exercice}

\begin{corrige}[Exercice n°3 - Contraction guidée (Proposition)]
\textbf{Le Lion et la Perle oppose deux visions du mariage à Ilujinle.} \textbf{Lakunle, chantre d'une modernité importée, condamne la dot comme une aliénation barbare et prône l'union libre et monogame.} \textbf{Sidi, gardienne de l'honneur collectif, rejette cette modernité stérile qui la priverait de valeur sociale.} \textbf{Baroka, figure de la tradition dynamique, use de ruse et de force pour conquérir la Perle.} \textbf{Il intègre le progrès technique (timbre) sans rompre le pacte coutumier (dot payée).} \textbf{Vaincue physiquement, séduite intellectuellement, Sidi choisit l'alliance traditionnelle, assurant la continuité du village.} \textbf{La modernité théorique s'efface devant l'intelligence pratique du pouvoir ancestral.}
\par
\textit{(8 lignes / env. 95 mots. Verbes argumentatifs : oppose, condamne, prône, rejette, use, intègre, rompe, choisit, assure, s'efface. Synonymies : aliénation/barbarie, stérile/vide, dynamique/vivante, pacte/coutume).}
\end{corrige}

\coursec{Synonymie et antonymie : outils de la reformulation}

Dans la section précédente, notre lecture méthodique du \emph{Lion et la Perle} nous a montré que Wole Soyinka oppose, à travers les personnages de Baroka et Lakunle, deux visions du monde : la tradition et la modernité. Remarquons-le : pour parler de cette opposition sans recopier le texte, nous avons déjà utilisé des mots de sens contraire. C'est exactement le travail qui nous attend maintenant. Pour résumer un texte, il ne suffit pas de repérer les idées essentielles : il faut encore les \cle{reformuler}, c'est-à-dire les exprimer avec ses propres mots. Les deux grands outils de la reformulation sont la \cle{synonymie} et l'\cle{antonymie}. Cette section vous donne la maîtrise complète de ces outils, indispensables à l'épreuve de contraction de texte du Probatoire et du Baccalauréat technique.

\courssub{La synonymie : des mots de sens proche}

\textbf{Intuition.} Imaginez que vous racontiez à un ami une histoire que vous avez lue. Vous ne récitez pas le texte mot à mot : vous le redites avec votre vocabulaire. Là où l'auteur écrivait « école », vous dites peut-être « établissement scolaire » ; là où il écrivait « ruse », vous dites « astuce ». Le sens reste le même, mais les mots changent. Vous venez, sans le savoir, d'utiliser des synonymes.

\begin{definition}[Synonymie et antonymie]
Deux mots sont \cle{synonymes} lorsqu'ils ont un sens très proche et peuvent se substituer l'un à l'autre dans un contexte donné. Deux mots sont \cle{antonymes} lorsqu'ils ont des sens opposés.
\end{definition}

\begin{exemplebox}[Premières paires de synonymes]
\begin{itemize}
\item \emph{école} / \emph{établissement scolaire} : « L'école du quartier est rénovée » = « L'établissement scolaire du quartier est rénové ».
\item \emph{ruse} / \emph{astuce} : dans le \emph{Lion et la Perle}, Baroka use de ruse ; on peut dire aussi qu'il use d'astuce.
\item \emph{entreprendre} / \emph{commencer} ; \emph{nombreux} / \emph{beaucoup de} ; \emph{région} / \emph{province} (selon le contexte).
\end{itemize}
\end{exemplebox}

\begin{attention}[Les synonymes ne sont jamais parfaitement interchangeables]
Deux synonymes ont un sens \emph{proche}, rarement \emph{identique}. Il existe presque toujours une nuance de registre de langue, d'intensité ou de contexte d'emploi. Comparez :
\begin{itemize}
\item \emph{mourir} : registre courant, neutre ;
\item \emph{décéder} : registre soutenu, administratif (« le défunt est décédé à l'hôpital ») ;
\item \emph{trépasser} : registre littéraire ou religieux, plus solennel.
\end{itemize}
Dans un résumé, choisissez le synonyme adapté au ton du texte de départ : un article de journal sérieux appelle « décéder » ou « mourir », jamais « crever » (registre familier).
\end{attention}

\textbf{Les trois types de nuances entre synonymes.}

\begin{center}
\begin{tabularx}{\linewidth}{|l|X|X|}
\hline
\rowcolor{popblueL} \textbf{Nuance} & \textbf{Explication} & \textbf{Exemple} \\
\hline
Registre de langue & familier / courant / soutenu & \emph{boulot} / \emph{travail} / \emph{emploi} \\
\hline
Intensité & le sens est plus ou moins fort & \emph{aimer} / \emph{adorer} ; \emph{fâché} / \emph{furieux} \\
\hline
Contexte d'emploi & le mot ne convient qu'à certaines situations & \emph{décéder} (administratif) / \emph{mourir} (tous contextes) \\
\hline
\end{tabularx}
\end{center}

\courssub{L'antonymie : des mots de sens contraire}

\textbf{Intuition.} Pour dire qu'une idée s'oppose à une autre, le français dispose de mots de sens contraire. Dans le \emph{Lion et la Perle}, l'opposition centrale peut se formuler ainsi : \emph{tradition} / \emph{modernité}. Baroka \emph{accepte} les coutumes ; Lakunle les \emph{rejette}. Ces paires de contraires sont précieuses pour le résumé, car elles permettent de condenser une phrase négative en une phrase affirmative, plus courte.

\begin{exemplebox}[Paires d'antonymes utiles]
\begin{itemize}
\item \emph{tradition} / \emph{modernité} ; \emph{accepter} / \emph{rejeter} ;
\item \emph{progrès} / \emph{régression} ; \emph{unité} / \emph{division} ;
\item \emph{richesse} / \emph{pauvreté} ; \emph{paix} / \emph{conflit}.
\end{itemize}
\end{exemplebox}

\textbf{Les antonymes formés par préfixation.} Le français fabrique de nombreux contraires en ajoutant un préfixe négatif devant un mot de base. Connaître ces préfixes multiplie votre stock d'antonymes :

\begin{center}
\begin{tabularx}{\linewidth}{|l|X|X|}
\hline
\rowcolor{popblueL} \textbf{Préfixe} & \textbf{Valeur} & \textbf{Exemples} \\
\hline
in- / im- / il- / ir- & négation & possible / \emph{im}possible ; lisible / \emph{il}lisible ; régulier / \emph{ir}régulier \\
\hline
dé- / dés- & inversion, suppression & accord / \emph{dés}accord ; faire / \emph{dé}faire \\
\hline
mé- & mauvais, erreur & content / \emph{mé}content ; connaissance / \emph{mé}connaissance \\
\hline
anti- & opposition & social / \emph{anti}social \\
\hline
\end{tabularx}
\end{center}

\begin{propriete}[L'antonymie avec négation : un outil de condensation]
Toute phrase de la forme « il n'est pas + antonyme » peut être remplacée par « il est + adjectif positif » (et inversement). Cette transformation raccourcit la phrase sans en changer le sens :
\begin{itemize}
\item « il n'est pas hostile » = « il est favorable » ;
\item « le projet n'est pas impossible » = « le projet est possible » ;
\item « elle ne refuse pas le dialogue » = « elle accepte le dialogue ».
\end{itemize}
\end{propriete}

\courssub{Les autres procédés de reformulation}

La synonymie et l'antonymie ne suffisent pas toujours. Le bon résumeur dispose de quatre procédés complémentaires.

\textbf{1. La conversion grammaticale (changement de catégorie).} On transforme un nom en verbe, un adjectif en nom, etc. C'est le procédé le plus puissant pour alléger une phrase :
\begin{itemize}
\item « la \emph{construction} des routes » $\rightarrow$ « \emph{construire} des routes » (nom $\rightarrow$ verbe) ;
\item « la \emph{modernisation} du pays » $\rightarrow$ « \emph{moderniser} le pays » ;
\item « le pays est \emph{prospère} » $\rightarrow$ « la \emph{prospérité} du pays » (adjectif $\rightarrow$ nom).
\end{itemize}

\textbf{2. Le changement de voix.} On passe de la voix active à la voix pronominale ou passive (et inversement) pour changer la formulation : « le gouvernement construit des routes » $\rightarrow$ « des routes se construisent ».

\textbf{3. La fusion de phrases.} Deux phrases courtes qui expriment une même idée peuvent être fondues en une seule : « Le texte oppose Baroka et Lakunle. Baroka défend la tradition. » $\rightarrow$ « Le texte oppose Lakunle à Baroka, défenseur de la tradition. »

\textbf{4. La suppression des expansions inutiles.} Les adjectifs descriptifs, les exemples, les répétitions et les compléments accessoires disparaissent du résumé : « le vaste et magnifique palais du sultan, construit au siècle dernier » $\rightarrow$ « le palais du sultan ».

\begin{popfigure}
\begin{tikzpicture}[
  grow=right,
  level distance=4.6cm,
  level 1/.style={sibling distance=2.6cm},
  edge from parent/.style={draw=popmuted,thick},
  every node/.style={font=\small}
]
\node[draw=popdark,fill=popblueL,rounded corners,align=center,font=\small\bfseries] {Procédés de\\ reformulation}
  child { node[draw=popblue,fill=popblueL!50,rounded corners,align=center] {Fusion de phrases\\ \emph{deux phrases $\rightarrow$ une seule}} }
  child { node[draw=popgreen,fill=popgreenL,rounded corners,align=center] {Conversion grammaticale\\ \emph{la construction $\rightarrow$ construire}} }
  child { node[draw=poporange,fill=poporangeL,rounded corners,align=center] {Antonymie + négation\\ \emph{pas hostile $\rightarrow$ favorable}} }
  child { node[draw=poppurple,fill=poppurpleL,rounded corners,align=center] {Synonymie\\ \emph{ruse $\rightarrow$ astuce}} };
\end{tikzpicture}
\legende{Arbre des quatre grands procédés de reformulation, avec un exemple par branche. Au résumé, on combine généralement plusieurs procédés dans une même phrase.}
\end{popfigure}

\begin{methode}[Reformuler une phrase en quatre étapes]
\begin{enumerate}
\item \textbf{Identifier} les mots-clés de la phrase (ceux qui portent l'idée essentielle).
\item \textbf{Remplacer} ces mots-clés par des synonymes de même registre, ou par un antonyme accompagné d'une négation.
\item \textbf{Transformer} les noms en verbes quand cela allège la phrase (conversion grammaticale).
\item \textbf{Supprimer} les expansions inutiles (adjectifs descriptifs, exemples, compléments accessoires) et vérifier que le sens est rigoureusement identique.
\end{enumerate}
\end{methode}

\begin{exemplebox}[Exemple chiffré : gagner neuf mots sans perdre le sens]
Phrase de départ (16 mots) : « Le gouvernement camerounais a entrepris la construction de nombreuses infrastructures routières dans les dix régions. »\\
Reformulation (7 mots) : « Le Cameroun construit de nombreuses routes partout. »\\
Procédés utilisés : \emph{gouvernement camerounais} $\rightarrow$ \emph{le Cameroun} (synonymie) ; \emph{a entrepris la construction} $\rightarrow$ \emph{construit} (conversion nom $\rightarrow$ verbe) ; \emph{infrastructures routières} $\rightarrow$ \emph{routes} (simplification) ; \emph{dans les dix régions} $\rightarrow$ \emph{partout} (condensation).
\end{exemplebox}

\begin{attention}[L'erreur capitale : la copie]
Recopier des phrases entières du texte, même en en supprimant quelques mots, n'est pas un résumé : c'est de la \cle{copie}, lourdement sanctionnée à l'examen. Le correcteur surligne les « collages » (suites de quatre mots ou plus identiques au texte) et pénalise chacun d'eux. Règle d'or : le résumé doit être \textbf{entièrement reformulé} ; seuls les noms propres et les termes techniques sans équivalent peuvent être conservés.
\end{attention}

\begin{savaistu}[La règle du quart]
Aux épreuves de contraction de texte du Probatoire et du Baccalauréat technique, la norme la plus fréquente est la \cle{règle du quart} : le résumé doit mesurer environ un quart de la longueur du texte original, avec une marge de tolérance de plus ou moins 10 \%. Pour un texte de 400 mots, visez donc 100 mots (entre 90 et 110). Comptez vos mots et indiquez leur nombre à la fin de votre copie : c'est une exigence de l'examen.
\end{savaistu}

\courssub{Tableau de paires utiles au résumé}

\begin{center}
\begin{tabularx}{\linewidth}{|l|X|X|}
\hline
\rowcolor{popblueL} \textbf{Mot du texte} & \textbf{Synonyme(s)} & \textbf{Antonyme (+ négation)} \\
\hline
école & établissement scolaire & --- \\
\hline
ruse & astuce, habileté & franchise (sans ruse) \\
\hline
tradition & coutume, héritage & modernité \\
\hline
accepter & approuver, admettre & rejeter, refuser \\
\hline
entreprendre & commencer, lancer & abandonner \\
\hline
nombreux & beaucoup de, de multiples & rare (peu nombreux) \\
\hline
progrès & avancée, développement & régression, recul \\
\hline
richesse & prospérité, fortune & pauvreté, misère \\
\hline
hostile & défavorable, opposé & favorable (pas hostile) \\
\hline
possible & réalisable, envisageable & impossible \\
\hline
\end{tabularx}
\end{center}

\courssub{Exercices guidés de reformulation}

\begin{exoresolu}[Reformulation d'une phrase de presse camerounaise]
Énoncé. Reformulez la phrase suivante, extraite d'un article camerounais, sans en changer le sens : « Les autorités ont décidé la modernisation des marchés ruraux afin d'améliorer les revenus des producteurs. » (17 mots)
\tcblower
Solution détaillée. Mots-clés : \emph{autorités}, \emph{modernisation}, \emph{améliorer}, \emph{revenus}. Conversion grammaticale : \emph{la modernisation} $\rightarrow$ \emph{moderniser}. Synonymie : \emph{autorités} $\rightarrow$ \emph{l'État} ; \emph{améliorer les revenus} $\rightarrow$ \emph{augmenter les gains}. Reformulation (12 mots) : « L'État veut moderniser les marchés ruraux pour augmenter les gains des producteurs. » Le sens est identique, aucune suite de quatre mots du texte n'est conservée.
\end{exoresolu}

\begin{exercice}[Reformuler cinq phrases d'un article camerounais]
Reformulez chacune des phrases suivantes sans en changer le sens, en utilisant au moins deux procédés par phrase. Indiquez le nombre de mots avant et après.
\begin{enumerate}
\item « Le gouvernement a annoncé la construction de nouveaux hôpitaux dans les chefs-lieux de département. »
\item « Les jeunes diplômés ne trouvent pas facilement un emploi dans les grandes villes du pays. »
\item « La tradition n'est pas hostile au progrès, affirme le chef traditionnel de la localité. »
\item « Les élèves des lycées techniques reçoivent une formation pratique qui facilite leur insertion professionnelle. »
\item « Malgré de nombreuses difficultés, les transporteurs continuent d'assurer la liaison entre Yaoundé et Douala. »
\end{enumerate}
\end{exercice}

\begin{corrige}[Exercice : reformulation de cinq phrases]
Voici des reformulations possibles (d'autres sont acceptables si le sens est respecté et si aucune phrase n'est copiée) :
\begin{enumerate}
\item « L'État va construire de nouveaux hôpitaux dans les chefs-lieux de département. » (12 mots contre 15 ; synonymie \emph{gouvernement} $\rightarrow$ \emph{État}, conversion \emph{construction} $\rightarrow$ \emph{construire}.)
\item « Les jeunes diplômés peinent à trouver un emploi dans les grandes villes. » (12 mots ; antonymie + négation : \emph{ne trouvent pas facilement} $\rightarrow$ \emph{peinent à trouver} ; suppression de \emph{du pays}.)
\item « Pour le chef de la localité, la tradition est favorable au progrès. » (12 mots ; antonymie + négation : \emph{n'est pas hostile} $\rightarrow$ \emph{est favorable} ; conversion \emph{affirme} $\rightarrow$ \emph{pour}.)
\item « La formation pratique des lycées techniques facilite l'insertion professionnelle des élèves. » (12 mots ; fusion et réorganisation, suppression de \emph{reçoivent une}.)
\item « Les transporteurs assurent toujours la liaison Yaoundé--Douala en dépit des difficultés. » (11 mots ; synonymie \emph{malgré} $\rightarrow$ \emph{en dépit de}, \emph{continuent d'assurer} $\rightarrow$ \emph{assurent toujours}.)
\end{enumerate}
\end{corrige}

\begin{aretenir}[L'essentiel de la section]
\begin{itemize}
\item Deux mots sont \cle{synonymes} s'ils ont un sens très proche ; ils sont \cle{antonymes} s'ils ont des sens opposés. Les synonymes se distinguent par le registre, l'intensité et le contexte.
\item Les préfixes \emph{in-/im-}, \emph{dé-/dés-}, \emph{mé-} fabriquent des antonymes.
\item Quatre procédés de reformulation : synonymie, antonymie avec négation, conversion grammaticale, fusion de phrases et suppression des expansions.
\item Interdiction absolue : copier des phrases du texte. Le résumé est entièrement reformulé et respecte la règle du quart.
\end{itemize}
\end{aretenir}

\coursec{La reformulation des idées essentielles : méthode complète du résumé}

Dans la section précédente, nous avons étudié la \cle{synonymie} et l'\cle{antonymie} comme outils du vocabulaire : dire la même chose avec d'autres mots, ou dire le contraire. Ces outils trouvent maintenant leur emploi le plus noble et le plus rentable à l'examen : la \cle{reformulation}. Car résumer un texte, c'est précisément redire l'essentiel de la pensée d'un auteur \emph{avec ses propres mots}. Nous allons construire ensemble, étape par étape, la méthode complète du résumé, telle qu'elle est exigée au Probatoire et au Baccalauréat technique.

\courssub{Qu'est-ce que reformuler les idées essentielles ?}

\begin{definition}[Résumé et reformulation]
Le \cle{résumé} (ou \cle{contraction de texte}) est un texte bref, rédigé avec les mots du candidat, qui restitue fidèlement les \cle{idées essentielles} d'un texte plus long, sans rien ajouter ni rien déformer. \cle{Reformuler}, c'est exprimer une idée avec d'autres mots que ceux de l'auteur, en conservant exactement le même sens.
\end{definition}

L'intuition est simple : imagine qu'un camarade était absent lors d'une réunion de classe. Tu ne peux pas lui répéter mot pour mot tout ce qui a été dit pendant deux heures ; tu lui dis \emph{l'essentiel} en quelques phrases, avec tes propres mots, sans inventer. C'est exactement le travail du résumé.

\begin{propriete}[Règle d'or du résumé]
Un résumé ne contient \textbf{ni exemple, ni citation, ni opinion personnelle, ni phrase copiée du texte}. Toute phrase recopiée telle quelle du texte est sanctionnée, car elle prouve que le candidat n'a pas reformulé. Le résumé doit être rédigé à la troisième personne et respecter le nombre de mots imposé.
\end{propriete}

\courssub{Les six étapes du résumé}

\begin{methode}[Les 6 étapes du résumé]
\begin{enumerate}
\item \textbf{Lire} le texte deux fois : une première lecture globale pour saisir le thème, une seconde lecture attentive pour dégager la thèse de l'auteur.
\item \textbf{Découper} le texte en parties et souligner les idées essentielles.
\item \textbf{Éliminer} les idées secondaires : exemples, citations, répétitions, comparaisons, anecdotes, chiffres illustratifs.
\item \textbf{Reformuler} chaque idée essentielle avec ses propres mots.
\item \textbf{Rédiger} le résumé : texte suivi, cohérent, à la troisième personne, avec des connecteurs logiques.
\item \textbf{Vérifier} : compter les mots, contrôler la fidélité, corriger la grammaire et l'orthographe.
\end{enumerate}
\end{methode}

\begin{popfigure}
\begin{center}
\begin{tikzpicture}[
node distance=0.55cm,
etape/.style={rectangle, rounded corners, draw=popblue, fill=popblueL, text width=8.6cm, align=center, font=\small, inner sep=5pt},
fleche/.style={-{Stealth[length=3mm]}, thick, poporange}]
\node[etape] (e1) {\textbf{1. Lecture} : lire le texte deux fois ; dégager le thème et la thèse};
\node[etape, below=of e1] (e2) {\textbf{2. Découpage} : repérer les parties ; souligner une idée essentielle par paragraphe};
\node[etape, below=of e2] (e3) {\textbf{3. Élimination} : rayer exemples, citations, répétitions, anecdotes, chiffres};
\node[etape, below=of e3] (e4) {\textbf{4. Reformulation} : traduire chaque idée avec ses propres mots (synonymes, conversion, condensation)};
\node[etape, below=of e4] (e5) {\textbf{5. Rédaction} : texte suivi, 3\textsuperscript{e} personne, connecteurs logiques, aucun jugement};
\node[etape, below=of e5] (e6) {\textbf{6. Vérification} : compter les mots, vérifier la fidélité, corriger la langue};
\draw[fleche] (e1) -- (e2);
\draw[fleche] (e2) -- (e3);
\draw[fleche] (e3) -- (e4);
\draw[fleche] (e4) -- (e5);
\draw[fleche] (e5) -- (e6);
\end{tikzpicture}
\end{center}
\legende{Organigramme des six étapes de la contraction de texte : chaque étape prépare la suivante ; aucune ne peut être sautée.}
\end{popfigure}

\courssub{Étape 1 : la lecture globale (deux lectures)}

On ne résume pas un texte qu'on a lu une seule fois, en diagonale. La méthode impose \textbf{deux lectures} :

\begin{itemize}
\item \textbf{Première lecture (globale)} : lire le texte entièrement, sans stylo, pour répondre à deux questions : \emph{De quoi parle le texte ?} (c'est le \cle{thème}) et \emph{Que veut démontrer l'auteur ?} (c'est la \cle{thèse}). Le titre, s'il existe, est un indice précieux.
\item \textbf{Seconde lecture (attentive)} : relire lentement, crayon à la main, en observant la progression : comment l'auteur passe-t-il d'une idée à l'autre ? Où sont les arguments ? Où est la conclusion ?
\end{itemize}

\begin{exemplebox}[Thème et thèse]
Dans un texte où l'auteur affirme que le téléphone portable distrait les élèves, donne des exemples de classes perturbées et propose de l'interdire en cours :
\begin{itemize}
\item le \textbf{thème} est : le téléphone portable à l'école ;
\item la \textbf{thèse} est : le téléphone portable nuit aux apprentissages et doit être réglementé.
\end{itemize}
Confondre thème et thèse est une erreur : le thème est un sujet, la thèse est une prise de position.
\end{exemplebox}

\courssub{Étape 2 : le découpage en parties et le soulignage des idées essentielles}

Un texte argumentatif est construit : il possède une \cle{structure} (introduction, développement, conclusion) que nous avons étudiée dans les sections précédentes. Le découpage consiste à repérer cette structure et à attribuer à chaque partie son idée maîtresse.

\begin{methode}[Dégager la structure]
\begin{enumerate}
\item Repérer les \textbf{marqueurs de structure} : \emph{d'abord, ensuite, enfin, en conclusion} ; les alinéas ; les paragraphes.
\item Pour chaque paragraphe (ou groupe de paragraphes qui développent la même idée), formuler \textbf{une seule idée essentielle} en une phrase courte.
\item Vérifier que les idées dégagées se succèdent logiquement : elles doivent raconter, à elles seules, le raisonnement de l'auteur.
\end{enumerate}
\end{methode}

\begin{aretenir}[Une idée par paragraphe]
En général, \textbf{un paragraphe = une idée essentielle}. Si deux paragraphes disent la même chose sous des formes différentes, ils ne comptent que pour une idée. Un texte de 300 à 400 mots contient habituellement \textbf{3 à 5 idées essentielles}.
\end{aretenir}

\courssub{Étape 3 : la suppression des idées secondaires}

C'est l'étape décisive : le résumé naît de ce qu'on \emph{enlève}. Les idées secondaires sont tout ce qui illustre, répète ou décore la pensée sans la faire avancer.

\begin{center}
\begin{tabularx}{\linewidth}{|l|X|X|}\hline
\rowcolor{popblueL} & \textbf{À SUPPRIMER (idées secondaires)} & \textbf{À GARDER (idées essentielles)} \\ \hline
Exemples & \emph{« Ainsi, à Douala, un élève sur deux... »} & --- \\ \hline
Citations & \emph{« Comme le dit un proverbe... »} & --- \\ \hline
Répétitions & la même idée exprimée deux fois & garder l'idée une seule fois \\ \hline
Comparaisons, anecdotes & récits, images, souvenirs personnels de l'auteur & --- \\ \hline
Chiffres illustratifs & statistiques données à titre d'exemple & --- \\ \hline
Thèse & --- & la position défendue par l'auteur \\ \hline
Arguments & --- & les raisons qui soutiennent la thèse \\ \hline
Conclusion & --- & la solution ou le bilan proposé \\ \hline
\end{tabularx}
\end{center}

\begin{attention}[Ce qui doit disparaître]
Les élèves hésitent souvent à supprimer les \textbf{exemples} et les \textbf{citations} parce qu'ils sont frappants et faciles à retenir. C'est un piège : l'exemple ne fait qu'\emph{illustrer} l'idée ; c'est l'idée qu'il faut garder, jamais l'exemple. De même, une statistique spectaculaire (\emph{« 70 \% des jeunes... »}) n'est qu'un appui : elle disparaît du résumé.
\end{attention}

\courssub{Étape 4 : la reformulation de chaque idée essentielle}

Reformuler, c'est traduire chaque idée essentielle dans sa propre langue. Trois techniques, qui s'appuient sur la synonymie étudiée dans la section précédente, sont à maîtriser :

\begin{enumerate}
\item \textbf{La substitution par synonymes} : remplacer les mots de l'auteur par des équivalents. \emph{« Les jeunes délaissent la lecture »} devient \emph{« Les adolescents abandonnent les livres »}.
\item \textbf{La conversion (changement grammatical)} : transformer un nom en verbe, une phrase active en phrase passive, etc. \emph{« La lecture est en recul »} devient \emph{« on lit de moins en moins »}.
\item \textbf{La condensation} : exprimer en quelques mots ce que l'auteur dit en une phrase longue. \emph{« Les élèves qui passent leurs soirées devant les écrans ne trouvent plus le temps ni le goût d'ouvrir un livre »} se condense en \emph{« les écrans détournent les élèves de la lecture »}.
\end{enumerate}

\begin{exoresolu}[Reformuler une phrase du texte]
\textbf{Énoncé.} Reformulez, sans en changer le sens, la phrase suivante extraite d'un texte à résumer : \emph{« L'école, qui devrait être le sanctuaire du savoir, devient trop souvent, hélas, le théâtre de toutes les distractions. »}
\tcblower
\textbf{Solution détaillée.}
\begin{enumerate}
\item \textbf{Idée essentielle} de la phrase : l'école perd sa fonction d'apprentissage à cause des distractions.
\item \textbf{Éléments secondaires à éliminer} : l'image \emph{« sanctuaire du savoir »}, l'image \emph{« théâtre »}, l'exclamation \emph{« hélas »} (sentiment de l'auteur exprimé de façon oratoire).
\item \textbf{Reformulation par synonymes et condensation} : \emph{« L'école ne remplit plus toujours sa mission d'enseignement, car les distractions y prennent trop de place. »}
\end{enumerate}
Le sens est identique, les mots sont différents, les images ont disparu : la reformulation est réussie.
\end{exoresolu}

\courssub{Étape 5 : la rédaction du résumé}

Les idées reformulées ne doivent pas être juxtaposées comme une liste : elles doivent former un \textbf{texte suivi et cohérent}.

\begin{propriete}[Contraintes de rédaction]
Le résumé est rédigé :
\begin{itemize}
\item en \textbf{phrases liées} par des connecteurs logiques (\emph{d'abord, ensuite, en effet, cependant, c'est pourquoi, en conclusion}) qui reproduisent la progression du texte ;
\item à la \textbf{troisième personne} : on écrit \emph{« l'auteur montre que... »}, jamais \emph{« je pense que... »} ;
\item \textbf{sans aucun jugement personnel} : le résumeur n'approuve ni ne critique ; il rapporte fidèlement ;
\item dans le \textbf{nombre de mots imposé}, avec une marge de tolérance de 10 \% en plus ou en moins.
\end{itemize}
\end{propriete}

\begin{exemplebox}[Un résumé chiffré]
Un texte de \textbf{320 mots} doit être résumé en \textbf{80 mots} (le quart, proportion habituelle à l'examen). Le texte contient environ \textbf{4 idées essentielles} : le candidat dispose donc d'environ \textbf{20 mots par idée}. Ce calcul simple permet de doser la longueur de chaque phrase avant même de rédiger.
\end{exemplebox}

\begin{attention}[Le comptage des mots est obligatoire]
Erreur fréquente et lourdement sanctionnée : dépasser ou sous-estimer fortement le nombre de mots imposé. Un résumé de 150 mots pour 80 demandés prouve que le candidat n'a pas éliminé le secondaire ; un résumé de 30 mots prouve qu'il a perdu des idées essentielles. \textbf{Compter les mots et l'indiquer à la fin} (par exemple : \emph{« 78 mots »}) fait partie de l'exercice. Le non-respect de la consigne coûte des points.
\end{attention}

\courssub{Étape 6 : la vérification finale}

Avant de recopier au propre, le candidat relit son résumé en se posant trois questions :

\begin{enumerate}
\item \textbf{Le nombre de mots est-il respecté ?} Compter mot par mot, ajuster en condensant ou en développant légèrement une idée.
\item \textbf{Le résumé est-il fidèle ?} Comparer avec la liste des idées essentielles dégagées à l'étape 2 : aucune idée ne doit manquer, aucune idée étrangère au texte ne doit apparaître.
\item \textbf{La langue est-elle correcte ?} Vérifier les accords, la conjugaison, l'orthographe et la ponctuation : une expression fautive dévalorise un résumé pourtant fidèle.
\end{enumerate}

\courssub{Les qualités d'un bon résumé}

\begin{aretenir}[Les cinq qualités du résumé]
Un bon résumé se reconnaît à cinq qualités :
\begin{itemize}
\item la \textbf{fidélité} : il restitue exactement la pensée de l'auteur, sans la déformer ;
\item l'\textbf{objectivité} : il ne contient ni opinion, ni émotion, ni appréciation personnelle ;
\item la \textbf{brièveté} : il respecte le nombre de mots imposé ;
\item la \textbf{cohérence} : les idées s'enchaînent logiquement grâce aux connecteurs ;
\item la \textbf{clarté} : il se comprend sans avoir lu le texte original.
\end{itemize}
\end{aretenir}

\begin{exoresolu}[Résumer un court paragraphe]
\textbf{Énoncé.} Résumez en une phrase d'environ 20 mots le passage suivant (environ 60 mots) : \emph{« Le paludisme demeure un fléau dans nos régions. Chaque année, des milliers d'enfants, comme le petit Moussa, âgé de cinq ans, dont le cas a ému tout son quartier, sont emportés par cette maladie. Selon les statistiques récentes, près de 40 \% des consultations hospitalières lui sont dues. Pourtant, des moyens simples existent : les moustiquaires imprégnées, l'assainissement des quartiers et le traitement rapide des fièvres. »}
\tcblower
\textbf{Solution détaillée.}
\begin{enumerate}
\item \textbf{Idées essentielles} : (a) le paludisme reste grave ; (b) des moyens simples de prévention et de traitement existent.
\item \textbf{À supprimer} : l'anecdote du petit Moussa (exemple), le chiffre de 40 \% (statistique illustrative), le mot \emph{« fléau »} (image).
\item \textbf{Résumé proposé} (20 mots) : \emph{« Le paludisme demeure meurtrier, mais des moyens simples de prévention et de traitement permettent de le combattre efficacement. »}
\end{enumerate}
\end{exoresolu}

\courssub{Le résumé dans la vie quotidienne camerounaise}

La contraction de texte n'est pas un exercice scolaire artificiel : c'est une compétence de vie.

\begin{itemize}
\item \textbf{Compte rendu d'une réunion de tontine} : la secrétaire ne note pas chaque parole ; elle consigne les décisions essentielles (montant des cotisations, bénéficiaire du tour, date de la prochaine séance) en quelques lignes fidèles et objectives.
\item \textbf{Synthèse d'un discours du chef de village} : le notable chargé d'informer les absents retient la thèse du chef (par exemple l'organisation des travaux communautaires), les arguments et les décisions, et laisse de côté les anecdotes et les salutations.
\item \textbf{Résumé d'une émission de radio} pour un camarade absent : on rapporte le sujet, les positions des invités et la conclusion, sans y mêler son propre avis.
\end{itemize}

\begin{savaistu}[Le compte rendu, un métier]
Dans la vie professionnelle technique — secrétariat, gestion, administration, comptabilité — le \textbf{compte rendu de réunion} est une application directe de la contraction de texte. Un bon secrétaire de direction est précisément quelqu'un qui sait écouter deux heures de débats et en restituer l'essentiel en une page claire, fidèle et objective. Maîtriser le résumé au lycée, c'est déjà préparer son insertion professionnelle.
\end{savaistu}

\begin{exercice}[Entraînement à la reformulation]
Reformulez chacune des phrases suivantes en une phrase plus courte, avec vos propres mots, sans en changer le sens :
\begin{enumerate}
\item \emph{« L'agriculture, qui nourrit depuis toujours nos villages, est aujourd'hui délaissée par une jeunesse qui rêve des lumières de la ville. »}
\item \emph{« Comme le rappelle un ancien adage, c'est goutte à goutte que l'eau creuse la pierre : la persévérance finit toujours par vaincre les obstacles les plus durs. »}
\item \emph{« Dans notre établissement, où l'on compte cette année 1 200 élèves, la bibliothèque ne reçoit en moyenne que 15 visiteurs par jour, ce qui montre bien le désintérêt croissant pour le livre. »}
\end{enumerate}
\end{exercice}

\begin{corrige}[Exercice : entraînement à la reformulation]
\begin{enumerate}
\item Idée essentielle : les jeunes abandonnent l'agriculture pour la ville. Reformulation possible : \emph{« La jeunesse délaisse l'agriculture villageoise, attirée par la vie urbaine. »} (L'image \emph{« lumières de la ville »} et la formule \emph{« depuis toujours »} sont supprimées.)
\item Idée essentielle : la persévérance triomphe des difficultés. Reformulation possible : \emph{« La persévérance permet de surmonter les plus grandes difficultés. »} (Le proverbe cité est une citation : il disparaît.)
\item Idée essentielle : les élèves lisent de moins en moins. Reformulation possible : \emph{« La faible fréquentation de la bibliothèque traduit le désintérêt des élèves pour la lecture. »} (Les chiffres 1 200 et 15 sont des illustrations : ils sont supprimés.)
\end{enumerate}
\end{corrige}

En résumé, la contraction de texte obéit à une méthode rigoureuse en six étapes — lire, découper, éliminer, reformuler, rédiger, vérifier — dont chacune prépare la suivante. Celui qui la maîtrise dispose d'un outil puissant, à l'examen comme dans la vie professionnelle. La section suivante, consacrée au compte rendu et à la remédiation, permettra de s'entraîner dans les conditions réelles de l'épreuve.

\coursec{Compte rendu, remédiation et entraînement à l'examen}

Dans la section précédente, nous avons construit pas à pas la méthode complète du résumé : repérer la thèse, dégager les idées essentielles, supprimer les exemples et les répétitions, puis reformuler avec ses propres mots en respectant le volume imposé. Mais une méthode n'est vraiment acquise que si l'élève sait l'utiliser seul, le jour de l'examen, sans le professeur à ses côtés. Cette dernière section a donc trois objectifs : \cle{vérifier} que chacun maîtrise la démarche (compte rendu), \cle{corriger} les difficultés qui subsistent (remédiation) et \cle{s'entraîner} dans les conditions réelles du Probatoire et du Baccalauréat technique.

\courssub{Le compte rendu collectif : vérifier les acquis}

\begin{definition}[Compte rendu et remédiation]
Le \cle{compte rendu} est la restitution, orale ou écrite, de ce qui a été appris : il permet à l'élève de formuler lui-même la méthode et au professeur de vérifier qu'elle est comprise. La \cle{remédiation} est un enseignement correctif qui reprend, \emph{autrement}, les notions non acquises après une évaluation : on ne répète pas le même cours, on change d'approche, d'exemple ou d'exercice pour débloquer la difficulté.
\end{definition}

\textbf{Restitution orale de la méthode en 6 étapes.} Le professeur interroge la classe : « Racontez, sans regarder votre cahier, les six étapes du résumé. » Chaque élève cité formule une étape avec ses propres mots, et la classe complète ou corrige. On obtient collectivement le tableau suivant, qui fait la synthèse de toute l'unité :

\begin{center}
\begin{tabularx}{\linewidth}{|c|l|X|}
\hline
\rowcolor{popblueL} \textbf{Étape} & \textbf{Action} & \textbf{Question de contrôle} \\
\hline
1 & Lire le texte deux fois & Ai-je compris le sens général ? \\
\hline
2 & Repérer la thèse & Quelle est l'idée principale défendue par l'auteur ? \\
\hline
3 & Dégager les idées essentielles & Quels sont les arguments qui soutiennent la thèse ? \\
\hline
4 & Supprimer le superflu & Ai-je éliminé exemples, citations, répétitions, détails ? \\
\hline
5 & Reformuler & Ai-je utilisé mes propres mots (synonymes, changements de tournures) ? \\
\hline
6 & Compter et rédiger & Ai-je respecté le nombre de mots imposé et relu ma copie ? \\
\hline
\end{tabularx}
\end{center}

\begin{attention}[Ce que le compte rendu doit révéler]
Si un élève confond l'étape 2 (la thèse) avec l'étape 3 (les arguments), ou s'il ne sait pas dire \emph{pourquoi} on supprime les exemples, c'est le signe que la notion n'est pas acquise : il faudra une remédiation ciblée. Le compte rendu n'est pas une simple récitation : c'est un \cle{diagnostic}.
\end{attention}

\courssub{Correction détaillée d'un résumé-type : identifier les erreurs}

L'observation la plus formatrice, à ce stade, consiste à corriger \emph{ensemble} un résumé rédigé par un élève (anonymisé) ou fabriqué par le professeur avec des erreurs volontaires. La classe joue le rôle du correcteur d'examen. Voici un cas d'école.

\begin{exoresolu}[Corriger un résumé fautif]
\textbf{Énoncé.} Texte de départ (environ 100 mots) : « Le travail manuel est méprisé dans nos sociétés modernes. Pourtant, il est indispensable. Qui répare nos routes, construit nos maisons, cultive nos champs ? Le maçon, le menuisier, l'agriculteur. Par exemple, à Bafoussam, de jeunes menuisiers fabriquent des meubles admirables que toute la ville s'arrache. Le travail manuel développe aussi l'intelligence pratique : celui qui fabrique un objet comprend le monde mieux que celui qui l'achète. Il faut donc réhabiliter le travail manuel, lui rendre sa dignité, sa noblesse, sa valeur. » Consigne : résumez ce texte en 30 mots (± 10 \%). Résumé d'un élève : « Je pense que le travail manuel est méprisé. Qui répare nos routes, construit nos maisons, cultive nos champs ? Par exemple, à Bafoussam, de jeunes menuisiers fabriquent des meubles admirables. Il faut réhabiliter le travail manuel. » (43 mots)
\tcblower
\textbf{Correction détaillée.} On relève quatre erreurs types :
\begin{enumerate}
\item \textbf{Avis personnel} : « Je pense que » est interdit dans un résumé, qui doit rester objectif.
\item \textbf{Phrase copiée} : la question « Qui répare nos routes... ? » est recopiée mot pour mot du texte.
\item \textbf{Exemple conservé} : l'exemple de Bafoussam est un détail illustratif : il devait être supprimé.
\item \textbf{Dépassement du volume} : 43 mots au lieu de 30 (± 3) : hors barème.
\end{enumerate}
\textbf{Résumé corrigé} (27 mots) : « Le travail manuel, bien que méprisé, reste indispensable à la vie sociale et développe l'intelligence pratique. Il faut donc urgemment lui rendre sa dignité et sa valeur. »
\end{exoresolu}

\begin{attention}[Erreur fréquente : l'avis personnel]
Ne jamais introduire son opinion dans un résumé : « je pense que », « à mon avis », « je trouve que l'auteur a raison » sont des fautes graves. Le résumé est un exercice de \cle{fidélité} : il restitue la pensée de l'auteur, totalement \cle{objectivement}, sans la juger ni la commenter.
\end{attention}

\courssub{La remédiation ciblée : exercices différenciés}

Tous les élèves n'ont pas la même difficulté. La remédiation efficace est donc \cle{différenciée} : à chaque faiblesse diagnostiquée correspond un exercice précis.

\begin{center}
\begin{tabularx}{\linewidth}{|l|X|X|}
\hline
\rowcolor{popblueL} \textbf{Difficulté observée} & \textbf{Cause fréquente} & \textbf{Exercice de remédiation} \\
\hline
Repérage de la thèse & L'élève confond le sujet et l'idée défendue & Souligner la thèse dans 3 textes courts ; choisir la bonne formulation parmi 3 propositions \\
\hline
Reformulation & Vocabulaire pauvre, peur de « trahir » le texte & Remplacer 10 mots soulignés par des synonymes ; transformer des phrases nominales en phrases verbales \\
\hline
Comptage des mots & Méthode de comptage inconnue & Compter les mots de 5 phrases données ; réduire une phrase de 20 mots à 12 mots sans perdre le sens \\
\hline
Suppression du superflu & L'élève croit que « tout est important » & Barrer, dans un paragraphe, les exemples, citations et répétitions, puis vérifier que le sens demeure \\
\hline
\end{tabularx}
\end{center}

\begin{exercice}[Remédiation : réduire sans trahir]
Réduisez chaque phrase à environ la moitié de sa longueur, sans perdre l'idée essentielle :
\begin{enumerate}
\item « L'éducation, qui est un droit fondamental reconnu à tous les enfants du monde, demeure malheureusement inaccessible à des millions d'enfants africains. »
\item « Le mensonge, c'est-à-dire le fait de dire volontairement ce qui est faux, détruit peu à peu la confiance qui unit les hommes. »
\end{enumerate}
\end{exercice}

\begin{corrige}[Exercice de remédiation]
\begin{enumerate}
\item « L'éducation, droit fondamental, reste inaccessible à des millions d'enfants africains. » (11 mots au lieu de 21 : la définition du droit est supprimée, l'idée conservée.)
\item « Le mensonge détruit la confiance entre les hommes. » (8 mots au lieu de 22 : la définition et la gradation « peu à peu » sont supprimées.)
\end{enumerate}
\end{corrige}

\courssub{Le barème d'évaluation du résumé au Baccalauréat technique}

Pour s'entraîner utilement, il faut connaître les critères du correcteur. Un résumé est noté sur cinq critères complémentaires, schématisés ci-dessous :

\begin{popfigure}
\begin{center}
\begin{tikzpicture}[x=1cm,y=1cm]
% Barres horizontales du barème
\foreach \y/\crit/\pts/\col in {4/{Fidélité au texte}/4/popblue, 3/{Reformulation}/4/popgreen, 2/{Cohérence et enchaînement}/4/poporange, 1/{Langue (orthographe, syntaxe)}/4/poppurple, 0/{Respect du volume}/4/poppink}{
  \fill[\col!70] (0,\y) rectangle (\pts*0.9,\y+0.6);
  \node[anchor=east, font=\small] at (-0.1,\y+0.3) {\crit};
  \node[font=\small\bfseries, text=popdark] at (\pts*0.9+0.4,\y+0.3) {\pts\ pts};
}
\draw[->, thick] (0,-0.4) -- (4.6,-0.4);
\node[font=\small, below] at (2.3,-0.5) {points attribués (total : 20)};
\end{tikzpicture}
\end{center}
\legende{Barème indicatif d'évaluation d'un résumé : cinq critères de 4 points chacun. La fidélité (le sens de l'auteur est respecté), la reformulation (aucune phrase copiée), la cohérence (les idées s'enchaînent logiquement), la langue et le respect du volume comptent à parts égales.}
\end{popfigure}

\begin{aretenir}[Les cinq critères du correcteur]
\begin{itemize}
\item \textbf{Fidélité (4 pts)} : la thèse et les arguments de l'auteur sont restitués sans déformation ni ajout.
\item \textbf{Reformulation (4 pts)} : les idées sont exprimées avec les mots de l'élève ; toute phrase copiée est sanctionnée.
\item \textbf{Cohérence (4 pts)} : les connecteurs logiques (\emph{d'abord, en effet, cependant, donc}) assurent l'enchaînement.
\item \textbf{Langue (4 pts)} : orthographe, grammaire, ponctuation.
\item \textbf{Respect du volume (4 pts)} : le nombre de mots imposé (± 10 \%) est respecté et indiqué à la fin.
\end{itemize}
\end{aretenir}

\begin{savaistu}[Le poids de la langue]
La note de langue (orthographe, syntaxe, ponctuation) peut représenter jusqu'à \emph{un quart} des points de l'épreuve de contraction de texte. Un résumé parfait sur le fond mais truffé de fautes perd donc une grande partie de sa valeur : la \cle{relecture} finale n'est pas un luxe, c'est une étape qui rapporte des points.
\end{savaistu}

\courssub{Entraînement en conditions d'examen}

\begin{exemplebox}[Sujet type Baccalauréat technique]
« Lisez attentivement le texte suivant (300 mots), puis résumez-le \cle{au quart de sa longueur}, soit 75 mots (± 10 \%, donc entre 68 et 83 mots). Vous indiquerez obligatoirement le nombre de mots à la fin de votre résumé. Durée : 45 minutes. »
\end{exemplebox}

\begin{methode}[Gérer ses 45 minutes le jour de l'examen]
\begin{enumerate}
\item \textbf{5 minutes} : lire le texte deux fois et souligner la thèse.
\item \textbf{10 minutes} : dégager au brouillon les 3 ou 4 idées essentielles ; barrer exemples et répétitions.
\item \textbf{15 minutes} : rédiger le résumé au brouillon en reformulant, avec des connecteurs logiques.
\item \textbf{5 minutes} : compter les mots et ajuster (supprimer ou condenser si l'on dépasse).
\item \textbf{10 minutes} : recopier au propre et relire l'orthographe ; inscrire le nombre de mots à la fin.
\end{enumerate}
\end{methode}

\begin{exoresolu}[Résumé en conditions d'examen]
\textbf{Énoncé.} Résumez en 50 mots (± 10 \%) le passage suivant : « La lecture est délaissée par les jeunes générations, qui préfèrent les écrans. Cette évolution est regrettable, car lire développe des capacités irremplaçables. D'abord, la lecture enrichit le vocabulaire : un lecteur régulier possède trois fois plus de mots qu'un non-lecteur. Ensuite, elle entraîne la concentration, qualité indispensable dans les études et dans la vie professionnelle. Enfin, elle ouvre l'esprit : par les livres, un enfant de Yaoundé peut voyager au Japon, vivre au Moyen Âge, comprendre le cœur humain. Prenons l'exemple de mon voisin, instituteur, qui lit un roman chaque semaine : sa culture est remarquable. Il faut donc, à l'école comme en famille, redonner aux jeunes le goût de lire. »
\tcblower
\textbf{Solution.} Thèse : il faut redonner aux jeunes le goût de lire. Arguments : la lecture enrichit le vocabulaire, entraîne la concentration, ouvre l'esprit. À supprimer : la statistique « trois fois plus », l'exemple du voisin instituteur, les voyages au Japon et au Moyen Âge (illustrations). Résumé (48 mots) : « Les jeunes délaissent aujourd'hui la lecture au profit des écrans. Pourtant, la lecture régulière enrichit considérablement le vocabulaire, développe la concentration indispensable et ouvre l'esprit sur le monde et sur autrui. L'école comme la famille doivent donc vraiment redonner aux jeunes le goût de la lecture. »
\end{exoresolu}

\courssub{L'auto-évaluation : la grille de relecture finale}

Le candidat autonome est celui qui sait se corriger lui-même. Avant de rendre la copie, il passe son résumé au crible d'une grille de relecture.

\begin{methode}[La grille de relecture finale en 6 points]
\begin{enumerate}
\item La \textbf{thèse} du texte est présente dans mon résumé.
\item Je n'ai gardé \textbf{que les idées essentielles} (ni exemples, ni citations, ni répétitions).
\item \textbf{Aucune phrase n'est copiée} du texte : j'ai reformulé avec mes mots.
\item Des \textbf{connecteurs logiques} relient mes idées.
\item Le \textbf{nombre de mots} est respecté et indiqué à la fin.
\item J'ai \textbf{vérifié l'orthographe} et la ponctuation.
\end{enumerate}
\end{methode}

\begin{popfigure}
\begin{center}
\begin{tikzpicture}[font=\small]
\foreach \y/\txt in {5.4/{La thèse du texte est présente.}, 4.5/{Seules les idées essentielles sont gardées.}, 3.6/{Aucune phrase copiée : tout est reformulé.}, 2.7/{Des connecteurs logiques relient les idées.}, 1.8/{Le nombre de mots est respecté et indiqué.}, 0.9/{Orthographe et ponctuation vérifiées.}}{
  \draw[thick, popblue] (0,\y) rectangle (0.45,\y+0.45);
  \node[anchor=west, text=popdark] at (0.7,\y+0.22) {\txt};
}
\node[font=\small\itshape, text=popmuted] at (3.4,0.2) {Cochez chaque case avant de rendre la copie.};
\end{tikzpicture}
\end{center}
\legende{Grille d'auto-évaluation du résumé : six cases à cocher. Si une case reste vide, le résumé n'est pas prêt à être rendu.}
\end{popfigure}

\begin{exercice}[Entraînement à la maison]
Choisissez un éditorial de journal d'environ 300 mots. Résumez-le en 75 mots en 45 minutes chronométrées, puis évaluez-vous avec le barème des cinq critères (sur 20). Recommencez une semaine plus tard avec un autre texte et comparez vos notes.
\end{exercice}

\courssub{Pour aller plus loin : du résumé au commentaire (hors programme strict)}

\begin{savaistu}[Complément utile au Bac : le commentaire composé]
\emph{Notion signalée à titre indicatif, hors programme strict de Première technique.} Le résumé exige une objectivité totale : restituer la pensée d'autrui sans la juger. Le \cle{commentaire composé}, lui, prolonge ce travail par une \emph{analyse personnelle} : il s'agit d'expliquer comment le texte est construit, quels procédés l'auteur utilise, et d'en apprécier la force persuasive. Celui qui maîtrise le résumé possède déjà le premier geste du commentaire : comprendre fidèlement un texte avant de le discuter.
\end{savaistu}

\begin{aretenir}[L'essentiel de la section]
Le compte rendu vérifie que la méthode en six étapes est acquise ; la remédiation corrige, par des exercices différenciés, les difficultés restantes (thèse, reformulation, comptage). Le jour de l'examen, le candidat gère ses 45 minutes, respecte les cinq critères du barème (fidélité, reformulation, cohérence, langue, volume) et relit sa copie à l'aide de la grille en six points. Le résumé est un exercice d'objectivité et de rigueur : c'est l'entraînement régulier, en conditions réelles, qui transforme la méthode en réflexe.
\end{aretenir}

\newpage

\coursec{Activité d'intégration}

\coursec{Leçon 19 : Formuler les idées essentielles d'un texte à résumer}

\courssub{Objectifs de la leçon}

À l'issue de cette leçon, l'élève doit être capable de :

\begin{itemize}
\item analyser la \cle{structure argumentative} d'un texte (thèse, arguments, exemples, réfutation, conclusion) ;
\item identifier et justifier la \cle{tonalité} dominante (comique, dramatique, satirique, tragique) par des indices linguistiques ;
\item distinguer les \cle{idées essentielles} des idées secondaires (exemples, citations, chiffres, répétitions, développements) ;
\item \cle{reformuler} les idées essentielles en mobilisant la synonymie, l'antonymie et la conversion grammaticale ;
\item rédiger un \cle{résumé} conforme aux exigences du Baccalauréat technique : fidélité, condensation (environ le quart du texte), cohérence, correction, respect de la contrainte de longueur ;
\item auto-évaluer sa production à l'aide d'une grille critériée et transférer cette compétence à l'étude de l'œuvre au programme (\emph{Le Lion et la Perle}).
\end{itemize}

\courssub{Le texte argumentatif : caractéristiques et structure argumentative}

\begin{definition}[Texte argumentatif]
Un texte argumentatif a pour but de \textbf{convaincre} ou de \textbf{persuader} un destinataire. Il défend une \cle{thèse} (point de vue) à l'aide d'\cle{arguments} (raisons logiques) illustrés par des \cle{exemples} (faits, chiffres, citations, anecdotes). Il anticipe les \cle{objections} (contre-arguments) qu'il \cle{réfute}, pour aboutir à une \cle{conclusion} qui réaffirme la thèse ou ouvre une perspective.
\end{definition}

\begin{propriete}[Structure canonique d'un éditorial / article d'opinion]
\begin{enumerate}
\item \textbf{Accroche / Constat :} Présentation du sujet, souvent par un paradoxe ou une situation choquante.
\item \textbf{Thèse :} Affirmation claire de la position défendue (souvent au § 2).
\item \textbf{Arguments + Exemples :} Preuves factuelles, statistiques, témoignages, comparaisons (souvent § 3).
\item \textbf{Réfutation (concession + réfutation) :} Prise en compte des objections adverses pour mieux les démonter (§ 4).
\item \textbf{Conclusion :} Synthèse, appel à l'action, ouverture solennelle (§ 5).
\end{enumerate}
\end{propriete}

\begin{methode}[Repérer la structure argumentative]
\begin{enumerate}
\item Numéroter les paragraphes.
\item Pour chaque paragraphe, identifier la \textbf{fonction} (constat, thèse, argument, exemple, objection, réfutation, conclusion).
\item Formuler l'\textbf{idée essentielle} en une phrase simple (sujet + verbe + complément), sans recopier le texte.
\item Éliminer mentalement les développements illustratifs (noms propres, chiffres, citations, détails).
\end{enumerate}
\end{methode}

\courssub{Les tonalités : comique, dramatique, satirique, tragique}

\begin{aretenir}[Les quatre tonalités majeures]
\begin{center}
\begin{tabularx}{\linewidth}{|l|X|X|}\hline
\rowcolor{popblueL} \textbf{Tonalité} & \textbf{Effet visé} & \textbf{Indices linguistiques repérables} \\ \hline
\textbf{Comique} & Faire rire, divertir & Jeux de mots, quiproquos, situations burlesques, hyperboles légères, lexique du ridicule. \\ \hline
\textbf{Satirique} & Dénoncer, corriger par la moquerie & \cle{Ironie} (dire le contraire de ce qu'on pense), \cle{antithèse} violente, \cle{caricature}, humour noir, verbes dépréciatifs, guillemets distanciateurs. \\ \hline
\textbf{Dramatique} & Émouvoir, susciter la pitié & Lexique du malheur/douleur, pathos, exclamations, questions rhétoriques angoissées, temps courts, focalisation sur la victime. \\ \hline
\textbf{Tragique} & Montrer la fatalité, l'inéluctable & Lexique de la fatalité/destin, impuissance du héros, ton solennel, généralisations, passé simple/imparfait, absence d'issue. \\ \hline
\end{tabularx}
\end{center}
\end{aretenir}

\begin{attention}[Piège fréquent]
Ne pas confondre \textbf{ton} (attitude de l'énonciateur : ironique, lyrique, pathétique) et \textbf{tonalité} (catégorie générique : comique, satirique...). Un texte satirique \emph{utilise} un ton ironique. Justifiez toujours par des \textbf{citations précises} (mots, tournures, figures de style), jamais par une impression.
\end{attention}

\courssub{Synonymie et antonymie : outils de la reformulation}

\begin{propriete}[Procédés de reformulation (paraphrase)]
Pour dire autrement sans trahir le sens :
\begin{itemize}
\item \textbf{Synonymie :} Remplacer un mot par un équivalent (\emph{urgence} $\rightarrow$ \emph{nécessité, priorité, impératif} ; \emph{obligatoire} $\rightarrow$ \emph{imposé, requis, indispensable}).
\item \textbf{Antonymie + Négation :} Utiliser le contraire avec une négation (\emph{pas une relégation} $\rightarrow$ \emph{loin d'être un échec} ; \emph{coûte cher} $\rightarrow$ \emph{n'est pas gratuit}).
\item \textbf{Conversion grammaticale (transcatégorisation) :} Changer la nature grammaticale.
  \begin{itemize}
  \item Nom $\rightarrow$ Verbe : \emph{la solution existe} $\rightarrow$ \emph{il est possible de \textbf{résoudre} le problème}.
  \item Verbe $\rightarrow$ Nom : \emph{nos ateliers manquent de techniciens} $\rightarrow$ \emph{la \textbf{pénurie} de techniciens}.
  \item Actif $\rightarrow$ Passif (ou inversement) : \emph{rendre l'enseignement obligatoire} $\rightarrow$ \emph{l'enseignement \textbf{est imposé}}.
  \item Phrase $\rightarrow$ Subordonnée / Groupe nominal : \emph{le pays forme des chômeurs} $\rightarrow$ \emph{la \textbf{formation de chômeurs} par le pays}.
  \end{itemize}
\item \textbf{Changement de point de vue / voix :} \emph{L'auteur propose} $\rightarrow$ \emph{Il est proposé par l'auteur}.
\end{itemize}
\end{propriete}

\courssub{La reformulation des idées essentielles : méthode complète du résumé}

\begin{methode}[La méthode du résumé en 6 étapes (exigible au Bac Tech)]
\begin{enumerate}
\item \textbf{Lecture active :} Lire 2 fois (globale puis détaillée). Souligner les connecteurs, les mots-clés, la thèse.
\item \textbf{Analyse structurelle :} Tableau paragraphe / fonction / idée essentielle (étape 1 de l'activité).
\item \textbf{Sélection / Élagage :} Supprimer exemples, chiffres, citations, noms propres, répétitions, adjectifs qualificatifs non essentiels, commentaires de l'auteur (\emph{« il est temps »}).
\item \textbf{Reformulation :} Réécrire chaque idée essentielle en appliquant synonymie, antonymie, conversion (étape 4 de l'activité). \textbf{Ne rien recopier.}
\item \textbf{Rédaction :} Enchaîner les phrases reformulées avec des \cle{connecteurs logiques} variés (au moins 3 : \emph{d'abord, en effet, cependant, donc, c'est pourquoi, enfin, par conséquent...}). Respecter l'ordre des idées du texte.
\item \textbf{Contrôle / Comptage :} Compter les mots (mots-outils comptés pour 1). Vérifier : fidélité, cohérence, correction, contrainte de longueur ($\pm 10\%$).
\end{enumerate}
\end{methode}

\courssub{Lecture méthodique : Le Lion et la Perle (n°4 et n°5)}

\begin{experience}[Transfert de compétence : Résumer une scène de \emph{Le Lion et la Perle}]
L'œuvre au programme (\emph{Le Lion et la Perle} de Wole Soyinka) offre un terrain d'application idéal.
\begin{itemize}
\item \textbf{Lecture n°4 (Acte 2, Scène 1 - La danse du voyageur perdu) :} Dégager la structure argumentative implicite de la mime : \textbf{Constat} (la voiture en panne) $\rightarrow$ \textbf{Thèse satirique} (l'administration coloniale est inefficace) $\rightarrow$ \textbf{Arguments par l'image} (le fonctionnaire qui ne fait rien, le chauffeur qui répare) $\rightarrow$ \textbf{Conclusion} (le rire du village triomphe de l'autorité).
\item \textbf{Lecture n°5 (Acte 3 - Le retour de Lakunle / Le triomphe de Baroka) :} Résumer en 60 mots ($\pm 10\%$) le monologue final de Baroka ou la scène de la fête finale. Appliquer la méthode en 6 étapes : élaguer les répliques secondaires, reformuler la thèse de Baroka (\emph{la tradition absorbe la modernité}), respecter la tonalité (satirique / comique).
\end{itemize}
\textbf{Consigne d'entraînement :} En binôme, produire ce résumé de 60 mots, l'afficher, l'auto-évaluer avec la grille ci-dessous.
\end{experience}

\courssub{Situation d'intégration : « Lu dans la presse »}

\courssub{Situation}

Le journal du lycée technique de Bafoussam prépare son numéro spécial « Orientation ». La rédaction, dirigée par les élèves de Première technique, souhaite publier une rubrique « Lu dans la presse » : il s'agit de proposer aux lecteurs pressés le résumé d'un éditorial récent. C'est ta classe qui a été choisie pour produire ce résumé. Le rédacteur en chef impose une contrainte stricte : le résumé devra comporter exactement \textbf{100 mots ($\pm$ 10 \%)}, soit entre 90 et 110 mots, pour tenir dans l'encadré prévu en page 2.

Voici l'éditorial à résumer (400 mots), paru dans la rubrique « Opinion » d'un quotidien camerounais :

\begin{exemplebox}[Éditorial à résumer : « L'enseignement technique, une urgence nationale »]
\textbf{§ 1.} Depuis des années, nos élites répètent que l'avenir du Cameroun se joue dans les livres. Soit. Mais dans quels livres ? Pendant que des milliers de bacheliers généraux encombrent les amphithéâtres sans débouchés précis, nos ateliers manquent de techniciens, nos chantiers manquent de conducteurs de travaux, nos usines manquent de main-d'œuvre qualifiée. Le paradoxe est cinglant : un pays qui forme des chômeurs diplômés tout en important des compétences.

\textbf{§ 2.} La solution existe pourtant : rendre l'enseignement technique obligatoire dès le collège. Dès la sixième, chaque élève devrait toucher un outil, démonter un moteur, câbler un circuit, tenir une comptabilité. Non pas pour enfermer les enfants dans un métier, mais pour leur donner le goût du concret et le respect du travail manuel, si injustement méprisé chez nous.

\textbf{§ 3.} Les faits plaident en ce sens. À Douala, le lycée technique de Deïdo place chaque année plus de 80 \% de ses bacheliers dans l'emploi ou la formation supérieure, quand certains lycées généraux voient leurs lauréats grossir les rangs des mototaxis. En Allemagne, pays cité en exemple par tous nos ministres, la filière professionnelle accueille près de la moitié des adolescents et constitue le pilier de la puissance industrielle. Un chef d'entreprise de Bonabéri confiait récemment : « Je préfère embaucher un jeune technicien camerounais bien formé que d'importer un expert à prix d'or. »

\textbf{§ 4.} Certains objectent que l'enseignement technique coûte cher : ateliers, machines, matières consommables. L'argument est recevable, mais il est à courte vue. Un technicien employé coûte moins cher à la nation qu'un diplômé désœuvré. D'autres craignent une sélection précoce et injuste. C'est oublier que l'orientation technique n'est pas une relégation : les passerelles vers l'enseignement général et supérieur existent et doivent être renforcées.

\textbf{§ 5.} Il est donc temps, il est même grand temps, que le législateur inscrive l'enseignement technique obligatoire au collège dans nos textes. Nos enfants méritent mieux que des diplômes décoratifs : ils méritent des compétences. La nation, elle, y gagnera son industrialisation tant annoncée.
\end{exemplebox}

\courssub{Consignes}

Par groupes de quatre élèves, et en suivant l'ordre des questions :

\begin{enumerate}
\item Numérote les cinq paragraphes de l'éditorial, puis remplis le tableau de la structure argumentative : pour chaque paragraphe, indique sa fonction (constat, thèse, argument, exemple, réfutation d'objection, conclusion) et l'idée essentielle en une phrase.
\item Quelle est la \cle{tonalité} dominante de cet éditorial ? Justifie ta réponse en citant \textbf{trois indices linguistiques} précis (mots, tournures, procédés) relevés dans le texte.
\item Relève dans les § 3 et § 4 les éléments \textbf{secondaires} à éliminer dans un résumé : exemples, chiffres, citations, noms propres. Justifie chaque élimination.
\item Reformule chacune des quatre idées essentielles dégagées (Constat, Thèse, Argument principal, Réfutation/Conclusion) en utilisant, pour chacune, \textbf{au moins deux synonymes} et \textbf{une conversion grammaticale} (nom $\leftrightarrow$ verbe, actif $\leftrightarrow$ passif, phrase $\leftrightarrow$ GN).
\item Rédige collectivement le résumé en \textbf{100 mots ($\pm$ 10 \%)}, en employant au moins \textbf{trois connecteurs logiques} différents (\emph{d'abord, en effet, cependant, donc, c'est pourquoi, enfin...}). Compte les mots et note le total obtenu.
\item Auto-évaluez votre production à l'aide de la grille de relecture ci-dessous, puis affichez votre résumé au tableau. La classe comparera les versions et élira le résumé le plus fidèle et le plus condensé, en justifiant le vote.
\end{enumerate}

\begin{center}
\begin{tabularx}{\linewidth}{|l|X|c|}\hline
\rowcolor{popblueL} \textbf{Critère} & \textbf{Question de relecture} & \textbf{Oui/Non} \\ \hline
Fidélité & Le résumé respecte-t-il la thèse de l'auteur sans ajouter d'opinion personnelle ? & \\ \hline
Essentiel & Toutes les idées essentielles sont-elles présentes ? Les exemples et chiffres ont-ils disparu ? & \\ \hline
Reformulation & Aucune phrase du texte n'est-elle recopiée telle quelle ? & \\ \hline
Cohérence & Les connecteurs logiques rendent-ils l'enchaînement clair ? & \\ \hline
Contrainte & Le nombre de mots est-il compris entre 90 et 110 ? & \\ \hline
Correction & Orthographe, grammaire et ponctuation sont-elles soignées ? & \\ \hline
\end{tabularx}
\end{center}

\courssub{Ressources à mobiliser}

\begin{aretenir}[Boîte à outils du résumé]
\begin{itemize}
\item \textbf{Structure argumentative} : thèse (point de vue défendu) $\rightarrow$ arguments (raisons) $\rightarrow$ exemples (illustrations, à éliminer) $\rightarrow$ réfutation des objections $\rightarrow$ conclusion.
\item \textbf{Tonalités} : comique (faire rire), satirique (dénoncer par la moquerie), dramatique (émouvoir par le malheur), tragique (fatalité irrémédiable). On les repère par le lexique, les exclamations, l'ironie, l'hyperbole.
\item \textbf{Reformulation} : synonymie (\emph{urgence} = \emph{nécessité, priorité}), antonymie avec négation (\emph{pas une relégation} = \emph{loin d'être un échec}), conversion grammaticale (\emph{la solution existe} $\rightarrow$ \emph{il est possible de résoudre}).
\item \textbf{Règles du résumé} : ne rien recopier, ne rien ajouter, supprimer exemples/chiffres/citations, respecter le quart du texte environ (ici : 100 mots pour 400), rédiger en phrases liées par des connecteurs.
\end{itemize}
\end{aretenir}

\courssub{Démarche et corrigé commenté}

\begin{exoresolu}[Résolution complète de l'activité]
Énoncé : résumer l'éditorial « L'enseignement technique, une urgence nationale » (400 mots) en 100 mots ($\pm$ 10 \%), après analyse de la structure, de la tonalité et reformulation des idées essentielles.
\tcblower
\textbf{Étape 1 : tableau de la structure argumentative.}

\begin{center}
\begin{tabularx}{\linewidth}{|c|l|X|}\hline
\rowcolor{popblueL} \textbf{§} & \textbf{Fonction} & \textbf{Idée essentielle} \\ \hline
1 & Constat / Accroche & Le Cameroun forme des diplômés au chômage alors qu'il manque cruellement de techniciens. \\ \hline
2 & Thèse & Il faut rendre l'enseignement technique obligatoire dès le collège pour valoriser le concret. \\ \hline
3 & Arguments + Exemples & La filière technique assure de réels débouchés (preuves camerounaises et étrangères). \\ \hline
4 & Réfutation (Concession + Réfutation) & Le coût et le risque de sélection précoce ne sont pas des obstacles insurmontables. \\ \hline
5 & Conclusion / Appel & Le législateur doit instituer cet enseignement pour l'avenir des jeunes et de la nation. \\ \hline
\end{tabularx}
\end{center}

\textbf{Étape 2 : tonalité dominante.} La tonalité est \cle{satirique} (avec une pointe dramatique sur le sort des jeunes). Indices linguistiques précis :
\begin{enumerate}
\item \textbf{Ironie} (§ 1) : « un pays qui forme des chômeurs diplômés » (antithèse choc diplôme/chômage).
\item \textbf{Antithèse méprisante} (§ 5) : « diplômes décoratifs » vs « compétences » (dévalorisation du général, valorisation du technique).
\item \textbf{Hyperbole critique} (§ 1) : « encombrent les amphithéâtres » (verbe péjoratif, image de l'inutilité).
\end{enumerate}
L'auteur dénonce une aberration nationale en ridiculisant le mépris social pour le manuel.

\textbf{Étape 3 : élimination des idées secondaires (§ 3 et § 4).}
On supprime impérativement :
\begin{itemize}
\item \textbf{Noms propres / Lieux :} « lycée technique de Deïdo », « Douala », « Allemagne », « Bonabéri » (ancrage contextuel, non idée).
\item \textbf{Chiffres / Statistiques :} « plus de 80 \% », « près de la moitié des adolescents » (illustration quantitative).
\item \textbf{Citation / Témoignage :} « Je préfère embaucher... » (exemple vivant, parole rapportée).
\item \textbf{Détails concrets :} « toucher un outil, démonter un moteur, câbler un circuit, tenir une comptabilité » (§ 2, liste illustrative) ; « ateliers, machines, matières consommables » (§ 4, liste illustrative).
\end{itemize}
On garde l'idée générale : \emph{La filière technique conduit efficacement à l'emploi, comme le montrent l'expérience nationale et internationale.}

\textbf{Étape 4 : reformulation des idées essentielles (exemples de productions élèves).}
\begin{itemize}
\item \textbf{Idée 1 (Constat) :} « nos ateliers manquent de techniciens » $\rightarrow$ \textbf{« le pays souffre d'une pénurie de techniciens »} (Synonymes : \emph{manquer} = \emph{souffrir d'une pénurie} / \emph{être en déficit} ; Conversion : Verbe \emph{manquer} $\rightarrow$ Nom \emph{pénurie}).
\item \textbf{Idée 2 (Thèse) :} « rendre l'enseignement technique obligatoire dès le collège » $\rightarrow$ \textbf{« l'enseignement technique devrait être imposé dès le collège »} (Synonymes : \emph{obligatoire} = \emph{imposé} / \emph{requis} ; Conversion : Actif \emph{rendre obligatoire} $\rightarrow$ Passif \emph{être imposé}).
\item \textbf{Idée 3 (Argument) :} « la solution existe » $\rightarrow$ \textbf{« il est possible de remédier à ce problème »} (Conversion : Nom \emph{solution} $\rightarrow$ Verbe \emph{remédier} ; Synonyme : \emph{existe} = \emph{est possible}).
\item \textbf{Idée 4 (Réfutation/Conclusion) :} « l'orientation technique n'est pas une relégation » $\rightarrow$ \textbf{« l'orientation technique est loin d'être un échec »} (Antonymie + Négation : \emph{relégation} $\rightarrow$ \emph{échec}).
\end{itemize}

\textbf{Étape 5 : résumé rédigé (100 mots exacts).}

\emph{Le Cameroun forme de nombreux diplômés sans emploi alors que le pays souffre d'une pénurie de techniciens. C'est pourquoi l'éditorialiste propose que l'enseignement technique soit imposé dès le collège, afin de donner aux élèves le goût du concret. En effet, cette filière offre de véritables débouchés, comme le prouvent les expériences camerounaises et étrangères. Certes, son coût et le risque de sélection précoce sont évoqués, cependant ces objections ne résistent pas à l'examen : l'orientation technique est loin d'être un échec. L'auteur conclut donc que le législateur doit instituer cet enseignement, profitable aux jeunes comme à la nation toute entière.}

\textbf{Vérification :}
\begin{itemize}
\item \textbf{Mots :} 100 exactement (contrainte 90-110 respectée).
\item \textbf{Connecteurs (4) :} \emph{C'est pourquoi} (conséquence), \emph{afin de} (but), \emph{En effet} (renforcement), \emph{Certes... cependant} (concession/opposition), \emph{donc} (conclusion).
\item \textbf{Fidélité :} Thèse respectée, ordre des idées conservé.
\item \textbf{Reformulation :} Aucune phrase recopiée ; vocabulaire propre (pénurie, imposé, remédier, loin d'être un échec).
\item \textbf{Élagage :} Noms propres, chiffres, citation, listes techniques disparus.
\end{itemize}

\textbf{Étape 6 : auto-évaluation.} Les six critères de la grille sont validés ; le groupe peut afficher sa copie et défendre ses choix de reformulation devant la classe.
\end{exoresolu}

\courssub{Compte rendu, remédiation et entraînement à l'examen}

\begin{exercice}[Entraînement Probatoire : Résumé d'un texte court]
Résumez le texte suivant en \textbf{50 mots ($\pm$ 10 \%)} (soit 45 à 55 mots). Appliquez la méthode en 6 étapes.

\textit{« L'intelligence artificielle (IA) transforme nos métiers. Certes, elle automatise des tâches répétitives, libérant du temps pour la créativité. Cependant, elle menace des emplois peu qualifiés. Face à ce défi, la formation continue devient indispensable. Les gouvernements doivent investir massivement dans l'éducation numérique dès l'école primaire. Seule une adaptation constante des compétences permettra à la jeunesse africaine de saisir les opportunités de la quatrième révolution industrielle, au lieu de la subir. » (Environ 85 mots)}
\end{exercice}

\begin{corrige}[Exercice n°1 : Entraînement Probatoire]
\textbf{Structure :} §1 Constat/Thèse (IA transforme, double visage) $\rightarrow$ §2 Argument (menace/opportunité) $\rightarrow$ §3 Solution/Conclusion (Formation continue, investissement État).
\textbf{Idées essentielles :} 1. L'IA double les métiers (gain créativité / perte emplois). 2. La formation continue est la réponse. 3. L'État doit éduquer au numérique tôt. 4. L'adaptation conditionne l'avenir de la jeunesse africaine.
\textbf{Résumé proposé (51 mots) :}
\emph{L'intelligence artificielle bouleverse le travail : elle libère la créativité mais supprime des emplois peu qualifiés. Par conséquent, la formation continue s'impose comme une nécessité. Les États doivent donc instaurer l'éducation numérique dès le primaire. Seule une mise à jour permanente des compétences permettra à la jeunesse africaine de maîtriser la quatrième révolution industrielle au lieu de la subir.}
\textbf{Vérification :} 51 mots. Connecteurs : \emph{mais, par conséquent, donc, au lieu de}. Reformulation : \emph{transforme} $\rightarrow$ \emph{bouleverse} ; \emph{automatise... libérant} $\rightarrow$ \emph{libère} ; \emph{investir massivement} $\rightarrow$ \emph{instaurer} ; \emph{saisir les opportunités... au lieu de la subir} $\rightarrow$ \emph{maîtriser... au lieu de la subir}.
\end{corrige}

\begin{aretenir}[Check-list finale pour l'épreuve du Bac Tech (Résumé)]
\begin{itemize}
\item [ ] Ai-je lu le texte deux fois ?
\item [ ] Ai-je fait le tableau structure / idées essentielles ?
\item [ ] Ai-je barré au brouillon tous les exemples, chiffres, citations, noms propres ?
\item [ ] Ai-je reformulé \textbf{chaque} idée essentielle (synonymie / conversion) ?
\item [ ] Ai-je utilisé \textbf{au moins 3 connecteurs} variés ?
\item [ ] Ai-je compté mes mots \textbf{deux fois} (brouillon + propre) ?
\item [ ] Suis-je dans la fourchette $\pm 10\%$ ?
\item [ ] Relu : orthographe, accords, ponctuation, fidélité au texte ?
\end{itemize}
\end{aretenir}

\courssub{Bilan de la leçon}

Cette leçon a permis de mobiliser, dans une situation authentique de la vie scolaire (« Lu dans la presse »), l'ensemble des savoir-faire de l'unité. Elle a d'abord montré qu'un éditorial n'est pas un simple alignement de phrases : il obéit à une \cle{structure argumentative} repérable (constat, thèse, arguments, réfutation, conclusion), et c'est cette architecture qui guide le dégagement des idées essentielles. Elle a ensuite confirmé que l'identification de la \cle{tonalité} — ici satirique — repose sur des indices linguistiques observables (ironie, antithèse, hyperbole) et non sur une impression vague. Elle a enfin entraîné les élèves à la \cle{reformulation} rigoureuse : synonymie, antonymie et conversion grammaticale permettent de dire autrement sans trahir, condition première d'un résumé réussi au Baccalauréat technique.

La comparaison collective des résumés affichés a mis en lumière un point essentiel : plusieurs condensations correctes sont possibles, mais toutes doivent respecter les mêmes exigences — fidélité à la pensée de l'auteur, suppression des exemples et des chiffres, cohérence assurée par les connecteurs, respect strict du nombre de mots. En prolongement, chaque élève pourra transférer cette compétence hors du cadre scolaire : résumer en cinq lignes, pour un camarade absent, le compte rendu d'une réunion de la coopérative scolaire ou de la tontine de quartier, en appliquant la même démarche en six étapes. La même rigueur s'applique à l'analyse de \emph{Le Lion et la Perle} : résumer une scène, c'est en extraire la structure dramaturgique (thèse du personnage, arguments scéniques, tonalité) pour la reformuler avec précision.

\newpage

\coursec{Méthodes et exercices résolus}

\coursec{Leçon 19 : Formuler les idées essentielles d'un texte à résumer}

\courssub{Le texte argumentatif : caractéristiques et structure}

\begin{definition}[Texte argumentatif]
Un \cle{texte argumentatif} a pour but de convaincre ou de persuader le lecteur en défendant une thèse à l'aide d'arguments organisés logiquement. Il se distingue du texte explicatif (qui informe) et du texte narratif (qui raconte).
\end{definition}

\begin{propriete}[Caractéristiques du texte argumentatif]
\begin{itemize}
\item \textbf{Une thèse} : idée principale, opinion défendue (explicite ou implicite).
\item \textbf{Des arguments} : raisons logiques qui soutiennent la thèse (faits, valeurs, autorité).
\item \textbf{Des exemples/illustrations} : cas concrets, anecdotes, chiffres, citations qui appuient les arguments.
\item \textbf{Des connecteurs logiques} : marqueurs de la progression (cause, conséquence, opposition, énumération).
\item \textbf{Une structure organisée} : plan déductif (thèse $\rightarrow$ arguments), inductif (arguments $\rightarrow$ thèse), ou dialectique (thèse/antithèse/synthèse).
\end{itemize}
\end{propriete}

\begin{methode}[Repérer la structure argumentative d'un texte]
Pour dégager la structure argumentative d'un texte avant de le résumer, suivre les étapes suivantes :
\begin{enumerate}
\item \textbf{Repérer la thèse} : c'est l'idée principale défendue par l'auteur. Elle se trouve souvent au début (plan déductif) ou à la fin (plan inductif). Chercher les marques : \emph{je pense que}, \emph{il faut admettre que}, \emph{en vérité}, \emph{certes}...
\item \textbf{Dégager les arguments} : ce sont les raisons qui justifient la thèse. Les repérer grâce aux connecteurs logiques : \emph{d'abord, ensuite, enfin, car, parce que, en effet, puisque, par conséquent}.
\item \textbf{Identifier les exemples} : ils illustrent les arguments (faits, anecdotes, citations, chiffres). Ils sont introduits par \emph{par exemple, ainsi, c'est le cas de, notamment}.
\item \textbf{Schématiser la progression} : noter sur un tableau à trois colonnes : thèse / arguments / exemples.
\item \textbf{Vérifier la cohérence} : chaque argument doit servir la thèse ; un élément qui ne sert pas la démonstration sera éliminé au moment du résumé.
\end{enumerate}
\textbf{Formule à retenir} : Thèse + Arguments + Exemples = structure argumentative. Au résumé, on garde la thèse et les arguments, on supprime les exemples.
\end{methode}

\begin{exoresolu}[Dégager la structure argumentative d'un paragraphe]
\textbf{Énoncé.} Lisez le texte suivant, puis dégagez sa structure argumentative (thèse, arguments, exemples) sous forme de tableau.

\emph{« L'école demeure le meilleur instrument de promotion sociale en Afrique. D'abord, elle donne à chaque enfant, quelle que soit son origine, les mêmes chances de réussite : un fils de paysan peut devenir médecin ou ingénieur. Ensuite, elle forme les cadres dont le pays a besoin pour son développement ; sans techniciens ni enseignants qualifiés, aucune industrie ne peut prospérer. Enfin, elle éduque le citoyen en lui apprenant ses droits et ses devoirs. C'est pourquoi tout investissement dans l'éducation est un investissement dans l'avenir. »}

\tcblower
\textbf{Solution détaillée.}

\textbf{Étape 1 : repérage de la thèse.} La première phrase affirme une position générale : « L'école demeure le meilleur instrument de promotion sociale en Afrique. » C'est la \cle{thèse}. La dernière phrase (« C'est pourquoi... ») la reformule en conclusion, ce qui confirme le repérage.

\textbf{Étape 2 : repérage des arguments.} Les connecteurs \emph{D'abord}, \emph{Ensuite}, \emph{Enfin} introduisent trois arguments :
\begin{itemize}
\item Argument 1 : l'école donne les mêmes chances de réussite à tous (égalité des chances) ;
\item Argument 2 : l'école forme les cadres nécessaires au développement (formation) ;
\item Argument 3 : l'école éduque le citoyen (citoyenneté).
\end{itemize}

\textbf{Étape 3 : repérage des exemples.} « Un fils de paysan peut devenir médecin ou ingénieur » illustre l'argument 1 ; « sans techniciens ni enseignants qualifiés, aucune industrie ne peut prospérer » illustre l'argument 2. L'argument 3 n'a pas d'exemple explicite.

\textbf{Étape 4 : tableau récapitulatif.}
\begin{center}
\begin{tabularx}{\linewidth}{|l|X|X|}\hline
\rowcolor{popblueL} \textbf{Élément} & \textbf{Contenu} & \textbf{Repère dans le texte} \\ \hline
Thèse & L'école est le meilleur instrument de promotion sociale & Première et dernière phrases \\ \hline
Argument 1 & Égalité des chances & « D'abord » \\ \hline
Argument 2 & Formation des cadres & « Ensuite » \\ \hline
Argument 3 & Éducation du citoyen & « Enfin » \\ \hline
Exemples & Fils de paysan devenu médecin ; techniciens et industrie & Après arguments 1 et 2 \\ \hline
\end{tabularx}
\end{center}

\textbf{Conclusion.} Le texte suit un plan déductif : la thèse est annoncée dès la première phrase, puis démontrée par trois arguments hiérarchisés, et reprise en conclusion. Pour le résumé, on conservera la thèse et les trois arguments ; on éliminera les exemples.
\end{exoresolu}

\courssub{La reformulation des idées essentielles : synonymie, antonymie et techniques}

\begin{methode}[Reformuler sans trahir la pensée de l'auteur]
La reformulation est le cœur du résumé. Procéder ainsi :
\begin{enumerate}
\item \textbf{Comprendre avant de transformer} : lire le passage jusqu'à pouvoir l'expliquer avec ses propres mots à l'oral.
\item \textbf{Remplacer les mots par des synonymes de même registre} : \emph{demeure} $\rightarrow$ reste ; \emph{instrument} $\rightarrow$ outil, moyen ; \emph{prospérer} $\rightarrow$ se développer.
\item \textbf{Changer la construction des phrases} : passer de la voix active à la voix passive, d'une phrase affirmative à une tournure nominale, fusionner deux phrases courtes en une seule, transformer une phrase négative en positive.
\item \textbf{Généraliser les exemples} : « un fils de paysan peut devenir médecin » $\rightarrow$ « tout enfant peut accéder aux professions valorisées ».
\item \textbf{Supprimer les procédés d'amplification} : répétitions, images, énumérations, questions rhétoriques, adresses au lecteur, concessions.
\item \textbf{Vérifier la fidélité} : la reformulation ne doit ni ajouter une idée personnelle, ni affaiblir, ni exagérer la pensée de l'auteur.
\end{enumerate}
\textbf{Réflexes} : jamais de « je » personnel ; jamais de copie de plus de trois mots consécutifs du texte ; toujours relire en comparant au texte source.
\end{methode}

\begin{methode}[Mobiliser synonymes et antonymes]
\begin{enumerate}
\item \textbf{Identifier le mot-clé} de la phrase à reformuler (nom, verbe ou adjectif porteur du sens).
\item \textbf{Chercher un synonyme de même registre} : un synonyme doit avoir un sens proche \emph{et} le même niveau de langue. On ne remplace pas \emph{ouvrage} (soutenu) par \emph{bouquin} (familier).
\item \textbf{Utiliser l'antonymie par la négation} : reformuler une idée par la négation du contraire : \emph{« il est difficile de réussir sans travailler »} $\rightarrow$ \emph{« le travail est indispensable à la réussite »} (difficile/sans $\rightarrow$ indispensable).
\item \textbf{Vérifier dans le contexte} : deux mots synonymes ne sont pas toujours interchangeables ; \emph{regarder} et \emph{observer} ne conviennent pas aux mêmes situations.
\item \textbf{Éviter la répétition} : dans un résumé, ne jamais répéter le même terme dans deux phrases voisines ; varier grâce aux synonymes.
\end{enumerate}
\textbf{À retenir} : synonymie = sens proche, même registre ; antonymie = sens contraire, utile avec la négation pour changer de construction.
\end{methode}

\begin{exoresolu}[Synonymes, antonymes et reformulation de phrases]
\textbf{Énoncé.}
1) Donnez deux synonymes de registre courant pour chacun des mots suivants : \emph{demeurer}, \emph{prospérer}, \emph{instruire}, \emph{instrument}.
2) Reformulez les phrases suivantes en changeant le vocabulaire et la construction, sans en altérer le sens.
a) \emph{« L'école demeure le meilleur instrument de promotion sociale en Afrique. »}
b) \emph{« Sans techniciens ni enseignants qualifiés, aucune industrie ne peut prospérer. »}
c) \emph{« Comment un peuple pourrait-il se développer s'il refuse d'instruire ses enfants ? »}
d) \emph{« L'ignorance est l'ennemie du progrès. »} (par antonymie)

\tcblower
\textbf{Solution détaillée.}

\textbf{1) Synonymes :}
\begin{itemize}
\item \emph{demeurer} : rester, continuer d'être ;
\item \emph{prospérer} : se développer, fleurir ;
\item \emph{instruire} : éduquer, former ;
\item \emph{instrument} : outil, moyen.
\end{itemize}
Chaque synonyme proposé appartient au registre courant, comme les mots de départ : ils sont donc interchangeables dans un résumé sans faute de niveau de langue.

\textbf{2) Reformulations :}
\begin{itemize}
\item \textbf{a)} On remplace \emph{demeure} par \emph{reste}, \emph{instrument} par \emph{moyen}, \emph{promotion sociale} par \emph{ascension sociale}, et on inverse la construction : « En Afrique, l'ascension sociale repose avant tout sur l'école. »
\item \textbf{b)} On transforme la tournure négative en tournure positive et on remplace \emph{prospérer} par \emph{se développer} : « Le développement de l'industrie exige des techniciens et des enseignants compétents. »
\item \textbf{c)} La question rhétorique est un procédé d'amplification : on la transforme en affirmation : « Un peuple qui néglige l'instruction de ses enfants ne peut pas se développer. »
\item \textbf{d)} L'antonyme de \emph{ignorance} est \emph{connaissance} (ou \emph{instruction}) ; l'antonyme de \emph{ennemie} est \emph{alliée} (ou \emph{condition}). En passant par la négation : « Sans instruction, il n'y a pas de progrès. » Ou par le contraire positif : « La connaissance est la condition du progrès. »
\end{itemize}
\textbf{Vérification.} Chaque reformulation respecte les trois exigences du résumé : fidélité au sens, changement des mots et des constructions, suppression des procédés expressifs.
\end{exoresolu}

\courssub{Les tonalités : comique, dramatique, satirique, tragique}

\begin{definition}[Tonalité]
La \cle{tonalité} est l'émotion dominante que le texte cherche à provoquer chez le lecteur (rire, peur, pitié, indignation...). Elle résulte du choix des mots, des figures de style, de la syntaxe et du registre de langue.
\end{definition}

\begin{methode}[Reconnaître et justifier les quatre tonalités au programme]
\begin{enumerate}
\item \textbf{Le comique} : il provoque le rire ou le sourire. \textbf{Repères} : exagération, répétition, quiproquo, jeux de mots, ridicule d'un personnage, chute inattendue. \textbf{Effet} : détente, amusement.
\item \textbf{Le satirique} : il rit \emph{pour critiquer} un défaut, un vice ou une institution. \textbf{Repères} : ironie, caricature, dénonciation sous le rire. \textbf{Question-clé} : « Qui ou quoi l'auteur critique-t-il ? » \textbf{Effet} : rire critique, prise de conscience.
\item \textbf{Le dramatique} : il crée la tension et l'émotion devant un conflit grave. \textbf{Repères} : opposition entre personnages, enjeu important, suspense, passion, dialogue vif. \textbf{Effet} : angoisse, empathie, tension.
\item \textbf{Le tragique} : il montre un personnage écrasé par une force supérieure (destin, fatalité, dieux, société), sans issue possible. \textbf{Repères} : mort annoncée, fatalité, impossibilité d'échapper, registre pathétique, impuissance du héros. \textbf{Effet} : pitié, terreur, sentiment d'inéluctable.
\end{enumerate}
\textbf{Démarche de justification (exigible au Bac)} : Nommer la tonalité $\rightarrow$ Citer un procédé précis du texte $\rightarrow$ Expliquer l'effet produit sur le lecteur. Formule : \emph{« La tonalité est ... car le texte utilise ... ce qui provoque chez le lecteur ... »}
\end{methode}

\begin{exoresolu}[Identifier et justifier une tonalité]
\textbf{Énoncé.} Pour chacun des passages suivants, indiquez la tonalité dominante et justifiez votre réponse en vous appuyant sur un procédé précis.

a) \emph{« Le ministre déclara solennellement que la caisse était vide, tout en remplissant ses poches devant les caméras. »}

b) \emph{« Elle savait que le bateau partirait à l'aube et que, quoi qu'elle fît, elle ne reverrait jamais son fils. »}

c) \emph{« Monsieur Pandore entra dans la salle en saluant profondément... une chaise qu'il avait prise pour le directeur. »}

\tcblower
\textbf{Solution détaillée.}

\textbf{a) Tonalité satirique.} Le passage dénonce l'hypocrisie et la corruption d'un homme politique par le contraste ironique entre sa déclaration (« la caisse était vide ») et son acte (« remplissant ses poches »). Le rire sert ici la critique sociale : c'est la marque du satirique et non du simple comique.

\textbf{b) Tonalité tragique.} Le personnage est confronté à une fatalité sans issue : « quoi qu'elle fît, elle ne reverrait jamais son fils ». L'adverbe \emph{jamais} et l'impuissance du personnage créent le pathétique propre au tragique : le lecteur éprouve pitié devant un malheur inévitable.

\textbf{c) Tonalité comique.} Le rire naît du quiproquo : le personnage salue une chaise prise pour le directeur. La chute inattendue de la phrase et le ridicule de la méprise relèvent du comique de situation.

\textbf{Conclusion.} Pour identifier une tonalité, il ne suffit pas de la nommer : il faut citer le procédé (ironie, fatalité, quiproquo) et expliquer l'effet produit (critique, pitié, rire). C'est cette justification qui rapporte des points à l'examen.
\end{exoresolu}

\courssub{Lecture méthodique de \emph{La Lionne et la Perle} (Wole Soyinka) — n°4 et n°5}

\begin{methode}[Lecture méthodique d'un extrait théâtral]
Pour analyser un extrait de \emph{La Lionne et la Perle} (ou tout extrait théâtral au programme) :
\begin{enumerate}
\item \textbf{Situer l'extrait} : replacer le passage dans l'œuvre (qui parle ? à quel moment de l'intrigue ? quel acte, quelle scène ?).
\item \textbf{Dégager l'enjeu} : quel conflit oppose les personnages ? (ex : tradition contre modernité, le chef Baroka contre l'instituteur Lakunle au sujet de Sidi).
\item \textbf{Analyser les procédés} : tonalité (comique, satirique...), registre de langue (soutenu, familier), gestuelle et didascalies (indications scéniques), images, proverbes.
\item \textbf{Dégager la signification} : que critique ou que défend l'auteur ? Quelle leçon pour la société ?
\item \textbf{Faire le lien avec la vie courante} : appliquer la leçon du texte à une situation concrète (débat tradition/modernité au village, rôle de la femme, éducation, dot).
\end{enumerate}
\textbf{Réflexe Bac} : une lecture méthodique réussie va toujours du texte (procédés) vers le sens (enjeu, signification), jamais l'inverse.
\end{methode}

\begin{exoresolu}[Lecture méthodique d'un extrait et application à une situation concrète]
\textbf{Énoncé.} Dans un extrait de \emph{La Lionne et la Perle}, Lakunle, jeune instituteur, refuse de payer la dot pour épouser Sidi, qu'il accuse d'être « une fille arriérée » attachée aux coutumes ; Sidi lui réplique avec vivacité et défend la tradition du village.

1) Dégagez l'enjeu de cet extrait. 2) Identifiez la tonalité dominante et justifiez-la. 3) Montrez comment ce débat éclaire une situation de la vie courante au Cameroun.

\tcblower
\textbf{Solution détaillée.}

\textbf{1) L'enjeu.} L'extrait met en scène le conflit entre la \cle{modernité} incarnée par Lakunle (école, rejet de la dot, idées venues d'ailleurs, langage ampoulé) et la \cle{tradition} défendue par Sidi (coutumes du village, fierté de son statut de jeune fille, bon sens concret). L'enjeu est donc culturel et social : faut-il abandonner les coutumes au nom d'un progrès mal compris ?

\textbf{2) La tonalité.} La tonalité dominante est \cle{satirique} : Soyinka se moque de Lakunle, dont la prétendue modernité n'est que vanité et verbiage — il humilie celle qu'il prétend aimer en la traitant d'« arriérée ». L'ironie consiste à faire du « moderniste » un personnage ridicule, tandis que Sidi, la « traditionnelle », fait preuve de bon sens, de dignité et de fierté. L'auteur critique ainsi le faux modernisme qui méprise les valeurs africaines sans les comprendre.

\textbf{3) Application à la vie courante.} Ce débat éclaire des situations camerounaises actuelles : les discussions sur la dot dans les familles (montant, signification), le retour au village des jeunes diplômés qui méprisent parfois les coutumes ancestrales, ou encore les conflits entre générations sur le mariage et le rôle de la femme. La leçon de Soyinka est qu'une modernité authentique ne consiste pas à rejeter en bloc la tradition, mais à dialoguer avec elle : on peut être instruit \emph{et} respectueux de sa culture.

\textbf{Conclusion.} La lecture méthodique a permis de passer du texte (échange Lakunle/Sidi) au sens (critique du faux modernisme), puis d'appliquer la leçon de l'œuvre à des situations concrètes de la vie courante, comme l'exige le savoir-faire du programme.
\end{exoresolu}

\courssub{La contraction de texte (résumé) : méthode complète et application}

\begin{methode}[Les étapes du résumé au Probatoire / Baccalauréat technique]
\begin{enumerate}
\item \textbf{Lecture analytique} (2 lectures minimum) : souligner la thèse, les arguments, les connecteurs ; écarter exemples, répétitions, digressions, concessions, questions rhétoriques.
\item \textbf{Plan schématique} : rédiger le schéma thèse/arguments au brouillon (tableau ou liste).
\item \textbf{Reformulation} : rédiger en phrases complètes, avec ses propres mots (synonymes, antonymes, changement de construction), en reliant les idées par des connecteurs (\emph{d'abord, ensuite, enfin, c'est pourquoi, par conséquent}).
\item \textbf{Respect de la longueur imposée} : compter les mots du texte source. Calculer la longueur cible (ex : quart = total ÷ 4). Une marge de \textbf{10 \%} est tolérée (ex : pour 30 mots, entre 27 et 33 mots).
\item \textbf{Neutralité absolue} : aucun jugement personnel, aucun « je », aucun commentaire, aucune connaissance extérieure ; le résumé rend compte de la pensée de l'auteur, pas de la vôtre.
\item \textbf{Relecture finale} : vérifier la fidélité au sens, l'orthographe, la liaison des idées (cohérence), et le nombre de mots exact.
\end{enumerate}
\textbf{Formule à retenir} : Résumé = Thèse + Arguments essentiels, reformulés, liés, à la longueur imposée, sans exemples ni opinion personnelle.
\end{methode}

\begin{exoresolu}[Résumer un texte argumentatif au quart de sa longueur]
\textbf{Énoncé.} Résumez le texte suivant (120 mots) en 30 mots environ.

\emph{« Le sport scolaire n'est pas une perte de temps, contrairement à ce que prétendent certains parents. En effet, il développe d'abord la santé physique des élèves : un adolescent qui court, qui nage, qui joue au ballon fortifie son cœur et ses muscles. Ensuite, il enseigne des valeurs indispensables comme le respect des règles, l'esprit d'équipe et le goût de l'effort, valeurs que la vie professionnelle exigera demain. Enfin, il améliore les résultats scolaires, car un corps sain abrite un esprit vigilant : les études montrent que les élèves sportifs concentrent mieux leur attention en classe. Il faut donc, loin de le supprimer, renforcer le sport dans nos établissements. »}

\tcblower
\textbf{Solution détaillée.}

\textbf{Étape 1 : structure du texte.}
\begin{itemize}
\item Thèse : le sport scolaire est utile et doit être renforcé (première et dernière phrases).
\item Argument 1 : il développe la santé physique.
\item Argument 2 : il enseigne des valeurs (respect, esprit d'équipe, effort).
\item Argument 3 : il améliore les résultats scolaires.
\item À éliminer : la concession (« contrairement à... »), les exemples (courir, nager, ballon), la référence aux études, les amplifications (« un corps sain... »).
\end{itemize}

\textbf{Étape 2 : reformulation et rédaction.} On reformule chaque élément et on relie par des connecteurs :

« Loin d'être inutile, le sport scolaire est bénéfique : il améliore la santé des élèves, leur transmet des valeurs essentielles et favorise leur réussite ; il convient donc de le développer. »

\textbf{Étape 3 : comptage.}
Loin(1) d'être(2) inutile(3) le(4) sport(5) scolaire(6) est(7) bénéfique(8) il(9) améliore(10) la(11) santé(12) des(13) élèves(14) leur(15) transmet(16) des(17) valeurs(18) essentielles(19) et(20) favorise(21) leur(22) réussite(23) il(24) convient(25) donc(26) de(27) le(28) développer(29).
Soit \textbf{29 mots}, conforme à la consigne (30 mots $\pm$ 10 \%, soit 27 à 33 mots).

\textbf{Étape 4 : vérification.} La thèse et les trois arguments sont présents ; aucun exemple n'est conservé ; aucune opinion personnelle n'est ajoutée ; aucun groupe de plus de trois mots n'est copié du texte.

\textbf{Conclusion.} Le résumé produit est fidèle, bref, personnel dans sa formulation et conforme à la longueur exigée : il remplit les quatre critères d'évaluation de la contraction de texte.
\end{exoresolu}

\courssub{Compte rendu, remédiation et entraînement à l'examen}

\begin{attention}[Les 5 erreurs qui coûtent des points au Probatoire / Bac Technique]
\begin{enumerate}
\item \textbf{Copier des phrases du texte} dans le résumé : reprendre plus de trois mots consécutifs sans reformuler est sanctionné. Reformulez toujours avec des synonymes et une construction nouvelle.
\item \textbf{Donner son avis personnel} : « je pense que », « à mon avis », « selon moi » sont interdits dans un résumé. Le résumé rend compte de la pensée de l'auteur, pas de la vôtre.
\item \textbf{Conserver les exemples et les détails} : anecdotes, chiffres, énumérations, comparaisons, citations et concessions doivent disparaître ; seuls la thèse et les arguments essentiels subsistent.
\item \textbf{Négliger la longueur imposée} : ne pas compter les mots ou dépasser largement la limite (plus de 10 \% d'écart) fait perdre des points. Comptez les mots du texte, divisez selon la consigne, puis comptez les mots de votre résumé.
\item \textbf{Nommer une tonalité sans la justifier} : écrire « la tonalité est comique » sans citer le procédé (quiproquo, ironie, fatalité) ni expliquer l'effet produit ne rapporte qu'une partie des points. Toujours appliquer la formule : \textbf{Tonalité + Procédé cité + Effet sur le lecteur}.
\end{enumerate}
\end{attention}

\begin{aretenir}[Check-list avant de rendre sa copie]
\begin{itemize}
\item[$\square$] Ai-je lu le texte deux fois ?
\item[$\square$] Ai-je identifié la thèse et les arguments (tableau au brouillon) ?
\item[$\square$] Ai-je reformulé \emph{toutes} les phrases (synonymes, structure changée) ?
\item[$\square$] Ai-je supprimé exemples, répétitions, « je », opinions ?
\item[$\square$] Ai-je compté les mots de mon résumé et vérifié la marge de 10 \% ?
\item[$\square$] Pour la tonalité : ai-je cité le procédé et expliqué l'effet ?
\item[$\square$] Relecture orthographe / grammaire / connecteurs logiques.
\end{itemize}
\end{aretenir}

\begin{exercice}[Entraînement 1 : Structure et reformulation]
\textbf{Texte support (85 mots) :}
\emph{« La déforestation menace gravement l'équilibre écologique de notre pays. En premier lieu, la disparition des arbres laisse le sol nu face aux pluies torrentielles, ce qui provoque une érosion massive. De plus, la perte du couvert forestier réduit la capacité de la nature à absorber le dioxyde de carbone, aggravant ainsi le réchauffement climatique. Enfin, la biodiversité s'effondre : les animaux perdent leur habitat naturel et disparaissent. Il est donc urgent de reboiser massivement pour sauver notre environnement. »}

\textbf{Consignes :}
1) Dégagez la thèse et les arguments sous forme de tableau.
2) Reformulez l'argument 2 en une seule phrase (changement de vocabulaire et de construction).
3) Proposez un titre au texte.
\end{exercice}

\begin{corrige}[Exercice 1]
\textbf{1) Structure :}
\begin{center}
\begin{tabularx}{\linewidth}{|l|X|X|}\hline
\rowcolor{popblueL} \textbf{Élément} & \textbf{Contenu} & \textbf{Repère} \\ \hline
Thèse & La déforestation menace l'équilibre écologique / Urgence de reboiser & 1ère et dernière phrases \\ \hline
Arg 1 & Érosion des sols (sol nu + pluies) & « En premier lieu » \\ \hline
Arg 2 & Aggravation réchauffement (moins d'absorption CO2) & « De plus » \\ \hline
Arg 3 & Effondrement biodiversité (perte habitat) & « Enfin » \\ \hline
\end{tabularx}
\end{center}

\textbf{2) Reformulation arg 2 :} « La réduction de la forêt diminue l'absorption du CO2, ce qui accélère le réchauffement de la planète. » (ou : « Le réchauffement climatique s'intensifie car la forêt, puits de carbone, disparaît. »)

\textbf{3) Titre :} « Les dangers de la déforestation » ou « Reboisement : une urgence écologique ».
\end{corrige}

\begin{exercice}[Entraînement 2 : Contraction de texte (Probatoire)]
\textbf{Consigne :} Résumez le texte ci-dessous en \textbf{25 mots environ} (texte source : 100 mots).

\emph{« L'usage excessif des écrans chez les jeunes inquiète les spécialistes de la santé. D'une part, la lumière bleue émise par les smartphones perturbe la sécrétion de mélatonine, l'hormone du sommeil, entraînant des insomnies chroniques. D'autre part, la position statique prolongée favorise les troubles musculo-squelettiques, notamment au niveau du cou et du dos. Par ailleurs, l'isolement social guette l'adolescent qui préfère le monde virtuel aux relations réelles. Face à ces risques, les médecins recommandent de limiter le temps d'écran à deux heures par jour et d'encourager les activités physiques. »}
\end{exercice}

\begin{corrige}[Exercice 2]
\textbf{Structure :} Thèse = Écrans = danger pour santé jeunes. Args = 1. Sommeil (lumière bleue/mélatonine). 2. Physique (posture/douleurs). 3. Social (isolement). Recommandation = Limiter à 2h + sport.

\textbf{Proposition de résumé (25 mots) :}
« L'abus d'écrans nuit gravement aux jeunes : troubles du sommeil par lumière bleue, douleurs musculaires par sédentarité, isolement relationnel ; les médecins préconisent une limitation à deux heures quotidiennes et la pratique sportive. »

\textbf{Comptage :} L'abus(1) d'écrans(2) nuit(3) gravement(4) aux(5) jeunes(6) : troubles(7) du(8) sommeil(9) par(10) lumière(11) bleue(12), douleurs(13) musculaires(14) par(15) sédentarité(16), isolement(17) relationnel(18) ; les(19) médecins(20) préconisent(21) une(22) limitation(23) à(24) deux(25) heures(26) quotidiennes(27) et(28) la(29) pratique(30) sportive(31). $\rightarrow$ 31 mots. Légèrement au-dessus (cible 25 $\pm$ 10\% = 22-28). Ajustement possible : « L'abus d'écrans nuit aux jeunes : insomnies (lumière bleue), douleurs (sédentarité), isolement ; médecins conseillent max 2h/jour et sport. » (24 mots). Conforme.
\end{corrige}

\begin{exercice}[Entraînement 3 : Tonalités et lecture d'image/texte court]
\textbf{Identifiez la tonalité de chaque extrait et justifiez par un procédé.}
a) \emph{« Le dictateur, dans son palais doré, signait des décrets de misère pour son peuple qu'il appelait "mes enfants chéris". »}
b) \emph{« Le jeune homme hésita une seconde au bord du gouffre. Le vent hurlait. Il savait qu'aucun filet ne l'attendait en bas. Il avança. »}
c) \emph{« Monsieur Dupont, pressé, mit ses chaussures... sur les mains. Il partit au travail, fier de son allure. »}
\end{exercice}

\begin{corrige}[Exercice 3]
\textbf{a) Satirique.} Ironie : contraste entre « palais doré » / « mes enfants chéris » et « décrets de misère ». Critique du pouvoir hypocrite.
\textbf{b) Tragique.} Fatalité : « aucun filet », « il avança » (inéluctable), « gouffre », « vent hurlait ». Personnage impuissant face à la mort annoncée. Pathétique.
\textbf{c) Comique.} Quiproquo / absurdité : chaussures sur les mains. Chute inattendue. Ridicule du personnage. Rire pur sans critique sociale profonde.
\end{corrige}

\newpage

\coursec{Exercices}

\courssub{Je vérifie mes connaissances}

\begin{exercice}[QCM : Le texte argumentatif et ses composantes]
Parmi les propositions suivantes, une seule est exacte pour chaque question. Recopiez le numéro et la lettre correspondante.

\textbf{1.} Dans un texte argumentatif, la \cle{thèse} correspond à :
\begin{itemize}
    \item A. L'opinion personnelle de l'auteur sans justification.
    \item B. L'idée principale que l'auteur défend et qu'il cherche à faire admettre.
    \item C. Un exemple concret illustrant le propos.
    \item D. La concession faite à l'adversaire.
\end{itemize}

\textbf{2.} Lequel de ces connecteurs n'exprime pas une relation de \cle{cause} ?
\begin{itemize}
    \item A. Parce que
    \item B. Puisque
    \item C. Cependant
    \item D. Comme
\end{itemize}

\textbf{3.} La \cle{structure argumentative} classique (thèse, arguments, exemples, conclusion) vise à :
\begin{itemize}
    \item A. Divertir le lecteur par une narration.
    \item B. Informer objectivement sans prendre parti.
    \item C. Convaincre ou persuader le destinataire.
    \item D. Décrire un paysage ou un personnage.
\end{itemize}

\textbf{4.} Dans une \cle{contraction de texte} (résumé), on doit :
\begin{itemize}
    \item A. Copier-coller les phrases les plus importantes.
    \item B. Conserver les exemples détaillés pour la preuve.
    \item C. Reformuler les idées essentielles en supprimant les détails et les répétitions.
    \item D. Donner son avis personnel sur le sujet traité.
\end{itemize}
\end{exercice}

\begin{exercice}[Vrai ou Faux justifié : Tonalités et Lexique]
Indiquez si les affirmations suivantes sont vraies ou fausses. Justifiez brièvement chaque réponse par une définition ou un contre-exemple.

\begin{enumerate}
    \item La \cle{tonalité satirique} vise uniquement à faire rire le lecteur sans aucune critique sociale.
    \item L'\cle{antonymie} est la relation sémantique qui unit deux mots de sens opposé (ex. : progrès / régression).
    \item Dans \emph{Le Lion et la Perle}, le personnage de Lakunle incarne exclusivement la tonalité tragique.
    \item La \cle{reformulation} par synonymie consiste à remplacer un mot par un autre de sens identique dans tous les contextes.
    \item Un résumé réussi doit respecter la limite de mots imposée (ex. : $\pm 10\%$) et l'ordre logique des idées du texte source.
\end{enumerate}
\end{exercice}

\begin{exercice}[Définitions et notions clés]
Donnez une définition précise et un exemple pour chacun des termes suivants :
\begin{enumerate}
    \item \cle{Argument d'autorité}
    \item \cle{Concession} (dans la structure argumentative)
    \item \cle{Tonalité dramatique}
    \item \cle{Conversion grammaticale} (outil de reformulation)
    \item \cle{Idée essentielle} (vs idée secondaire)
\end{enumerate}
\end{exercice}

\begin{exercice}[Application directe : Connecteurs et Structure]
Voici une liste de connecteurs logiques. Classez-les dans le tableau ci-dessous selon leur fonction logique : \textbf{Addition}, \textbf{Opposition}, \textbf{Cause}, \textbf{Conséquence}, \textbf{Concession}.

\medskip
\textit{Liste :} \textbf{De plus ; Or ; Par conséquent ; Bien que ; En effet ; Ainsi ; Néanmoins ; D'ailleurs ; Car ; Du fait que.}

\medskip
\begin{center}
\begin{tabularx}{\linewidth}{|l|X|}\hline
\textbf{Fonction logique} & \textbf{Connecteurs classés} \\ \hline
Addition & \\ \hline
Opposition & \\ \hline
Cause & \\ \hline
Conséquence & \\ \hline
Concession & \\ \hline
\end{tabularx}
\end{center}
\end{exercice}

\courssub{Je m'entraîne}

\begin{exercice}[Repérage dans un texte argumentatif : L'agriculture au Cameroun]
Lisez l'extrait d'éditorial suivant (environ 180 mots) :

\begin{quote}
\textbf{L'agriculture, pilier de notre économie}

Il est indéniable que l'agriculture demeure le moteur principal de l'économie camerounaise. \textbf{En effet}, elle emploie plus de 60\% de la population active et contribue à hauteur de 20\% au PIB national. \textbf{De plus}, elle assure la sécurité alimentaire du pays et fournit des matières premières à nos industries locales. \textbf{Cependant}, ce secteur fait face à des défis majeurs : le vieillissement des exploitations, l'insuffisance des infrastructures rurales et l'accès difficile au financement. \textbf{Néanmoins}, des solutions existent. \textbf{Par exemple}, la modernisation des techniques culturales et la formation des jeunes agriculteurs permettraient d'accroître les rendements. \textbf{Par conséquent}, l'État doit accorder la priorité budgétaire à ce secteur stratégique. \textbf{En somme}, investir dans l'agriculture, c'est investir dans l'avenir de la nation.
\end{quote}

\begin{enumerate}
    \item Soulignez dans le texte (ou recopiez) la \cle{thèse} défendue par l'auteur.
    \item Relevez \textbf{trois arguments} distincts développés pour soutenir cette thèse.
    \item Identifiez \textbf{deux exemples} ou illustrations précis donnés dans le texte.
    \item Relevez \textbf{dix connecteurs} logiques présents dans le texte et classez-les selon leur fonction (addition, opposition, cause, conséquence, illustration, conclusion).
\end{enumerate}
\end{exercice}

\begin{exercice}[Reformulation par synonymie et conversion grammaticale]
Reformulez les phrases suivantes en changeant la structure syntaxique (conversion) et/ou le vocabulaire (synonymie), sans en modifier le sens. La phrase reformulée doit être correcte syntaxiquement.

\begin{enumerate}
    \item Le développement de l'agriculture conditionne la prospérité du pays.
    \item Les traditions ancestrales sont respectées par les villageois.
    \item La ruse de Lakunle ne lui permet pas de conquérir Sidi.
    \item Le conflit entre tradition et modernité structure la pièce de Soyinka.
    \item Il est nécessaire que les jeunes s'intéressent à l'agro-industrie.
    \item Bien que le sol soit fertile, les récoltes restent insuffisantes.
    \item Le progrès technique a entraîné une hausse de la production.
    \item L'auteur démontre que l'éducation est la clé du développement.
\end{enumerate}
\end{exercice}

\begin{exercice}[Identification des tonalités dans \emph{Le Lion et la Perle}]
Pour chacun des extraits suivants (références à la pièce de Wole Soyinka), identifiez la \cle{tonalité} dominante (comique, satirique, dramatique, tragique) et justifiez votre réponse en relevant \textbf{deux procédés} (lexique, syntaxe, figure de style, situation, ironie, etc.).

\begin{enumerate}
    \item \textit{Extrait 1 (Scène de la danse du "Voyageur perdu") :} Lakunle mime de façon grotesque la voiture qui tombe en panne, imitant le bruit du moteur avec sa bouche, gesticulant de manière ridicule sous le regard amusé des villageoises.
    \item \textit{Extrait 2 (Monologue de Baroka sur la virilité) :} "Je suis le Lion... Je chasse, je mange, je dors... Et pourtant, la vieillesse me guette comme un voleur dans la nuit. Ma force s'en va, l'eau de la source tarit."
    \item \textit{Extrait 3 (Dialogue Lakunle / Sidi sur la dot) :} Lakunle : "Payer la dot ? C'est acheter une femme comme on achète une chèvre au marché ! C'est barbare, rétrograde, indigne d'un homme moderne !" Sidi : "C'est notre coutume. Sans dot, je suis une fille perdue, sans honneur."
    \item \textit{Extrait 4 (La ruse de Baroka / Sadiku) :} Sadiku danse, jubilant, brandissant l'image du Lion vaincu : "Les femmes l'ont eu ! Le grand Baroka est impuissant !" Le public sait que c'est un piège ourdi par le Bale.
\end{enumerate}
\end{exercice}

\begin{exercice}[Entraînement lexical : Synonymes et Antonymes]
Pour chacun des mots suivants, fréquents dans les textes de résumé (économie, société, littérature), donnez \textbf{un synonyme} et \textbf{un antonyme}, puis employez le mot initial dans une phrase courte montrant son sens contextuel.

\begin{center}
\begin{tabularx}{\linewidth}{|c|X|X|X|}\hline
\textbf{Mot} & \textbf{Synonyme} & \textbf{Antonyme} & \textbf{Phrase d'emploi} \\ \hline
Développement & & & \\ \hline
Tradition & & & \\ \hline
Progrès & & & \\ \hline
Ruse & & & \\ \hline
Conflit & & & \\ \hline
Modernité & & & \\ \hline
Stratégique & & & \\ \hline
Rendement & & & \\ \hline
Héritage & & & \\ \hline
Pérennité & & & \\ \hline
\end{tabularx}
\end{center}
\end{exercice}

\courssub{Je me prépare au Bac}

\begin{exercice}[Sujet type Baccalauréat Technique : Contraction de texte]
\textbf{Consigne :} Lisez attentivement le texte ci-dessous. Vous en ferez une \cle{contraction} de \textbf{80 mots ($\pm 10\%$)}. Vous indiquerez le nombre de mots utilisés à la fin de votre résumé.

\textbf{Texte support (320 mots) :}

\textit{L'agriculture camerounaise à l'épreuve de la modernisation}

Le Cameroun, souvent désigné comme le "grenier de l'Afrique centrale", possède des atouts naturels exceptionnels pour l'agriculture : une diversité climatique allant du sahelien au équatorial, des sols volcaniques fertiles, un réseau hydrographique dense. Pourtant, le secteur peine à jouer pleinement son rôle de levier de développement. La production reste largement dominée par l'agriculture familiale, pratiquée sur de petites exploitations avec des outils rudimentaires (houe, machette). Le rendement à l'hectare y est parmi les plus bas du continent.

Paradoxalement, le pays importe massivement des produits qu'il pourrait produire : riz, blé, poisson, lait. Cette dépendance alimentaire grève la balance commerciale et expose la nation aux fluctuations des cours mondiaux. Les causes sont structurelles : enclavement des zones de production par manque de routes rurales, difficulté d'accès au crédit pour l'achat d'intrants (engrais, semences améliorées), vieillissement de la paysannerie et désintérêt de la jeunesse pour un métier perçu comme pénible et peu rémunérateur.

Face à ce constat, les pouvoirs publics ont lancé plusieurs initiatives : la Société de Développement du Coton (SODECOTON), le Programme d'Appui à la Compétitivité de l'Agriculture (PACA), la création de la Banque Nationale de Développement (BND) pour financer l'agro-business. L'objectif affiché est la "deuxième génération de l'agriculture" : mécanisation, irrigation, transformation locale, insertion des jeunes diplômés.

Toutefois, la réussite de cette mutation suppose un changement de paradigme. Il ne suffit pas de distribuer des tracteurs ; il faut organiser les producteurs en coopératives solides, sécuriser le foncier rural, former aux nouvelles technologies et garantir des prix rémunérateurs. L'agriculture ne doit plus être une activité de subsistance mais une entreprise moderne, créatrice de richesse et d'emplois. C'est à cette condition que le Cameroun atteindra l'autosuffisance alimentaire et deviendra un acteur majeur de l'agro-exportation en Afrique centrale.

\begin{enumerate}
    \item Dégagez la \cle{thèse} principale de ce texte.
    \item Identifiez la \cle{structure argumentative} : thèse, arguments (au moins 3), exemples/illustrations, conclusion.
    \item Rédigez le \textbf{résumé} en 80 mots ($\pm 10\%$). Respectez l'ordre des idées, éliminez les détails, reformulez.
    \item Indiquez le nombre exact de mots de votre production.
\end{enumerate}
\end{exercice}

\begin{exercice}[Sujet type Probatoire : Résumé et Justification des suppressions]
\textbf{Consigne :} Après avoir lu le texte suivant (environ 240 mots), vous en ferez un \textbf{résumé en un quart de sa longueur (environ 60 mots)}. Ensuite, vous \textbf{justifierez en cinq lignes maximum} les principales suppressions opérées (exemples, détails, répétitions, commentaires personnels).

\textbf{Texte :}

\textit{La jeunesse face à l'entrepreneuriat agricole}

L'exode rural vide nos campagnes de leurs forces vives. Chaque année, des milliers de jeunes diplômés ou non affluent vers Douala et Yaoundé, grossissant les rangs du chômage urbain. Pourtant, la terre ne demande qu'à être travaillée. Le Ministère de l'Agriculture et du Développement Rural (MINADER) a mis en place le Programme "Jeunes Entrepreneurs Agricoles" (JEA) doté d'un fonds de 5 milliards de FCFA. Ce programme offre des formations techniques, un accompagnement à la création de business plans et un accès facilité au crédit via la BND.

Prenons l'exemple de Paul, 28 ans, titulaire d'un BTS en Gestion. Après deux ans de chômage, il a intégré le programme JEA. Il a obtenu un prêt de 2 millions pour démarrer un élevage de poulets de chair à 30 km de Yaoundé. Aujourd'hui, il emploie trois personnes et réalise un chiffre d'affaires mensuel de 500 000 FCFA. "La terre ne ment pas", dit-il. De même, Marie, 24 ans, a transformé sa plantation de plantains en une unité de fabrication de chips conditionnés, vendus en supermarchés.

Ces réussites prouvent que l'agriculture est un gisement d'emplois décents. Mais les obstacles persistent : la paperasse administrative pour obtenir les titres fonciers, la fluctuation des prix à la récolte, le manque d'électricité pour la conservation. L'État doit simplifier les procédures foncières et investir dans l'énergie rurale. En conclusion, l'entrepreneuriat agricole n'est pas une voie de garage, c'est une voie royale pour la jeunesse camerounaise, à condition d'un environnement des affaires favorable.

\begin{enumerate}
    \item Rédigez le résumé (environ 60 mots).
    \item Justifiez en 5 lignes les suppressions effectuées.
    \item Indiquez le nombre de mots du résumé.
\end{enumerate}
\end{exercice}

\begin{exercice}[Remédiation : Analyse et correction d'un résumé fautif]
Voici un texte source (100 mots) et une \textbf{proposition de résumé fautive} (65 mots) réalisée par un élève. Le résumé demandé était de \textbf{30 mots ($\pm 10\%$)}.

\textbf{Texte source :}
\textit{Le changement climatique menace la production cacaoyère en Afrique de l'Ouest. La hausse des températures et l'irrégularité des pluies réduisent les rendements et favorisent les maladies fongiques comme le "balai de sorcier". La Côte d'Ivoire et le Ghana, premiers producteurs mondiaux, voient leur économie fragilisée. Pour s'adapter, la recherche agronomique développe des variétés tolérantes à la sécheresse. L'agroforesterie, qui associe cacaoyers et arbres d'ombrage, régule le microclimat. Ces solutions techniques nécessitent un financement massif et une vulgarisation efficace auprès des millions de petits planteurs analphabètes.}

\textbf{Résumé élève (fautif) :}
\textit{Le réchauffement climatique menace le cacao en Afrique de l'Ouest. La chaleur et la pluie irrégulière baissent les rendements et apportent des maladies comme le balai de sorcier. La Côte d'Ivoire et le Ghana, les deux premiers producteurs du monde, sont en danger économique. Heureusement, la recherche crée des variétés résistantes à la sécheresse. L'agroforesterie, qui mélange cacaoyers et grands arbres, protège les plants. Mais ces solutions coûtent très cher et il faut les expliquer aux petits planteurs qui ne savent pas lire. À mon avis, l'État doit payer pour eux.}

\begin{enumerate}
    \item Relevez \textbf{quatre erreurs majeures} commises par l'élève par rapport aux consignes du résumé (nombre de mots, fidélité, objectivité, sélection).
    \item Proposez une \textbf{version corrigée} du résumé respectant la limite de 30 mots ($\pm 10\%$), l'objectivité et la reformulation.
    \item Indiquez le nombre de mots de votre version corrigée.
\end{enumerate}
\end{exercice}

\newpage

\coursec{Corrigés des exercices}

\coursec{III : Formuler les idées essentielles d'un texte à résumer}

\courssub{Je vérifie mes connaissances}

\begin{corrige}[Exercice 1 — QCM : Le texte argumentatif et ses composantes]
\textbf{1. B.} La \cle{thèse} est l'idée centrale, l'opinion que l'auteur soutient et tente de faire valider par le lecteur. Ce n'est pas une simple opinion gratuite (A), ni un exemple (C), ni une concession (D).

\textbf{2. C.} \textbf{Cependant} marque l'\cle{opposition} (ou la concession). \emph{Parce que}, \emph{puisque}, \emph{comme} introduisent une \cle{cause}.

\textbf{3. C.} La finalité du texte argumentatif est de \cle{convaincre} (par la raison) ou de \cle{persuader} (par l'émotion). Il ne vise pas à divertir (A), informer neutrement (B) ou décrire (D).

\textbf{4. C.} La \cle{contraction de texte} (ou résumé) impose une \cle{reformulation} personnelle des idées forces, en éliminant les détails, exemples, répétitions et commentaires subjectifs. Le copier-coller (A), la conservation des exemples (B) ou l'avis personnel (D) sont des erreurs sanctionnées.
\end{corrige}

\begin{corrige}[Exercice 2 — Vrai ou Faux justifié : Tonalités et Lexique]
\begin{enumerate}
    \item \textbf{Faux.} La \cle{tonalité satirique} utilise le rire, l'ironie, la caricature pour \textbf{critiquer} des travers sociaux, moraux ou politiques (vices, travers, injustices). Le rire est une arme, pas une fin en soi.
    \item \textbf{Vrai.} L'\cle{antonymie} unit deux termes de sens contraire (ex. : \emph{progrès / régression}, \emph{clair / obscur}). On distingue antonymes complémentaires (mort/vivant) et graduels (chaud/froid).
    \item \textbf{Faux.} Lakunle incarne la \cle{tonalité comique} (par son ridicule, sa maladresse, son langage ampoulé) et \cle{satirique} (moquerie de la modernité superficielle). La tonalité \cle{tragique} concerne plutôt le destin inéluctable (ex. : Baroka face à la vieillesse/mort).
    \item \textbf{Faux.} La \cle{synonymie} est une relation de \textbf{sens voisin} (proche), non \textbf{identique}. Les synonymes diffèrent souvent par le registre (familier/soutenu), la connotation (péjoratif/laudatif) ou l'intensité. Ex. : \emph{maison / demeure / baraque / logis}.
    \item \textbf{Vrai.} Deux contraintes formelles majeures du résumé : le \textbf{respect de la limite quantitative} (nombre de mots $\pm 10\%$) et le \textbf{respect de l'ordre logique} des idées du texte source (fidélité au plan).
\end{enumerate}
\end{corrige}

\begin{corrige}[Exercice 3 — Définitions et notions clés]
\begin{enumerate}
    \item \textbf{Argument d'autorité} : Argument qui s'appuie sur le témoignage ou l'avis d'un \cle{spécialiste}, d'un \cle{témoin qualifié} ou d'une \cle{source reconnue} (institution, ouvrage de référence) pour valider une thèse. \textit{Exemple : « Selon le rapport 2023 de la FAO, l'insécurité alimentaire recule. »}
    \item \textbf{Concession} : Procédé argumentatif où l'orateur \textbf{accorde un point} à l'adversaire (« Certes, il est vrai que... », « Admettons que... ») pour \textbf{mieux le réfuter ensuite} » (« mais », « cependant »). Cela montre l'honnêteté intellectuelle et renforce la crédibilité.
    \item \textbf{Tonalité dramatique} : Ton qui exprime la \cle{tension}, le \cle{conflit} intense, l'angoisse face à un destin souvent funeste. Elle suscite l'émotion (pitié, terreur) par des situations extrêmes, un pathétique fort, un lexique de la fatalité, de la souffrance, de la mort. \textit{Exemple : Le monologue de Baroka sur la vieillesse (extrait 2).}
    \item \textbf{Conversion grammaticale} (ou \cle{transposition}) : Technique de reformulation consistant à \textbf{changer la nature grammaticale} d'un mot ou la structure de la phrase sans en changer le sens. \textit{Exemples :} Verbe $\to$ Nom (\emph{il développe} $\to$ \emph{le développement}) ; Actif $\to$ Passif (\emph{Les villageois respectent les traditions} $\to$ \emph{Les traditions sont respectées par les villageois}) ; Proposition $\to$ Groupe nominal (\emph{Parce qu'il pleut} $\to$ \emph{À cause de la pluie}).
    \item \textbf{Idée essentielle} : Information \textbf{indispensable} à la compréhension du message global ; elle structure le texte (thèse, arguments principaux). \textbf{Idée secondaire} : Information \textbf{accessoire} (exemple, détail, précision, illustration, répétition) qui développe ou appuie l'idée essentielle mais peut être supprimée au résumé.
\end{enumerate}
\end{corrige}

\begin{corrige}[Exercice 4 — Application directe : Connecteurs et Structure]
\begin{center}
\begin{tabularx}{\linewidth}{|l|X|}\hline
\textbf{Fonction logique} & \textbf{Connecteurs classés} \\ \hline
Addition & De plus ; D'ailleurs \\ \hline
Opposition & Or ; Néanmoins \\ \hline
Cause & En effet ; Car ; Du fait que \\ \hline
Conséquence & Par conséquent ; Ainsi \\ \hline
Concession & Bien que \\ \hline
\end{tabularx}
\end{center}
\medskip
\emph{Note :} \emph{Or} marque une opposition forte (souvent en tête de phrase). \emph{Néanmoins} (adverbe) et \emph{Bien que} (conjonction + subjonctif) marquent l'opposition/concession. \emph{En effet}, \emph{Car}, \emph{Du fait que} expliquent la cause. \emph{Par conséquent}, \emph{Ainsi} tirent la conséquence. \emph{De plus}, \emph{D'ailleurs} ajoutent un argument.
\end{corrige}

\courssub{Je m'entraîne}

\begin{corrige}[Exercice 5 — Repérage dans un texte argumentatif : L'agriculture au Cameroun]
\begin{enumerate}
    \item \textbf{Thèse} : \textit{« Il est indéniable que l'agriculture demeure le moteur principal de l'économie camerounaise »} (ou la phrase finale : \textit{« En somme, investir dans l'agriculture, c'est investir dans l'avenir de la nation »}).
    \item \textbf{Trois arguments} :
    \begin{enumerate}
        \item Poids socio-économique : emploi (60\% population active) et contribution au PIB (20\%).
        \item Rôle vital : sécurité alimentaire et fourniture de matières premières aux industries.
        \item Existence de solutions : modernisation techniques et formation des jeunes pour accroître les rendements.
    \end{enumerate}
    \item \textbf{Deux exemples/illustrations} :
    \begin{enumerate}
        \item \textit{« la modernisation des techniques culturales »}
        \item \textit{« la formation des jeunes agriculteurs »}
    \end{enumerate}
    \item \textbf{Connecteurs logiques (7 identifiés dans le texte)} :
    \begin{center}
    \begin{tabularx}{\linewidth}{|l|l|X|}\hline
    \textbf{Connecteur} & \textbf{Fonction} & \textbf{Contexte} \\ \hline
    En effet & Cause / Justification & Justifie le poids économique (chiffres). \\ \hline
    De plus & Addition & Ajoute un 2e argument (sécurité alimentaire). \\ \hline
    Cependant & Opposition & Introduit les défis/difficultés. \\ \hline
    Néanmoins & Opposition / Concession & Passe aux solutions malgré les défis. \\ \hline
    Par exemple & Illustration & Introduit les solutions concrètes. \\ \hline
    Par conséquent & Conséquence & Déduit la recommandation (priorité budgétaire). \\ \hline
    En somme & Conclusion / Synthèse & Reformule la thèse finale. \\ \hline
    \end{tabularx}
    \end{center}
    \textit{Remarque : Le texte ne contient que 7 connecteurs logiques majeurs (hors coordinations « et »). La consigne en demandait 10 ; les 3 manquants sont absents du support fourni.}
\end{enumerate}
\end{corrige}

\begin{corrige}[Exercice 6 — Reformulation par synonymie et conversion grammaticale]
\textit{Plusieurs reformulations sont possibles. En voici des modèles corrects.}
\begin{enumerate}
    \item \textbf{La prospérité du pays dépend du développement agricole.} (Conversion : sujet/objet ; Synonymie : conditionne $\to$ dépend).
    \item \textbf{Les villageois respectent les traditions ancestrales.} (Conversion : Passif $\to$ Actif).
    \item \textbf{Malgré sa ruse, Lakunle échoue à conquérir Sidi.} (Conversion : Négation + structure ; Synonymie : ne permet pas $\to$ échoue à).
    \item \textbf{L'opposition entre tradition et modernité structure la pièce de Soyinka.} (Synonymie : conflit $\to$ opposition ; Conversion : Sujet/Objet).
    \item \textbf{Les jeunes doivent s'intéresser à l'agro-industrie.} (Conversion : Impersonnel « Il est nécessaire que » + Subjonctif $\to$ Verbe modal « devoir » + Infinitif).
    \item \textbf{Malgré la fertilité du sol, les récoltes restent insuffisantes.} (Conversion : Subordonnée concessive « Bien que » $\to$ Groupe prépositionnel « Malgré » + Nominalisation « fertile $\to$ fertilité »).
    \item \textbf{La production a augmenté grâce au progrès technique.} (Conversion : Actif $\to$ Passif / Cause $\to$ Conséquence ; Synonymie : entraîné une hausse $\to$ augmenté).
    \item \textbf{L'auteur prouve le rôle clé de l'éducation dans le développement.} (Synonymie : démontre $\to$ prouve ; clé $\to$ rôle clé ; conversion : proposition $\to$ GN).
\end{enumerate}
\end{corrige}

\begin{corrige}[Exercice 7 — Identification des tonalités dans \emph{Le Lion et la Perle}]
\begin{enumerate}
    \item \textbf{Tonalité : Comique (et Satirique).}
    \textbf{Procédés :}
    \begin{itemize}
        \item \textit{Mime grotesque / Gestes ridicules :} « mime de façon grotesque », « gesticulant de manière ridicule » $\to$ effet visuel burlesque, caricature de la modernité (la voiture en panne).
        \item \textit{Onomatopées / Bruitage oral :} « imitant le bruit du moteur avec sa bouche » $\to$ ridicule du personnage qui se prend pour une machine, satire de l'imitation servile de l'Occident.
    \end{itemize}
    \item \textbf{Tonalité : Dramatique (teintée de Tragique).}
    \textbf{Procédés :}
    \begin{itemize}
        \item \textit{Lexique de la finitude et de la perte :} « vieillesse », « guette », « voleur », « nuit », « force s'en va », « tarit » $\to$ champ lexical de la mort, du déclin inéluctable.
        \item \textit{Métaphore filée et personnification :} « la vieillesse me guette comme un voleur dans la nuit », « l'eau de la source tarit » $\to$ image poignante de la virilité qui s'épuise, conscience tragique du temps qui passe.
    \end{itemize}
    \item \textbf{Tonalité : Satirique (et Comique).}
    \textbf{Procédés :}
    \begin{itemize}
        \item \textit{Ironie et comparaison dévalorisante :} « C'est acheter une femme comme on achète une chèvre au marché ! » $\to$ critique violente de la coutume (dot) par le prisme de la modernité (marchandisation), ridiculisation de la tradition.
        \item \textit{Lexique péjoratif accumulé :} « barbare, rétrograde, indigne » $\to$ jugement de valeur excessif, caricatural, propre à la satire idéologique.
    \end{itemize}
    \item \textbf{Tonalité : Comique (Ironie dramatique).}
    \textbf{Procédés :}
    \begin{itemize}
        \item \textit{Ironie dramatique (savoir inégal) :} Le public sait que Baroka ment (piège), Sadiku et les femmes croient à sa défaite $\to$ effet comique de supériorité du spectateur.
        \item \textit{Scène burlesque / Carnavalisation :} « Sadiku danse, jubilant, brandissant l'image du Lion vaincu » $\to$ renversement temporaire du pouvoir, fête grotesque du triomphe féminin supposé, registre du carnaval.
    \end{itemize}
\end{enumerate}
\end{corrige}

\begin{corrige}[Exercice 8 — Entraînement lexical : Synonymes et Antonymes]
\begin{center}
\begin{tabularx}{\linewidth}{|c|X|X|X|}\hline
\textbf{Mot} & \textbf{Synonyme} & \textbf{Antonyme} & \textbf{Phrase d'emploi} \\ \hline
Développement & Essor, expansion, croissance & Régression, stagnation, sous-développement & Le \cle{développement} rural est la priorité du nouveau plan quinquennal. \\ \hline
Tradition & Coutume, héritage, usage & Modernité, innovation, rupture & La \cle{tradition} orale transmet l'histoire et les valeurs aux jeunes générations. \\ \hline
Progrès & Avancée, amélioration, perfectionnement & Régression, recul, déclin & Le \cle{progrès} technique a permis de mécaniser les cultures vivrières. \\ \hline
Ruse & Stratagème, tromperie, fourberie & Franchise, honnêteté, loyauté & La \cle{ruse} de Baroka lui permet de duper Sadiku et de conquérir Sidi. \\ \hline
Conflit & Opposition, affrontement, discorde & Accord, harmonie, concorde & Le \cle{conflit} entre Lakunle et Baroka illustre le choc des cultures. \\ \hline
Modernité & Nouveauté, contemporain, avant-garde & Tradition, archaïsme, passé & Lakunle incarne une \cle{modernité} superficielle, faite d'imitation vestimentaire. \\ \hline
Stratégique & Crucial, décisif, capital & Secondaire, accessoire, marginal & L'agriculture est un secteur \cle{stratégique} pour la souveraineté alimentaire. \\ \hline
Rendement & Productivité, rapport, efficacité & Gaspillage, perte, inefficacité & Le \cle{rendement} à l'hectare reste faible faute d'intrants adaptés. \\ \hline
Héritage & Legs, patrimoine, transmission & Oubli, rupture, dissipation & La préservation de l'\cle{héritage} culturel passe par l'éducation. \\ \hline
Pérennité & Durabilité, longévité, permanence & Éphémère, fugacité, instabilité & La \cle{pérennité} des exploitations agricoles nécessite une gestion durable des sols. \\ \hline
\end{tabularx}
\end{center}
\end{corrige}

\begin{methode}[Méthode de la contraction de texte (résumé)]
\begin{enumerate}
    \item \textbf{Lire et analyser} : Identifier le thème, la thèse, la structure (arguments, exemples, concessions), les connecteurs logiques.
    \item \textbf{Sélectionner} : Ne garder que les idées essentielles (thèse + arguments principaux). Éliminer : exemples, détails, chiffres précis, noms propres de programmes, répétitions, commentaires subjectifs.
    \item \textbf{Reformuler} : Appliquer la synonymie et la conversion grammaticale (transposition). Changer la structure des phrases, le vocabulaire, sans trahir le sens. Rester neutre (pas d'avis personnel).
    \item \textbf{Ordonner} : Respecter l'enchaînement logique du texte source (fidélité au plan).
    \item \textbf{Rédiger et compter} : Écrire le brouillon, compter les mots, ajuster pour respecter la fourchette ($\pm 10\%$). Soigner la cohérence (connecteurs, ponctuation).
    \item \textbf{Relire} : Vérifier fidélité, neutralité, longueur, correction linguistique.
\end{enumerate}
\end{methode}

\courssub{Je me prépare au Bac}

\begin{corrige}[Exercice 9 — Sujet type Baccalauréat Technique : Contraction de texte]
\textbf{Barème indicatif (sur 20) :}
\begin{itemize}
    \item Identification thèse + structure : 4 pts
    \item Qualité de la reformulation (synonymie, conversion, neutralité) : 8 pts
    \item Respect de la longueur (72--88 mots) : 4 pts
    \item Respect de l'ordre logique / cohérence : 4 pts
\end{itemize}
\textbf{Points de vigilance :} Ne pas citer les noms de programmes (SODECOTON, PACA, BND) ni les chiffres précis (60\%, 20\%). Grouper : « programmes publics », « initiatives étatiques ». Éliminer la liste des outils (houe, machette) et des importations (riz, blé, poisson, lait). Reformuler « deuxième génération » par « agriculture moderne/entreprise ».

\begin{enumerate}
    \item \textbf{Thèse principale :} Malgré des atouts naturels exceptionnels, l'agriculture camerounaise reste peu performante et dépendante ; sa modernisation en entreprise rentable, via des réformes structurelles (foncier, formation, organisation), est la condition de l'autosuffisance et de l'exportation.
    \item \textbf{Structure argumentative :}
    \begin{itemize}
        \item \textbf{Thèse} : Atouts naturels gâchés, secteur en difficulté (importations, faible productivité).
        \item \textbf{Argument 1 (Causes)} : Structurelles (enclavement, crédit, vieillissement, désintérêt jeunesse).
        \item \textbf{Argument 2 (Réponses)} : Initiatives publiques (programmes, « 2e génération » : mécanisation, transformation).
        \item \textbf{Argument 3 (Conditions)} : Changement de paradigme indispensable (coopératives, foncier, formation, prix).
        \item \textbf{Conclusion} : Agriculture = entreprise moderne $\to$ autosuffisance + agro-exportation.
    \end{itemize}
    \item \textbf{Résumé (86 mots) :}
    \begin{quote}
    Le Cameroun, grenier de l'Afrique centrale, possède d'immenses atouts agroécologiques. Pourtant, l'agriculture familiale, peu productive, ne couvre pas les besoins nationaux, entraînant d'importantes importations. Les causes sont structurelles : enclavement rural, accès difficile au crédit, vieillissement de la paysannerie, désintérêt des jeunes. L'État répond par des programmes publics visant une agriculture moderne : mécanisation, irrigation, transformation locale. Mais la réussite impose un changement de paradigme : organisation en coopératives, sécurisation foncière, formation adaptée, prix rémunérateurs. L'agriculture doit devenir une entreprise rentable pour assurer l'autosuffisance et l'agro-exportation.
    \end{quote}
    \item \textbf{Nombre de mots : 86 mots} (dans la fourchette 72--88).
\end{enumerate}
\end{corrige}

\begin{corrige}[Exercice 10 — Sujet type Probatoire : Résumé et Justification des suppressions]
\textbf{Barème indicatif (sur 20) :}
\begin{itemize}
    \item Résumé (contenu, forme, longueur 54--66 mots) : 12 pts
    \item Justification des suppressions (pertinence, concision 5 lignes) : 8 pts
\end{itemize}
\begin{enumerate}
    \item \textbf{Résumé (63 mots) :}
    \begin{quote}
    L'exode rural vide les campagnes au profit du chômage urbain. Le programme JEA soutient l'entrepreneuriat agricole (crédit, formation), générant des réussites créatrices d'emplois et revenus. Cependant, l'accès au foncier, la volatilité des prix et le déficit énergétique freinent l'essor. L'État doit sécuriser le foncier et électrifier les campagnes. L'agriculture devient une voie royale pour la jeunesse, à condition d'un cadre des affaires favorable.
    \end{quote}
    \item \textbf{Justification des suppressions (5 lignes) :}
    \begin{itemize}
        \item Suppression des exemples nominatifs (Paul, Marie) et chiffrés (montants, CA, employés) : détails illustratifs.
        \item Suppression des détails du programme JEA (5 milliards, business plan, BND) : précisions secondaires.
        \item Suppression de la citation subjective « La terre ne ment pas » : commentaire personnel.
        \item Regroupement des obstacles en trois catégories synthétiques (foncier, prix, énergie) : élimination du récit.
        \item Conservation de la thèse finale (« voie royale ») reformulée sans l'opposition « voie de garage ».
    \end{itemize}
    \item \textbf{Nombre de mots du résumé : 63 mots} (dans la fourchette 54--66).
\end{enumerate}
\end{corrige}

\begin{corrige}[Exercice 11 — Remédiation : Analyse et correction d'un résumé fautif]
\begin{enumerate}
    \item \textbf{Quatre erreurs majeures de l'élève :}
    \begin{enumerate}
        \item \textbf{Non-respect de la contrainte de longueur :} 65 mots produits pour 30 demandés ($\pm 10\%$ = 27--33 mots). Dépassement de plus de 100\%.
        \item \textbf{Subjectivité / Apport personnel :} Ajout de « Heureusement » (marque d'opinion) et surtout « À mon avis, l'État doit payer pour eux » (recommandation absente du texte source). Le résumé doit être \textbf{objectif} et \textbf{fidèle}.
        \item \textbf{Défaut de reformulation (quasi copier-coller) :} « la chaleur et la pluie irrégulière baissent les rendements » (source : « la hausse des températures et l'irrégularité des pluies réduisent les rendements ») ; « ces solutions coûtent très cher » (source : « nécessitent un financement massif ») ; « il faut les expliquer aux petits planteurs qui ne savent pas lire » (source : « vulgarisation efficace auprès des millions de petits planteurs analphabètes »). Absence de conversion grammaticale et de synonymie suffisante.
        \item \textbf{Mauvaise sélection (conservation du secondaire) :} Maintien des noms propres « Côte d'Ivoire et le Ghana » et de l'exemple précis « balai de sorcier » au lieu de généraliser (« premiers producteurs mondiaux », « maladies fongiques »).
    \end{enumerate}
    \item \textbf{Version corrigée (33 mots) :}
    \begin{quote}
    Le réchauffement menace la cacaoculture ouest-africaine (rendements en baisse, maladies). L'économie des premiers producteurs est fragilisée. Variétés résistantes et agroforesterie sont des solutions, mais leur diffusion exige financement massif et formation des planteurs.
    \end{quote}
    \item \textbf{Nombre de mots : 33 mots} (dans la fourchette 27--33).
\end{enumerate}
\end{corrige}

\newpage

\coursec{Fiche bilan}

\begin{popfigure}
\begin{tikzpicture}[mindmap, grow cyclic, every node/.style={concept, font=\small}, text width=2.6cm, level 1/.append style={sibling angle=51, level distance=3.6cm}, root concept/.append style={concept color=popblue, text=white, font=\small\bfseries}]
\node[root concept]{Résumé : formuler les idées essentielles}
  child[concept color=poporange!80]{ node{Texte argumentatif : thèse, arguments, exemples, connecteurs} }
  child[concept color=popgreen!80]{ node{Structure : introduction, développement, conclusion} }
  child[concept color=poppurple!80]{ node{Tonalités : comique, dramatique, satirique, tragique} }
  child[concept color=poppink!80]{ node{Le Lion et la Perle : lectures méthodiques 4 et 5} }
  child[concept color=popgold!85]{ node{Synonymie et antonymie : outils de la reformulation} }
  child[concept color=popblue!60]{ node{Méthode : repérer, hiérarchiser, reformuler, rédiger} }
  child[concept color=popgreen!60]{ node{Remédiation et entraînement au Probatoire/Bac} };
\end{tikzpicture}
\legende{Carte mentale de la leçon : les sept savoirs à mobiliser pour formuler les idées essentielles d'un texte à résumer.}
\end{popfigure}

\begin{bilan}
\textbf{1. Définitions clés}
\begin{itemize}
  \item \cle{Texte argumentatif} : texte qui défend une \cle{thèse} (point de vue) à l'aide d'\cle{arguments} (raisons) illustrés par des \cle{exemples}.
  \item \cle{Tonalité} : ton dominant d'un texte (comique, dramatique, satirique, tragique), repérable par le vocabulaire, le rythme et les procédés (ironie, hyperbole, registre pathétique).
  \item \cle{Synonymie} : relation entre mots de sens proche (\emph{important} / \emph{essentiel}) ; \cle{antonymie} : relation entre mots de sens opposé (\emph{vrai} / \emph{faux}).
  \item \cle{Reformulation} : dire autrement, avec ses propres mots, sans trahir la pensée de l'auteur.
\end{itemize}

\textbf{2. À connaître par cœur}
\begin{center}
\begin{tabularx}{\linewidth}{|l|X|}\hline
\rowcolor{popblueL} \textbf{Notion} & \textbf{Contenu à retenir} \\ \hline
Structure argumentative & Introduction (sujet + thèse), développement (arguments + exemples), conclusion (bilan + ouverture). \\ \hline
Connecteurs logiques & Addition (\emph{de plus}), opposition (\emph{cependant}), cause (\emph{car}), conséquence (\emph{donc}), illustration (\emph{par exemple}). \\ \hline
Règle du résumé & Environ \textbf{un quart} du texte ; objectivité totale ; aucune citation longue ; phrases personnelles. \\ \hline
Tonalités & Comique (rire), satirique (critique par la moquerie), dramatique (tension, conflit), tragique (fatalité, mort annoncée). \\ \hline
\end{tabularx}
\end{center}

\textbf{3. Méthode express du résumé}
\begin{enumerate}
  \item Lire deux fois ; repérer la thèse et le plan.
  \item Souligner les idées essentielles ; éliminer exemples, répétitions, détails.
  \item Hiérarchiser les idées dans l'ordre du texte.
  \item Reformuler avec synonymes, changements de tournures, regroupements.
  \item Rédiger un texte continu, lié par des connecteurs, sans avis personnel.
  \item Vérifier : longueur, fidélité, orthographe, absence de copie.
\end{enumerate}
\end{bilan}

\begin{aretenir}[Checklist avant le Bac]
\begin{itemize}
  \item[$\square$] Je sais définir un texte argumentatif et repérer sa thèse.
  \item[$\square$] Je distingue argument, exemple et simple détail.
  \item[$\square$] Je connais les trois parties de la structure argumentative.
  \item[$\square$] Je reconnais les connecteurs logiques et leur valeur.
  \item[$\square$] Je sais identifier les quatre tonalités (comique, dramatique, satirique, tragique).
  \item[$\square$] Je maîtrise la synonymie et l'antonymie pour reformuler.
  \item[$\square$] Je respecte la règle du quart de la longueur du texte.
  \item[$\square$] Je rédige avec mes propres mots, sans recopier de phrases du texte.
  \item[$\square$] Je reste objectif : aucun jugement personnel dans le résumé.
  \item[$\square$] Je relie mes phrases par des connecteurs variés.
  \item[$\square$] Je relis pour corriger orthographe, grammaire et ponctuation.
  \item[$\square$] Je sais appliquer la méthode aux extraits du \emph{Lion et la Perle}.
\end{itemize}
\end{aretenir}

\findecours

\end{document}
